\documentclass[11pt,a4paper]{article}
% =====================================================================
%  Preámbulo común para los apuntes de Física Electrónica
%  Uso (desde unidades/XX-nombre/archivo.tex):
%     \documentclass[11pt,a4paper]{article}
%     % =====================================================================
%  Preámbulo común para los apuntes de Física Electrónica
%  Uso (desde unidades/XX-nombre/archivo.tex):
%     \documentclass[11pt,a4paper]{article}
%     % =====================================================================
%  Preámbulo común para los apuntes de Física Electrónica
%  Uso (desde unidades/XX-nombre/archivo.tex):
%     \documentclass[11pt,a4paper]{article}
%     \input{../../comun/preambulo.tex}
% =====================================================================

% ---------- Idioma y tipografía ----------
\usepackage[T1]{fontenc}
\usepackage[utf8]{inputenc}
\usepackage[spanish,es-nodecimaldot,es-noshorthands,es-tabla]{babel}
\usepackage{lmodern}
\usepackage{microtype}
\usepackage[a4paper,margin=2.3cm,headheight=14pt]{geometry}
\usepackage{parskip}
\setlength{\parskip}{0.55em}

% ---------- Matemática y unidades ----------
\usepackage{amsmath,amssymb,mathtools}
\usepackage{bm}
\usepackage[per-mode=symbol,output-decimal-marker={.},separate-uncertainty]{siunitx}
\sisetup{range-phrase={\text{--}},range-units=single}
\DeclareSIUnit{\eV}{eV}

% ---------- Gráficos ----------
\usepackage{graphicx}
\usepackage{xcolor}
\usepackage{tikz}
\usetikzlibrary{arrows.meta,decorations.pathmorphing,decorations.markings,
  patterns,calc,positioning,shapes.geometric,fit,backgrounds,shadings}
\usepackage{pgfplots}
\pgfplotsset{compat=1.18}
\usepgfplotslibrary{fillbetween}
\usepackage[siunitx,RPvoltages]{circuitikz}

% ---------- Tablas, listas, figuras ----------
\usepackage{booktabs,array,tabularx,multirow}
\usepackage{enumitem}
\setlist{itemsep=0.2em,topsep=0.3em}
\usepackage{caption,subcaption}
\captionsetup{font=small,labelfont=bf,format=hang}
\usepackage{float}
\usepackage{wrapfig}
\usepackage{needspace}

% ---------- Colores ----------
\definecolor{azul}{HTML}{1F4E79}
\definecolor{azulclaro}{HTML}{DCE9F5}
\definecolor{verde}{HTML}{2E7D32}
\definecolor{verdeclaro}{HTML}{E3F2E1}
\definecolor{naranja}{HTML}{C55A11}
\definecolor{naranjaclaro}{HTML}{FBE9DB}
\definecolor{rojo}{HTML}{B71C1C}
\definecolor{rojoclaro}{HTML}{FDECEA}
\definecolor{violeta}{HTML}{5E35B1}
\definecolor{violetaclaro}{HTML}{EDE7F6}
\definecolor{gris}{HTML}{555555}
% colores para dibujos de bandas / portadores
\colorlet{colEc}{azul}
\colorlet{colEv}{rojo}
\colorlet{colElec}{azul}
\colorlet{colHueco}{rojo}
\colorlet{colFoton}{naranja}

% ---------- Cajas didácticas ----------
\usepackage[most]{tcolorbox}
\tcbset{
  base/.style={enhanced,breakable,arc=2pt,boxrule=0.6pt,left=6pt,right=6pt,
    top=4pt,bottom=4pt,fonttitle=\bfseries\small,before skip=8pt,after skip=8pt}
}
% Definición
\newtcolorbox{definicion}[1][]{base,colback=azulclaro,colframe=azul,
  title={Definición\ifx\relax#1\relax\else: #1\fi}}
% Resultado / fórmula clave
\newtcolorbox{clave}[1][]{base,colback=verdeclaro,colframe=verde,
  title={\ifx\relax#1\relax Resultado clave\else #1\fi}}
% Ejemplo resuelto
\newtcolorbox{ejemplo}[1][]{base,colback=white,colframe=violeta,
  coltitle=white,colbacktitle=violeta,
  title={Ejemplo\ifx\relax#1\relax\else: #1\fi}}
% Corrección / aclaración respecto de las diapositivas
\newtcolorbox{ojo}[1][]{base,colback=rojoclaro,colframe=rojo,
  title={\ensuremath{\triangle}\ Ojo\ifx\relax#1\relax\else: #1\fi}}
% Intuición física / comentario
\newtcolorbox{intuicion}[1][]{base,colback=naranjaclaro,colframe=naranja,
  title={Idea física\ifx\relax#1\relax\else: #1\fi}}

% ---------- Encabezados ----------
\usepackage{fancyhdr}
\pagestyle{fancy}
\fancyhf{}
\renewcommand{\headrulewidth}{0.4pt}
\fancyhead[L]{\small\color{gris}Física Electrónica}
\fancyhead[R]{\small\color{gris}\unidadtitulo}
\fancyfoot[C]{\small\thepage}
\newcommand{\unidadtitulo}{}
\newcommand{\unidad}[1]{\renewcommand{\unidadtitulo}{#1}}

\usepackage{titlesec}
\titleformat{\section}{\Large\bfseries\color{azul}}{\thesection}{0.6em}{}
\titleformat{\subsection}{\large\bfseries\color{azul}}{\thesubsection}{0.6em}{}
\titleformat{\subsubsection}{\normalsize\bfseries\color{azul!80!black}}{\thesubsubsection}{0.6em}{}

\usepackage[hidelinks,bookmarksnumbered]{hyperref}
\usepackage{bookmark}

% ---------- Macros físicas ----------
\newcommand{\EG}{E_G}
\newcommand{\Ec}{E_c}
\newcommand{\Ev}{E_v}
\newcommand{\EF}{E_F}
\newcommand{\Eph}{E_{\mathrm{ph}}}
\newcommand{\lG}{\lambda_G}
\newcommand{\VT}{V_T}
\newcommand{\Popt}{P_{\mathrm{opt}}}
\newcommand{\Resp}{\mathcal{R}}
\newcommand{\vsat}{v_{\mathrm{sat}}}
\newcommand{\Voc}{V_{oc}}
\newcommand{\Isc}{I_{sc}}
\newcommand{\diapo}[1]{{\small\textsf{[diapo~#1]}}}
\newcommand{\dd}{\,\mathrm{d}}

% Distintivo para resultados que ya están en la hoja de fórmulas oficial.
% Uso: en el título de una caja, \begin{clave}[Título \enhoja], o en texto con \enhojatexto.
\newcommand{\enhojatexto}{{\scriptsize\sffamily\colorbox{verdeclaro}{\textcolor{verde!70!black}{\ensuremath{\checkmark}\,en la hoja}}}}
\newcommand{\enhoja}{\hfill{\scriptsize\sffamily\colorbox{white}{\textcolor{verde!70!black}{\ensuremath{\checkmark}\,en la hoja}}}}

\input{../../comun/constantes.tex}

% ---------- Estilos TikZ comunes ----------
\tikzset{
  banda/.style={line width=1.1pt},
  bandac/.style={banda,colEc},
  bandav/.style={banda,colEv},
  fermi/.style={dashed,gris,line width=0.6pt},
  elec/.style={circle,fill=colElec,inner sep=0pt,minimum size=5pt},
  hueco/.style={circle,draw=colHueco,fill=white,line width=0.9pt,inner sep=0pt,minimum size=5pt},
  foton/.style={colFoton,line width=0.9pt,decorate,
    decoration={snake,amplitude=1.6pt,segment length=5pt,post length=3pt},-{Stealth[length=5pt]}},
  flecha/.style={-{Stealth[length=5pt]},line width=0.7pt},
}

% =====================================================================

% ---------- Idioma y tipografía ----------
\usepackage[T1]{fontenc}
\usepackage[utf8]{inputenc}
\usepackage[spanish,es-nodecimaldot,es-noshorthands,es-tabla]{babel}
\usepackage{lmodern}
\usepackage{microtype}
\usepackage[a4paper,margin=2.3cm,headheight=14pt]{geometry}
\usepackage{parskip}
\setlength{\parskip}{0.55em}

% ---------- Matemática y unidades ----------
\usepackage{amsmath,amssymb,mathtools}
\usepackage{bm}
\usepackage[per-mode=symbol,output-decimal-marker={.},separate-uncertainty]{siunitx}
\sisetup{range-phrase={\text{--}},range-units=single}
\DeclareSIUnit{\eV}{eV}

% ---------- Gráficos ----------
\usepackage{graphicx}
\usepackage{xcolor}
\usepackage{tikz}
\usetikzlibrary{arrows.meta,decorations.pathmorphing,decorations.markings,
  patterns,calc,positioning,shapes.geometric,fit,backgrounds,shadings}
\usepackage{pgfplots}
\pgfplotsset{compat=1.18}
\usepgfplotslibrary{fillbetween}
\usepackage[siunitx,RPvoltages]{circuitikz}

% ---------- Tablas, listas, figuras ----------
\usepackage{booktabs,array,tabularx,multirow}
\usepackage{enumitem}
\setlist{itemsep=0.2em,topsep=0.3em}
\usepackage{caption,subcaption}
\captionsetup{font=small,labelfont=bf,format=hang}
\usepackage{float}
\usepackage{wrapfig}
\usepackage{needspace}

% ---------- Colores ----------
\definecolor{azul}{HTML}{1F4E79}
\definecolor{azulclaro}{HTML}{DCE9F5}
\definecolor{verde}{HTML}{2E7D32}
\definecolor{verdeclaro}{HTML}{E3F2E1}
\definecolor{naranja}{HTML}{C55A11}
\definecolor{naranjaclaro}{HTML}{FBE9DB}
\definecolor{rojo}{HTML}{B71C1C}
\definecolor{rojoclaro}{HTML}{FDECEA}
\definecolor{violeta}{HTML}{5E35B1}
\definecolor{violetaclaro}{HTML}{EDE7F6}
\definecolor{gris}{HTML}{555555}
% colores para dibujos de bandas / portadores
\colorlet{colEc}{azul}
\colorlet{colEv}{rojo}
\colorlet{colElec}{azul}
\colorlet{colHueco}{rojo}
\colorlet{colFoton}{naranja}

% ---------- Cajas didácticas ----------
\usepackage[most]{tcolorbox}
\tcbset{
  base/.style={enhanced,breakable,arc=2pt,boxrule=0.6pt,left=6pt,right=6pt,
    top=4pt,bottom=4pt,fonttitle=\bfseries\small,before skip=8pt,after skip=8pt}
}
% Definición
\newtcolorbox{definicion}[1][]{base,colback=azulclaro,colframe=azul,
  title={Definición\ifx\relax#1\relax\else: #1\fi}}
% Resultado / fórmula clave
\newtcolorbox{clave}[1][]{base,colback=verdeclaro,colframe=verde,
  title={\ifx\relax#1\relax Resultado clave\else #1\fi}}
% Ejemplo resuelto
\newtcolorbox{ejemplo}[1][]{base,colback=white,colframe=violeta,
  coltitle=white,colbacktitle=violeta,
  title={Ejemplo\ifx\relax#1\relax\else: #1\fi}}
% Corrección / aclaración respecto de las diapositivas
\newtcolorbox{ojo}[1][]{base,colback=rojoclaro,colframe=rojo,
  title={\ensuremath{\triangle}\ Ojo\ifx\relax#1\relax\else: #1\fi}}
% Intuición física / comentario
\newtcolorbox{intuicion}[1][]{base,colback=naranjaclaro,colframe=naranja,
  title={Idea física\ifx\relax#1\relax\else: #1\fi}}

% ---------- Encabezados ----------
\usepackage{fancyhdr}
\pagestyle{fancy}
\fancyhf{}
\renewcommand{\headrulewidth}{0.4pt}
\fancyhead[L]{\small\color{gris}Física Electrónica}
\fancyhead[R]{\small\color{gris}\unidadtitulo}
\fancyfoot[C]{\small\thepage}
\newcommand{\unidadtitulo}{}
\newcommand{\unidad}[1]{\renewcommand{\unidadtitulo}{#1}}

\usepackage{titlesec}
\titleformat{\section}{\Large\bfseries\color{azul}}{\thesection}{0.6em}{}
\titleformat{\subsection}{\large\bfseries\color{azul}}{\thesubsection}{0.6em}{}
\titleformat{\subsubsection}{\normalsize\bfseries\color{azul!80!black}}{\thesubsubsection}{0.6em}{}

\usepackage[hidelinks,bookmarksnumbered]{hyperref}
\usepackage{bookmark}

% ---------- Macros físicas ----------
\newcommand{\EG}{E_G}
\newcommand{\Ec}{E_c}
\newcommand{\Ev}{E_v}
\newcommand{\EF}{E_F}
\newcommand{\Eph}{E_{\mathrm{ph}}}
\newcommand{\lG}{\lambda_G}
\newcommand{\VT}{V_T}
\newcommand{\Popt}{P_{\mathrm{opt}}}
\newcommand{\Resp}{\mathcal{R}}
\newcommand{\vsat}{v_{\mathrm{sat}}}
\newcommand{\Voc}{V_{oc}}
\newcommand{\Isc}{I_{sc}}
\newcommand{\diapo}[1]{{\small\textsf{[diapo~#1]}}}
\newcommand{\dd}{\,\mathrm{d}}

% Distintivo para resultados que ya están en la hoja de fórmulas oficial.
% Uso: en el título de una caja, \begin{clave}[Título \enhoja], o en texto con \enhojatexto.
\newcommand{\enhojatexto}{{\scriptsize\sffamily\colorbox{verdeclaro}{\textcolor{verde!70!black}{\ensuremath{\checkmark}\,en la hoja}}}}
\newcommand{\enhoja}{\hfill{\scriptsize\sffamily\colorbox{white}{\textcolor{verde!70!black}{\ensuremath{\checkmark}\,en la hoja}}}}

% =====================================================================
%  Constantes oficiales de la materia ("Tabla de constantes.pdf").
%  Usar SIEMPRE estos valores en los ejemplos numéricos, para que
%  coincidan con los que se usan en clase y en el examen.
% =====================================================================

% ---------- Constantes físicas ----------
\newcommand{\qval}{\num{1.60e-19}}          % C
\newcommand{\kval}{\num{8.61e-5}}           % eV/K
\newcommand{\kJval}{\num{1.38e-23}}         % J/K
\newcommand{\hval}{\num{6.62e-34}}          % J s
\newcommand{\heVval}{\num{4.13e-15}}        % eV s
\newcommand{\epsval}{\num{8.85e-14}}        % F/cm
\newcommand{\VTval}{\SI{26}{mV}}            % kT/q a 300 K (0.0259 V)
\newcommand{\epsrox}{3.9}                   % SiO2

% ---------- Silicio ----------
\newcommand{\epsrSi}{11.8}
\newcommand{\niSi}{\num{1.0e10}}            % cm^-3
\newcommand{\munSi}{1350}                   % cm^2/(V s)
\newcommand{\mupSi}{480}                    % cm^2/(V s)
\newcommand{\EGSi}{1.12}                    % eV
\newcommand{\NCSi}{\num{3.22e19}}           % cm^-3
\newcommand{\NVSi}{\num{1.83e19}}           % cm^-3

% ---------- Germanio ----------
\newcommand{\epsrGe}{15.8}
\newcommand{\niGe}{\num{2.4e13}}
\newcommand{\munGe}{3900}
\newcommand{\mupGe}{1900}
\newcommand{\EGGe}{0.66}

% ---------- GaAs ----------
\newcommand{\epsrGaAs}{13.1}
\newcommand{\niGaAs}{\num{2.1e6}}
\newcommand{\munGaAs}{8800}
\newcommand{\mupGaAs}{400}
\newcommand{\EGGaAs}{1.42}

% ---------- Tabla para incluir en los apuntes ----------
\newcommand{\tablaconstantes}{%
\begin{tabular}{lccc}
\toprule
 & Si & Ge & GaAs \\
\midrule
$\epsilon_r$                     & 11.8 & 15.8 & 13.1 \\
$\mu_n$ (\si{cm^2/V.s})          & 1350 & 3900 & 8800 \\
$\mu_p$ (\si{cm^2/V.s})          & 480  & 1900 & 400 \\
$n_i$ (\si{cm^{-3}}, 300\,K)     & \num{1.0e10} & \num{2.4e13} & \num{2.1e6} \\
$E_G$ (\si{\eV})                 & 1.12 & 0.66 & 1.42 \\
\bottomrule
\end{tabular}}


% ---------- Estilos TikZ comunes ----------
\tikzset{
  banda/.style={line width=1.1pt},
  bandac/.style={banda,colEc},
  bandav/.style={banda,colEv},
  fermi/.style={dashed,gris,line width=0.6pt},
  elec/.style={circle,fill=colElec,inner sep=0pt,minimum size=5pt},
  hueco/.style={circle,draw=colHueco,fill=white,line width=0.9pt,inner sep=0pt,minimum size=5pt},
  foton/.style={colFoton,line width=0.9pt,decorate,
    decoration={snake,amplitude=1.6pt,segment length=5pt,post length=3pt},-{Stealth[length=5pt]}},
  flecha/.style={-{Stealth[length=5pt]},line width=0.7pt},
}

% =====================================================================

% ---------- Idioma y tipografía ----------
\usepackage[T1]{fontenc}
\usepackage[utf8]{inputenc}
\usepackage[spanish,es-nodecimaldot,es-noshorthands,es-tabla]{babel}
\usepackage{lmodern}
\usepackage{microtype}
\usepackage[a4paper,margin=2.3cm,headheight=14pt]{geometry}
\usepackage{parskip}
\setlength{\parskip}{0.55em}

% ---------- Matemática y unidades ----------
\usepackage{amsmath,amssymb,mathtools}
\usepackage{bm}
\usepackage[per-mode=symbol,output-decimal-marker={.},separate-uncertainty]{siunitx}
\sisetup{range-phrase={\text{--}},range-units=single}
\DeclareSIUnit{\eV}{eV}

% ---------- Gráficos ----------
\usepackage{graphicx}
\usepackage{xcolor}
\usepackage{tikz}
\usetikzlibrary{arrows.meta,decorations.pathmorphing,decorations.markings,
  patterns,calc,positioning,shapes.geometric,fit,backgrounds,shadings}
\usepackage{pgfplots}
\pgfplotsset{compat=1.18}
\usepgfplotslibrary{fillbetween}
\usepackage[siunitx,RPvoltages]{circuitikz}

% ---------- Tablas, listas, figuras ----------
\usepackage{booktabs,array,tabularx,multirow}
\usepackage{enumitem}
\setlist{itemsep=0.2em,topsep=0.3em}
\usepackage{caption,subcaption}
\captionsetup{font=small,labelfont=bf,format=hang}
\usepackage{float}
\usepackage{wrapfig}
\usepackage{needspace}

% ---------- Colores ----------
\definecolor{azul}{HTML}{1F4E79}
\definecolor{azulclaro}{HTML}{DCE9F5}
\definecolor{verde}{HTML}{2E7D32}
\definecolor{verdeclaro}{HTML}{E3F2E1}
\definecolor{naranja}{HTML}{C55A11}
\definecolor{naranjaclaro}{HTML}{FBE9DB}
\definecolor{rojo}{HTML}{B71C1C}
\definecolor{rojoclaro}{HTML}{FDECEA}
\definecolor{violeta}{HTML}{5E35B1}
\definecolor{violetaclaro}{HTML}{EDE7F6}
\definecolor{gris}{HTML}{555555}
% colores para dibujos de bandas / portadores
\colorlet{colEc}{azul}
\colorlet{colEv}{rojo}
\colorlet{colElec}{azul}
\colorlet{colHueco}{rojo}
\colorlet{colFoton}{naranja}

% ---------- Cajas didácticas ----------
\usepackage[most]{tcolorbox}
\tcbset{
  base/.style={enhanced,breakable,arc=2pt,boxrule=0.6pt,left=6pt,right=6pt,
    top=4pt,bottom=4pt,fonttitle=\bfseries\small,before skip=8pt,after skip=8pt}
}
% Definición
\newtcolorbox{definicion}[1][]{base,colback=azulclaro,colframe=azul,
  title={Definición\ifx\relax#1\relax\else: #1\fi}}
% Resultado / fórmula clave
\newtcolorbox{clave}[1][]{base,colback=verdeclaro,colframe=verde,
  title={\ifx\relax#1\relax Resultado clave\else #1\fi}}
% Ejemplo resuelto
\newtcolorbox{ejemplo}[1][]{base,colback=white,colframe=violeta,
  coltitle=white,colbacktitle=violeta,
  title={Ejemplo\ifx\relax#1\relax\else: #1\fi}}
% Corrección / aclaración respecto de las diapositivas
\newtcolorbox{ojo}[1][]{base,colback=rojoclaro,colframe=rojo,
  title={\ensuremath{\triangle}\ Ojo\ifx\relax#1\relax\else: #1\fi}}
% Intuición física / comentario
\newtcolorbox{intuicion}[1][]{base,colback=naranjaclaro,colframe=naranja,
  title={Idea física\ifx\relax#1\relax\else: #1\fi}}

% ---------- Encabezados ----------
\usepackage{fancyhdr}
\pagestyle{fancy}
\fancyhf{}
\renewcommand{\headrulewidth}{0.4pt}
\fancyhead[L]{\small\color{gris}Física Electrónica}
\fancyhead[R]{\small\color{gris}\unidadtitulo}
\fancyfoot[C]{\small\thepage}
\newcommand{\unidadtitulo}{}
\newcommand{\unidad}[1]{\renewcommand{\unidadtitulo}{#1}}

\usepackage{titlesec}
\titleformat{\section}{\Large\bfseries\color{azul}}{\thesection}{0.6em}{}
\titleformat{\subsection}{\large\bfseries\color{azul}}{\thesubsection}{0.6em}{}
\titleformat{\subsubsection}{\normalsize\bfseries\color{azul!80!black}}{\thesubsubsection}{0.6em}{}

\usepackage[hidelinks,bookmarksnumbered]{hyperref}
\usepackage{bookmark}

% ---------- Macros físicas ----------
\newcommand{\EG}{E_G}
\newcommand{\Ec}{E_c}
\newcommand{\Ev}{E_v}
\newcommand{\EF}{E_F}
\newcommand{\Eph}{E_{\mathrm{ph}}}
\newcommand{\lG}{\lambda_G}
\newcommand{\VT}{V_T}
\newcommand{\Popt}{P_{\mathrm{opt}}}
\newcommand{\Resp}{\mathcal{R}}
\newcommand{\vsat}{v_{\mathrm{sat}}}
\newcommand{\Voc}{V_{oc}}
\newcommand{\Isc}{I_{sc}}
\newcommand{\diapo}[1]{{\small\textsf{[diapo~#1]}}}
\newcommand{\dd}{\,\mathrm{d}}

% Distintivo para resultados que ya están en la hoja de fórmulas oficial.
% Uso: en el título de una caja, \begin{clave}[Título \enhoja], o en texto con \enhojatexto.
\newcommand{\enhojatexto}{{\scriptsize\sffamily\colorbox{verdeclaro}{\textcolor{verde!70!black}{\ensuremath{\checkmark}\,en la hoja}}}}
\newcommand{\enhoja}{\hfill{\scriptsize\sffamily\colorbox{white}{\textcolor{verde!70!black}{\ensuremath{\checkmark}\,en la hoja}}}}

% =====================================================================
%  Constantes oficiales de la materia ("Tabla de constantes.pdf").
%  Usar SIEMPRE estos valores en los ejemplos numéricos, para que
%  coincidan con los que se usan en clase y en el examen.
% =====================================================================

% ---------- Constantes físicas ----------
\newcommand{\qval}{\num{1.60e-19}}          % C
\newcommand{\kval}{\num{8.61e-5}}           % eV/K
\newcommand{\kJval}{\num{1.38e-23}}         % J/K
\newcommand{\hval}{\num{6.62e-34}}          % J s
\newcommand{\heVval}{\num{4.13e-15}}        % eV s
\newcommand{\epsval}{\num{8.85e-14}}        % F/cm
\newcommand{\VTval}{\SI{26}{mV}}            % kT/q a 300 K (0.0259 V)
\newcommand{\epsrox}{3.9}                   % SiO2

% ---------- Silicio ----------
\newcommand{\epsrSi}{11.8}
\newcommand{\niSi}{\num{1.0e10}}            % cm^-3
\newcommand{\munSi}{1350}                   % cm^2/(V s)
\newcommand{\mupSi}{480}                    % cm^2/(V s)
\newcommand{\EGSi}{1.12}                    % eV
\newcommand{\NCSi}{\num{3.22e19}}           % cm^-3
\newcommand{\NVSi}{\num{1.83e19}}           % cm^-3

% ---------- Germanio ----------
\newcommand{\epsrGe}{15.8}
\newcommand{\niGe}{\num{2.4e13}}
\newcommand{\munGe}{3900}
\newcommand{\mupGe}{1900}
\newcommand{\EGGe}{0.66}

% ---------- GaAs ----------
\newcommand{\epsrGaAs}{13.1}
\newcommand{\niGaAs}{\num{2.1e6}}
\newcommand{\munGaAs}{8800}
\newcommand{\mupGaAs}{400}
\newcommand{\EGGaAs}{1.42}

% ---------- Tabla para incluir en los apuntes ----------
\newcommand{\tablaconstantes}{%
\begin{tabular}{lccc}
\toprule
 & Si & Ge & GaAs \\
\midrule
$\epsilon_r$                     & 11.8 & 15.8 & 13.1 \\
$\mu_n$ (\si{cm^2/V.s})          & 1350 & 3900 & 8800 \\
$\mu_p$ (\si{cm^2/V.s})          & 480  & 1900 & 400 \\
$n_i$ (\si{cm^{-3}}, 300\,K)     & \num{1.0e10} & \num{2.4e13} & \num{2.1e6} \\
$E_G$ (\si{\eV})                 & 1.12 & 0.66 & 1.42 \\
\bottomrule
\end{tabular}}


% ---------- Estilos TikZ comunes ----------
\tikzset{
  banda/.style={line width=1.1pt},
  bandac/.style={banda,colEc},
  bandav/.style={banda,colEv},
  fermi/.style={dashed,gris,line width=0.6pt},
  elec/.style={circle,fill=colElec,inner sep=0pt,minimum size=5pt},
  hueco/.style={circle,draw=colHueco,fill=white,line width=0.9pt,inner sep=0pt,minimum size=5pt},
  foton/.style={colFoton,line width=0.9pt,decorate,
    decoration={snake,amplitude=1.6pt,segment length=5pt,post length=3pt},-{Stealth[length=5pt]}},
  flecha/.style={-{Stealth[length=5pt]},line width=0.7pt},
}

\unidad{Unidad 11 -- Dispositivos de potencia}

% ---------- Macros propias de la unidad ----------
\newcommand{\Ecr}{\mathcal{E}_c}          % campo eléctrico crítico (ruptura)
\newcommand{\Ron}{R_{\mathrm{on}}}
\newcommand{\Ronsp}{R_{\mathrm{on,sp}}}
\newcommand{\VB}{V_{B}}
\newcommand{\VAK}{V_{AK}}
\newcommand{\IA}{I_{A}}
\newcommand{\IG}{I_{G}}
\newcommand{\RDSon}{R_{DS(\mathrm{on})}}
\newcommand{\Qg}{Q_{g}}
\newcommand{\Qrr}{Q_{rr}}
\newcommand{\trr}{t_{rr}}
\newcommand{\IRM}{I_{RM}}
\newcommand{\boff}{\beta_{\mathrm{apag}}}
\newcommand{\taueff}{\tau_{\mathrm{ef}}}

% ---------- Valores de literatura para materiales WBG ----------
% (no están en la tabla oficial de la materia; se toman de la diapo 29
%  y de Baliga, "Fundamentals of Power Semiconductor Devices")
\newcommand{\EcSi}{\num{3e5}}          % V/cm
\newcommand{\EcSiC}{\num{2e6}}         % V/cm (4H-SiC, diapo 29)
\newcommand{\EcGaN}{\num{3.3e6}}       % V/cm
\newcommand{\epsrSiC}{9.7}
\newcommand{\epsrGaN}{9.0}
\newcommand{\munSiC}{720}              % cm^2/(V s) (diapo 29)
\newcommand{\munGaN}{900}
\newcommand{\EGSiC}{3.26}
\newcommand{\EGGaN}{3.39}

\title{\color{azul}\textbf{Semiconductores y dispositivos de potencia}\\[0.3em]
  \large Diodo de potencia, tiristor, triac, GTO, MOSFET, IGBT y dispositivos de banda ancha (SiC, GaN)}
\author{Física Electrónica -- Resumen de la unidad 11}
\date{Basado en la clase ``Semiconductores y Dispositivos de Potencia'' (Ing.\ E.\ D.\ Báez), en
Pierret, \emph{Semiconductor Device Fundamentals} (caps.~6, 8, 13 y 18), en Mohan--Undeland--Robbins,
\emph{Electrónica de potencia} (caps.~20--25) y en Baliga, \emph{Fundamentals of Power Semiconductor Devices}}

\begin{document}
\maketitle
\thispagestyle{empty}

\begin{tcolorbox}[base,colback=white,colframe=gris,title={Cómo leer este apunte}]
\small
\begin{itemize}
  \item Cajas \textcolor{azul}{\textbf{azules}}: definiciones; \textcolor{verde}{\textbf{verdes}}: resultados
  clave; \textcolor{naranja}{\textbf{naranjas}}: intuición física; \textcolor{violeta}{\textbf{violetas}}: ejemplos;
  \textcolor{rojo}{\textbf{rojas}} (\emph{Ojo}): errores o imprecisiones de las diapositivas y confusiones típicas.
  \item \diapo{n} indica la diapositiva $n$ de la presentación de la unidad.
  \item El sello \enhojatexto{} indica que el resultado figura en la hoja de fórmulas oficial. La hoja
  no tiene una sección de potencia: lo que se usa de ella son los resultados de juntura, BJT
  (Ebers--Moll, $\alpha$, $I_{CB0}$) y MOSFET, sobre los que se construye toda la unidad.
  \item Los ejemplos con silicio usan la \emph{tabla de constantes} de la materia. Para SiC y GaN, que
  no están en la tabla, se usan los valores de la diapo~29 (aclarado en cada caso).
  \item Notación del tiristor (igual que las diapositivas y Mohan): capas $p_1n_1p_2n_2$ desde el ánodo,
  junturas $J_1$ ($p_1n_1$), $J_2$ ($n_1p_2$), $J_3$ ($p_2n_2$); $Q_1$ es el PNP ($p_1n_1p_2$) con
  ganancia $\alpha_1$ y $Q_2$ el NPN ($n_1p_2n_2$) con ganancia $\alpha_2$.
\end{itemize}
\end{tcolorbox}

\tableofcontents
\newpage

% =====================================================================
\section{Introducción: el semiconductor como llave}
\label{sec:intro}
% =====================================================================

En electrónica de potencia los semiconductores casi nunca se usan como amplificadores lineales: se
usan como \textbf{llaves} que conectan y desconectan una carga de una fuente, miles o millones de
veces por segundo (convertidores CC--CC, inversores, rectificadores controlados, variadores de
motores). Todos los dispositivos de la unidad \diapo{1} son distintas respuestas a la misma
pregunta: \emph{¿cómo hacer una llave que bloquee mucha tensión, conduzca mucha corriente con poca
caída, y pase de un estado al otro rápido y con poca energía de control?}

\begin{definicion}[Llave ideal]
\begin{itemize}
  \item \textbf{Apagada (bloqueo)}: soporta cualquier tensión, directa o inversa, con corriente nula.
  \item \textbf{Encendida (conducción)}: conduce cualquier corriente con caída de tensión nula.
  \item \textbf{Conmutación}: instantánea, y controlada con una señal de potencia despreciable.
\end{itemize}
Como $p=vi$ es siempre cero, una llave ideal no disipa.
\end{definicion}

Un dispositivo real se aparta en los tres puntos, y de ahí salen las tres fuentes de pérdidas:
\begin{itemize}
  \item \textbf{Pérdidas de conducción}: $P_{\mathrm{cond}}=V_{\mathrm{on}}I$ (o $\Ron I^2$). Dependen de
  cómo se \emph{conduce}: con un solo tipo de portador (MOSFET, Schottky) o con inyección de ambos
  (diodo PiN, tiristor, IGBT).
  \item \textbf{Pérdidas de conmutación}: durante la transición $v$ e $i$ coexisten (figura~\ref{fig:conmutacion}).
  Con conmutación lineal de duración $t_{\mathrm{c,enc}}$ y $t_{\mathrm{c,apag}}$ a frecuencia $f_s$,
  \[
     P_{\mathrm{conm}} \approx \tfrac12\,V_d I_o\,f_s\,(t_{\mathrm{c,enc}}+t_{\mathrm{c,apag}}).
  \]
  Crecen con la frecuencia y con la carga almacenada que hay que meter o sacar del dispositivo.
  \item \textbf{Pérdidas de bloqueo y de control}: fugas en bloqueo ($V I_{\mathrm{fuga}}$) y la
  potencia del circuito de disparo.
\end{itemize}

\begin{figure}[H]
\centering
\begin{tikzpicture}[x=1cm,y=1cm,>=Stealth]
  % ejes
  \draw[->] (0,0) -- (11,0) node[right] {$t$};
  \draw[->] (0,0) -- (0,3.1);
  % tensión
  \draw[rojo,line width=1pt] (0,2.4) -- (1.5,2.4) -- (2.7,0.12) -- (6.5,0.12) -- (7.7,2.4) -- (10.5,2.4);
  \node[rojo,right] at (10.5,2.4) {$v_T$};
  % corriente
  \draw[azul,line width=1pt] (0,0.05) -- (1.5,0.05) -- (2.7,2.0) -- (6.5,2.0) -- (7.7,0.05) -- (10.5,0.05);
  \node[azul,right] at (10.5,0.35) {$i_T$};
  \draw[dashed,gris] (0,2.0) node[left,black] {$I_o$} -- (1.5,2.0);
  \node[left] at (0,2.4) {$V_d$};
  % potencia
  \begin{scope}[yshift=-3cm]
    \draw[->] (0,0) -- (11,0) node[right] {$t$};
    \draw[->] (0,0) -- (0,2.2) node[left] {$p_T$};
    \fill[naranja!35] (1.5,0) .. controls (1.9,1.6) and (2.3,1.6) .. (2.7,0.1) -- (2.7,0) -- cycle;
    \fill[naranja!35] (6.5,0) .. controls (6.9,1.6) and (7.3,1.6) .. (7.7,0.1) -- (7.7,0) -- cycle;
    \draw[naranja!80!black,line width=1pt] (0,0.02) -- (1.5,0.02) .. controls (1.9,1.6) and (2.3,1.6) .. (2.7,0.25)
       -- (6.5,0.25) .. controls (6.9,1.6) and (7.3,1.6) .. (7.7,0.02) -- (10.5,0.02);
    \node[font=\small,align=center] at (2.1,1.75) {$W_{\mathrm{enc}}$};
    \node[font=\small,align=center] at (7.1,1.75) {$W_{\mathrm{apag}}$};
    \node[font=\small] at (4.6,0.55) {$V_{\mathrm{on}}I_o$ (conducción)};
    \draw[<->] (1.5,-0.3) -- node[below,font=\small] {$t_{\mathrm{c,enc}}$} (2.7,-0.3);
    \draw[<->] (6.5,-0.3) -- node[below,font=\small] {$t_{\mathrm{c,apag}}$} (7.7,-0.3);
  \end{scope}
\end{tikzpicture}
\caption{Conmutación de una carga inductiva idealizada. En cada transición la potencia instantánea
$p_T=v_Ti_T$ llega a valores del orden de $V_dI_o/4$; la energía de cada pico ($W_{\mathrm{enc}}$,
$W_{\mathrm{apag}}$) se pierde $f_s$ veces por segundo.}
\label{fig:conmutacion}
\end{figure}

\subsection{Dos familias: unipolares y bipolares}

La física de los dispositivos de potencia se ordena alrededor de un compromiso que se deriva en la
sección~\ref{sec:deriva}: \textbf{para bloquear mucha tensión hace falta una región gruesa y poco
dopada} (la \emph{región de arrastre} o \emph{de deriva}, $n^-$), y esa región, en conducción, es una
resistencia enorme. Hay dos maneras de vivir con ella:

\begin{table}[H]
\centering\small
\begin{tabularx}{\textwidth}{>{\raggedright\arraybackslash}p{2.7cm}>{\raggedright\arraybackslash}X>{\raggedright\arraybackslash}X}
\toprule
 & \textbf{Unipolares} (portadores mayoritarios) & \textbf{Bipolares} (inyección de minoritarios) \\
\midrule
Ejemplos & Schottky, MOSFET, JFET, HEMT & Diodo PiN, BJT, tiristor, GTO, IGBT \\
Región $n^-$ en conducción & Resistencia óhmica fija: $\Ron\propto \VB^2$ & Se \emph{inunda} de electrones y huecos
($\delta n=\delta p\gg N_D$): \textbf{modulación de conductividad} \\
Caída en conducción & $\Ron I$ (chica a baja tensión, enorme a alta tensión) & $\approx$ una juntura
($\sim\SI{1}{V}$) $+$ algo; casi independiente de $\VB$ \\
Velocidad & Muy rápida (sólo cargar capacidades) & Lenta: hay que \emph{sacar} la carga almacenada
($Q\approx I\tau$) \\
Rango natural & $\lesssim\SI{600}{V}$ en Si (hasta kV en SiC) & De cientos de V a $\sim\SI{10}{kV}$ \\
\bottomrule
\end{tabularx}
\caption{Las dos familias. Los materiales de banda ancha (SiC, GaN; secciones~\ref{sec:wbg} y
\ref{sec:gan}) desplazan el límite de los unipolares a tensiones mucho mayores.}
\label{tab:uni-bi}
\end{table}

\subsection{Mapa de los dispositivos}

\begin{figure}[H]
\centering
\begin{tikzpicture}
\begin{axis}[
  width=0.92\textwidth, height=7.2cm,
  xmode=log, ymode=log,
  xmin=10, xmax=3e7, ymin=10, ymax=1e9,
  xlabel={Frecuencia de conmutación $f_s$ (Hz)},
  ylabel={Potencia controlada (VA)},
  grid=major, grid style={gris!20},
  tick label style={font=\small},
  every node/.style={fill opacity=1,text opacity=1},
]
  \fill[rojo!25,rounded corners=6pt,fill opacity=0.65]    (axis cs:20,1e5)  rectangle (axis cs:1e3,4e8);
  \node[rojo!70!black,font=\small\bfseries] at (axis cs:140,8e7) {Tiristor};
  \fill[naranja!30,rounded corners=6pt,fill opacity=0.65] (axis cs:80,1e5)  rectangle (axis cs:3e3,6e7);
  \node[naranja!70!black,font=\small\bfseries] at (axis cs:500,3e6) {GTO / IGCT};
  \fill[azul!20,rounded corners=6pt,fill opacity=0.65]    (axis cs:500,3e2) rectangle (axis cs:1e5,3e7);
  \node[azul,font=\small\bfseries] at (axis cs:7e3,2e5) {IGBT};
  \fill[verde!25,rounded corners=6pt,fill opacity=0.65]   (axis cs:1e4,10)  rectangle (axis cs:2e6,2e5);
  \node[verde!60!black,font=\small\bfseries] at (axis cs:3e4,2.5e1) {MOSFET (Si)};
  \fill[violeta!25,rounded corners=6pt,fill opacity=0.65] (axis cs:2e4,1e3) rectangle (axis cs:1e6,3e6);
  \node[violeta,font=\small\bfseries] at (axis cs:1.5e5,8e5) {MOSFET SiC};
  \fill[gris!25,rounded corners=6pt,fill opacity=0.65]    (axis cs:1e5,10)  rectangle (axis cs:2e7,3e4);
  \node[gris!80!black,font=\small\bfseries] at (axis cs:2e6,1e4) {HEMT GaN};
\end{axis}
\end{tikzpicture}
\caption{Zonas de aplicación típicas (orientativas; los bordes se mueven año a año). A mayor
potencia, dispositivos bipolares lentos; a mayor frecuencia, unipolares. El SiC y el GaN empujan a los
unipolares hacia arriba y a la derecha.}
\label{fig:mapa}
\end{figure}

\begin{definicion}[Área de operación segura (SOA)]
Región del plano $(V,I)$ por la que puede pasar la trayectoria de conmutación sin destruir el
dispositivo. Está limitada por: la corriente máxima (hilos, metalización), la potencia máxima
(temperatura de juntura), la \emph{segunda ruptura} en bipolares (concentración de corriente con
realimentación térmica) y la tensión de ruptura por avalancha. Los circuitos \emph{snubber} existen
para mantener la trayectoria dentro de la SOA y limitar $\dd v/\dd t$ y $\dd i/\dd t$.
\end{definicion}

% =====================================================================
\section{Bloqueo de tensión y región de arrastre}
\label{sec:deriva}
% =====================================================================

Todos los dispositivos de potencia verticales tienen la misma columna vertebral: una juntura
$p^+n^-$ (o un contacto Schottky, o un canal MOS) seguida de una región $n^-$ gruesa y poco dopada y
de un sustrato $n^+$ que hace de contacto \diapo{2, 20}. La $n^-$ es la \textbf{región de arrastre
de drenaje} (\emph{drift region}): es la que \emph{soporta la tensión} en bloqueo.

\subsection{Ruptura por avalancha: el campo crítico}

En inversa el campo eléctrico es máximo en la juntura. Cuando llega al \textbf{campo crítico}
$\Ecr$ los portadores ganan entre choques energía suficiente para ionizar la red y la corriente
crece en avalancha (Pierret §6.3.1). En lo que sigue se usa el modelo simple ``ruptura cuando
$\mathcal{E}_{\max}=\Ecr$'' con $\Ecr$ constante: para Si $\Ecr\approx\SI{3e5}{V/cm}$, para 4H-SiC
$\Ecr\approx\SI{2e6}{V/cm}$ \diapo{29}.

\begin{intuicion}[¿Por qué un semiconductor de banda ancha aguanta más campo?]
Para ionizar hace falta que un portador gane del orden de $\EG$ (en rigor $\sim1.5\EG$) entre choques.
Con $\EG$ tres veces mayor, el campo necesario es mucho mayor: $\Ecr$ crece aproximadamente como
$\EG^{2}$ a $\EG^{2.5}$. Ésa es la principal razón de ser del SiC y del GaN en potencia.
\end{intuicion}

\subsection{Diseño sin perforación (NPT): campo triangular}

Juntura $p^+n^-$ abrupta con $N_A\gg N_D$: la zona desierta queda casi toda en la $n^-$, de ancho
$W$. De Poisson (hoja de fórmulas, juntura M--S o $p^+n$ \enhojatexto) el campo es lineal:
\[
  \mathcal{E}(x)=\frac{qN_D}{\epsilon}(W-x),\qquad \mathcal{E}_{\max}=\frac{qN_DW}{\epsilon},
  \qquad V=\frac12\,\mathcal{E}_{\max}W .
\]
Imponiendo $\mathcal{E}_{\max}=\Ecr$ en la ruptura:

\begin{clave}[Tensión de ruptura y ancho de la región de arrastre]
\[
  \VB=\frac{\epsilon\,\Ecr^2}{2qN_D},\qquad
  W_{\VB}=\frac{2\VB}{\Ecr},\qquad
  N_D=\frac{\epsilon\,\Ecr^2}{2q\VB}.
\]
Para bloquear más tensión hay que \textbf{dopar menos} ($N_D\propto1/\VB$) y \textbf{engrosar}
($W\propto\VB$) la región de arrastre. Ésta es la ``región de arrastre del espesor $W_d$'' que
depende de la tensión de ruptura en el diodo de la diapo~2.
\end{clave}

\subsection{La resistencia en conducción: el límite unipolar}

Si en conducción esa región sólo tiene sus electrones de equilibrio ($n=N_D$), su resistencia por
unidad de área es
\[
  \Ronsp=\frac{W}{q\mu_nN_D}
  =\frac{2\VB/\Ecr}{q\mu_n\,\epsilon\Ecr^2/(2q\VB)}
  \quad\Longrightarrow\quad
  \boxed{\;\Ronsp=\frac{4\VB^2}{\epsilon\,\mu_n\,\Ecr^3}\;}
\]
que es la expresión de la diapo~30. El denominador $\epsilon\mu_n\Ecr^3$ es la \textbf{figura de
mérito de Baliga} (BFOM) del material: depende sólo del semiconductor, y aparece al cubo el campo
crítico.

\begin{clave}[Límite unipolar]
$\Ronsp\propto\VB^2$ (y en la práctica $\propto\VB^{2.5}$, porque $\Ecr$ baja un poco con $N_D$).
Duplicar la tensión de bloqueo cuadruplica la resistencia: por eso el MOSFET de silicio deja de ser
competitivo arriba de $\sim\SIrange{200}{600}{V}$ y aparecen el IGBT (modulación de conductividad) y
el SiC (más $\Ecr$).
\end{clave}

\begin{ejemplo}[Región de arrastre para \SI{600}{V} en silicio]
Con $\epsilon=\epsrSi\,\epsilon_0$, $\Ecr=\EcSi\,\si{V/cm}$, $\mu_n=\munSi\,\si{cm^2/V.s}$:
\[
  N_D=\frac{11.8\cdot\num{8.85e-14}\cdot(\num{3e5})^2}{2\cdot\num{1.6e-19}\cdot600}
  =\SI{4.9e14}{cm^{-3}},\qquad
  W=\frac{2\cdot600}{\num{3e5}}=\SI{40}{\micro m},
\]
\[
  \Ronsp=\frac{W}{q\mu_nN_D}=\frac{\num{40e-4}}{\num{1.6e-19}\cdot1350\cdot\num{4.9e14}}
  =\SI{0.038}{\ohm.cm^2}.
\]
Un MOSFET de \SI{1}{cm^2} de área activa tendría, sólo por la deriva, $\SI{38}{m\ohm}$:
con \SI{30}{A} disipa $\SI{34}{W}$. Con 4H-SiC la misma tensión se logra con
$\Ronsp\approx\SI{0.29}{m\ohm.cm^2}$, unas 130 veces menos (sección~\ref{sec:wbg}).
\end{ejemplo}

\subsection{Diseño con perforación (PT): campo trapezoidal}

Si la $n^-$ se hace \emph{más delgada} que $W_{\VB}$ y se apoya en una capa $n^+$, la zona desierta
``perfora'' hasta la $n^+$ antes de la ruptura y el campo queda trapezoidal
(figura~\ref{fig:pt-npt}). La tensión es el área bajo la curva:
\[
  \VB^{PT}=\Ecr\,W_d-\frac{qN_DW_d^2}{2\epsilon}.
\]
Con el mismo $\Ecr$ se bloquea casi la misma tensión con una región más delgada (y por lo tanto
menos resistiva y con menos carga almacenada en dispositivos bipolares). En el límite $N_D\to0$
(capa ``intrínseca'') el campo es rectangular y $\VB=\Ecr W_d$: es la estructura PiN.

\begin{figure}[H]
\centering
\begin{tikzpicture}[x=1cm,y=1cm,>=Stealth]
  \foreach \dx/\tit in {0/{(a) NPT: $n^-$ gruesa}, 7.6/{(b) PT: $n^-$ delgada $+$ $n^+$}} {
    \begin{scope}[xshift=\dx cm]
      \node[font=\small\bfseries] at (2.9,3.6) {\tit};
      \draw[->] (0,0) -- (6.4,0) node[right] {$x$};
      \draw[->] (0,0) -- (0,3.1) node[left] {$\mathcal{E}$};
      \draw[dashed,gris] (0,2.6) node[left,black,font=\small] {$\Ecr$} -- (6,2.6);
    \end{scope}
  }
  % NPT
  \fill[rojo!15] (0.4,0) rectangle (0.8,-0.5);
  \fill[azul!8] (0.8,0) rectangle (6,-0.5);
  \node[font=\scriptsize] at (0.6,-0.25) {$p^+$};
  \node[font=\scriptsize] at (3.4,-0.25) {$n^-$};
  \fill[naranja!25] (0.8,0) -- (0.8,2.6) -- (5.2,0) -- cycle;
  \draw[naranja!80!black,line width=1pt] (0.8,0) -- (0.8,2.6) -- (5.2,0);
  \node[font=\small] at (2.0,0.8) {$\VB$};
  \draw[<->] (0.8,-0.75) -- node[below,font=\small] {$W_{\VB}=2\VB/\Ecr$} (5.2,-0.75);
  % PT
  \begin{scope}[xshift=7.6cm]
    \fill[rojo!15] (0.4,0) rectangle (0.8,-0.5);
    \fill[azul!8] (0.8,0) rectangle (3.6,-0.5);
    \fill[azul!30] (3.6,0) rectangle (4.4,-0.5);
    \node[font=\scriptsize] at (0.6,-0.25) {$p^+$};
    \node[font=\scriptsize] at (2.2,-0.25) {$n^-$};
    \node[font=\scriptsize] at (4.0,-0.25) {$n^+$};
    \fill[naranja!25] (0.8,0) -- (0.8,2.6) -- (3.6,1.0) -- (3.75,0) -- cycle;
    \draw[naranja!80!black,line width=1pt] (0.8,0) -- (0.8,2.6) -- (3.6,1.0) -- (3.75,0);
    \node[font=\small] at (2.0,0.8) {$\VB$};
    \draw[<->] (0.8,-0.75) -- node[below,font=\small] {$W_d<W_{\VB}$} (3.6,-0.75);
  \end{scope}
\end{tikzpicture}
\caption{Perfil de campo en la ruptura. El área sombreada es la tensión bloqueada. En el diseño PT la
pendiente $qN_D/\epsilon$ es la misma pero el triángulo se ``corta'' en la capa $n^+$ (que en el IGBT
se llama capa de compensación o \emph{buffer}, diapo~22).}
\label{fig:pt-npt}
\end{figure}

\begin{figure}[H]
\centering
\begin{tikzpicture}
\begin{axis}[
  width=0.78\textwidth, height=6.4cm,
  xmode=log, ymode=log,
  xmin=100, xmax=1e4, ymin=1e-2, ymax=1e4,
  xlabel={Tensión de ruptura $\VB$ (V)},
  ylabel={$\Ronsp$ (m$\Omega\cdot$cm$^2$)},
  grid=major, grid style={gris!20},
  legend pos=north west, legend style={font=\small},
  domain=100:1e4, samples=40,
]
  % Ron,sp = 4 V^2/(eps mu Ec^3), en mOhm cm^2 -> *1e3
  \addplot[rojo,line width=1.1pt]    {1e3*4*x^2/(11.8*8.85e-14*1350*(3e5)^3)};
  \addlegendentry{Si}
  \addplot[azul,line width=1.1pt]    {1e3*4*x^2/(9.7*8.85e-14*720*(2e6)^3)};
  \addlegendentry{4H-SiC}
  \addplot[verde!60!black,line width=1.1pt] {1e3*4*x^2/(9.0*8.85e-14*900*(3.3e6)^3)};
  \addlegendentry{GaN}
\end{axis}
\end{tikzpicture}
\caption{Límite unipolar $\Ronsp=4\VB^2/(\epsilon\mu_n\Ecr^3)$ con los parámetros de la diapo~29.
Un dispositivo unipolar no puede estar por debajo de su recta (salvo trucos de estructura como la
\emph{superjuntura}, sección~\ref{sec:mosfet}).}
\label{fig:limite-unipolar}
\end{figure}

% =====================================================================
\section{Diodo de potencia}
\label{sec:diodo}
% =====================================================================

\subsection{Estructura $p^+n^-n^+$ (PiN)}

\begin{figure}[H]
\centering
\begin{tikzpicture}[x=1cm,y=1cm,>=Stealth]
  \fill[rojo!20] (0,3.2) rectangle (6,4.0);
  \fill[azul!8]  (0,1.2) rectangle (6,3.2);
  \fill[azul!30] (0,0)   rectangle (6,1.2);
  \draw (0,0) rectangle (6,4.0);
  \draw (0,3.2) -- (6,3.2); \draw (0,1.2) -- (6,1.2);
  \fill[gris] (0,4.0) rectangle (6,4.12); \fill[gris] (0,-0.12) rectangle (6,0);
  \node at (1.0,3.6) {$p^+$}; \node[font=\small] at (3.9,3.6) {$N_A=\SI{1e19}{cm^{-3}}$};
  \node at (1.0,2.2) {$n^-$ (epi)}; \node[font=\small] at (3.9,2.2) {$N_D=\SI{1e14}{cm^{-3}}$};
  \node at (1.0,0.6) {$n^+$ (sustrato)}; \node[font=\small] at (3.9,0.6) {$N_D=\SI{1e19}{cm^{-3}}$};
  \draw (3,4.12) -- (3,4.6) node[circle,draw,inner sep=1.3pt,fill=white] {} node[right=3pt] {Ánodo};
  \draw (3,-0.12) -- (3,-0.6) node[circle,draw,inner sep=1.3pt,fill=white] {} node[right=3pt] {Cátodo};
  \draw[<->] (6.3,3.2) -- node[right,font=\small] {$\sim\SI{10}{\micro m}$} (6.3,4.0);
  \draw[<->] (6.3,1.2) -- node[right,font=\small,align=left] {$W_d$: depende de $\VB$\\(región de arrastre)} (6.3,3.2);
  \draw[<->] (6.3,0) -- node[right,font=\small] {$\sim\SI{250}{\micro m}$ (soporte mecánico)} (6.3,1.2);
\end{tikzpicture}
\caption{Diodo de potencia típico \diapo{2}. La $n^+$ inferior no bloquea nada: sólo da rigidez
mecánica y un buen contacto óhmico. Todo el diseño está en $N_D$ y $W_d$ de la capa epitaxial.}
\label{fig:diodo-estructura}
\end{figure}

Comparado con un diodo de señal, el de potencia tiene dos diferencias: la región $n^-$ (para
bloquear, sección~\ref{sec:deriva}) y la gran área (para conducir decenas o miles de amperes). La
primera es un problema en conducción: con $N_D=\SI{1e14}{cm^{-3}}$ y $W_d=\SI{100}{\micro m}$ la $n^-$
\emph{sin modular} tendría
\[
  \Ronsp=\frac{W_d}{q\mu_nN_D}=\frac{\num{1e-2}}{\num{1.6e-19}\cdot1350\cdot\num{1e14}}=\SI{0.46}{\ohm.cm^2},
\]
es decir \SI{46}{V} de caída a \SI{100}{A/cm^2}. Los diodos reales caen $\sim\SIrange{1}{2}{V}$:
la diferencia la hace la \textbf{alta inyección}.

\subsection{Baja y alta inyección; inyección doble}

\begin{itemize}
  \item \textbf{Baja inyección} ($\delta p\ll n_{n0}=N_D$): los huecos que inyecta la $p^+$ son pocos;
  los electrones de equilibrio de la $n^-$ neutralizan su carga sin dificultad \diapo{3}. Es el diodo
  de la unidad de junturas.
  \item \textbf{Alta inyección} ($\delta p\gtrsim N_D$): la carga de los huecos inyectados ya no puede
  ser neutralizada por los pocos electrones de la $n^-$. Esa carga positiva \emph{atrae electrones
  desde la $n^+$}: la juntura $n^+n^-$ inyecta electrones en la región de arrastre, con
  $\delta n=\delta p$ (neutralidad). Es la \textbf{inyección doble} \diapo{3}.
\end{itemize}

\begin{figure}[H]
\centering
\begin{tikzpicture}
\begin{axis}[
  width=0.82\textwidth, height=6cm,
  ymode=log, ymin=1e4, ymax=1e20,
  xmin=-0.3, xmax=1.3,
  xtick={0,1}, xticklabels={$0$,$W_d$},
  xlabel={$x$ en la región $n^-$},
  ylabel={concentración (cm$^{-3}$)},
  ytick={1e6,1e10,1e14,1e16,1e19},
  axis lines=left,
  every axis plot/.append style={line width=1.1pt},
  legend style={font=\small,at={(0.5,0.4)},anchor=center},
]
  % zonas
  \fill[rojo!10] (axis cs:-0.3,1e4) rectangle (axis cs:0,1e20);
  \fill[azul!12] (axis cs:1,1e4) rectangle (axis cs:1.3,1e20);
  \node[font=\small] at (axis cs:-0.15,3e19) {$p^+$};
  \node[font=\small] at (axis cs:1.15,3e19) {$n^+$};
  % alta inyección: perfil catenaria n=p
  \addplot[naranja!80!black,domain=0:1,samples=60] {3e17*cosh((x-0.5)/0.2)/cosh(0.5/0.2)};
  \addlegendentry{$n(x)\approx p(x)$ en alta inyección}
  % equilibrio
  \addplot[azul,dashed,domain=0:1] {1e14};
  \addlegendentry{$n_{n0}=N_D=10^{14}$}
  \addplot[rojo,dashed,domain=0:1] {1e6};
  \addlegendentry{$p_{n0}=n_i^2/N_D=10^{6}$}
\end{axis}
\end{tikzpicture}
\caption{Portadores en la región de arrastre \diapo{3}. En conducción a corriente alta, $n\approx p$
es tres órdenes de magnitud mayor que el dopaje: la $n^-$ deja de ser una resistencia de
\SI{0.46}{\ohm.cm^2}. El perfil es una catenaria (se deriva abajo).}
\label{fig:pin-perfil}
\end{figure}

\subsection{Modulación de conductividad}

Con $n\approx p\gg N_D$ la conductividad de la región de arrastre es
$\sigma=q(\mu_n+\mu_p)n$, controlada por la \emph{corriente} y no por el dopaje.

\begin{ejemplo}[Cuánto baja la resistencia]
Con $n=p=\SI{1e17}{cm^{-3}}$ y las movilidades de la tabla ($\mu_n=\munSi$, $\mu_p=\mupSi$):
\[
  \rho=\frac{1}{q(\mu_n+\mu_p)n}=\frac{1}{\num{1.6e-19}\cdot1830\cdot\num{1e17}}=\SI{0.034}{\ohm.cm},
  \qquad
  \rho W_d=\SI{3.4e-4}{\ohm.cm^2}.
\]
A \SI{100}{A/cm^2}: \SI{34}{mV} en vez de \SI{46}{V}. (Las movilidades bajan algo en alta
inyección por la dispersión portador--portador, pero el orden de magnitud no cambia.)
\end{ejemplo}

\paragraph{Perfil de portadores: difusión ambipolar.}
En la $n^-$, con $\delta n=\delta p$ y el campo determinado por la neutralidad, las dos ecuaciones
de continuidad se combinan en una sola, para el ``par'' electrón--hueco, con un coeficiente de
\textbf{difusión ambipolar}:
\[
  D_a\frac{\dd^2\delta p}{\dd x^2}-\frac{\delta p}{\tau_{HL}}=0,
  \qquad D_a=\frac{2D_nD_p}{D_n+D_p},\qquad L_a=\sqrt{D_a\tau_{HL}},
\]
donde $\tau_{HL}$ es la vida media en alta inyección (sección~\ref{sec:recomb}). La solución con
inyección desde los dos extremos es simétrica:
\[
  \delta p(x)=\delta n(x)\propto\cosh\!\left(\frac{x-W_d/2}{L_a}\right),
\]
la ``catenaria'' de la figura~\ref{fig:pin-perfil}. Si $W_d\lesssim 2L_a$ el perfil es casi plano y la
región está bien modulada; si $W_d\gg L_a$ el centro queda pobre en portadores y la caída crece
rápidamente. Por eso \textbf{una vida media larga baja la caída en conducción}.

\begin{clave}[Caída en directa del diodo de potencia]
\[
  V=V_j+V_d\;\approx\;V_j+R_{\mathrm{enc}}\,I \qquad\text{\diapo{3}}
\]
$V_j$: caída en las junturas $p^+n^-$ y $n^-n^+$ (logarítmica en $I$); $V_d$: caída en la región de
arrastre modulada, aproximadamente óhmica a corriente alta. En la característica $I$--$V$ aparece como
una recta de pendiente $1/R_{\mathrm{enc}}$ a partir de $\sim\SI{1}{V}$.
\end{clave}

\subsection{Recuperación inversa}

La conductividad modulada tiene un precio: en conducción la región de arrastre almacena una carga
$Q_F\approx I_F\,\tau_{HL}$ de portadores en exceso (control de carga, Pierret cap.~8). Para que el
diodo bloquee, esa carga tiene que \emph{desaparecer} --- por recombinación o extraída como corriente
inversa. Mientras haya carga cerca de la juntura, el diodo sigue ``encendido'' en sentido inverso.

\begin{figure}[H]
\centering
\begin{tikzpicture}[x=1cm,y=1cm,>=Stealth]
  \draw[->] (0,0) -- (11,0) node[right] {$t$};
  \draw[->] (0,-2.6) -- (0,2.4) node[left] {$i_D$};
  % corriente
  \draw[azul,line width=1.1pt] (0,1.8) -- (1.5,1.8) -- (4.2,-1.8)
     .. controls (4.6,-2.1) and (4.9,-0.9) .. (5.8,-0.2) .. controls (6.4,0) .. (10.5,0);
  \node[azul,left] at (0,1.8) {$I_F$};
  \draw[dashed,gris] (0,-1.98) node[left,black] {$-\IRM$} -- (4.4,-1.98);
  \node[font=\small,rotate=-53] at (3.1,0.6) {$-\dd i/\dd t$};
  % tiempos
  \draw[dashed,gris] (2.85,0) -- (2.85,-2.6);
  \draw[dashed,gris] (4.4,0) -- (4.4,-2.6);
  \draw[dashed,gris] (6.2,0) -- (6.2,-2.6);
  \draw[<->] (2.85,-2.45) -- node[below,font=\small] {$t_a$} (4.4,-2.45);
  \draw[<->] (4.4,-2.45) -- node[below,font=\small] {$t_b$} (6.2,-2.45);
  \draw[<->] (2.85,-3.15) -- node[below,font=\small] {$\trr$} (6.2,-3.15);
  % área Qrr
  \fill[azul!25] (2.85,0) -- (4.2,-1.8) .. controls (4.6,-2.1) and (4.9,-0.9) .. (5.8,-0.2) .. controls (6.0,-0.1) .. (6.2,0) -- cycle;
  \node[font=\small] at (4.6,-0.6) {$\Qrr$};
  % tensión
  \draw[rojo,line width=1.1pt] (0,0.12) -- (4.3,0.12) .. controls (4.6,-1.0) and (4.9,-1.6) .. (5.3,-1.45)
      .. controls (6.0,-1.2) .. (10.5,-1.2);
  \node[rojo,right] at (10.5,-1.2) {$v_D=-V_R$};
  \node[rojo,font=\small,align=left] at (8.2,-0.6) {sobretensión\\por $L\,\dd i/\dd t$};
\end{tikzpicture}
\caption{Apagado de un diodo de potencia con $\dd i/\dd t$ impuesto por el circuito ($\approx V_R/L$).
Durante $t_a$ la juntura sigue polarizada en directa (todavía hay carga en su borde): la corriente
inversa sólo está limitada por el circuito. Cuando la carga cerca de la juntura se agota, la juntura
toma tensión y la corriente se recupera en $t_b$.}
\label{fig:recup-inversa}
\end{figure}

\begin{definicion}[Parámetros de recuperación inversa]
\begin{itemize}
  \item $\trr=t_a+t_b$: tiempo de recuperación inversa.
  \item $\Qrr$: carga recuperada (área de la corriente negativa), $\Qrr\approx\tfrac12\IRM\trr$.
  \item $\IRM$: pico de corriente inversa. Si $t_b\ll t_a$: $t_a=\IRM/(\dd i/\dd t)$ y
  \[
    \IRM\approx\sqrt{2\,\Qrr\,\frac{\dd i}{\dd t}},\qquad \trr\approx\sqrt{\frac{2\Qrr}{\dd i/\dd t}}.
  \]
  \item Factor de suavidad $S=t_b/t_a$. Un diodo ``abrupto'' ($S\ll1$) corta la corriente de golpe y,
  con la inductancia del circuito, genera sobretensiones y ruido; uno ``suave'' ($S\sim1$) es más
  benigno.
\end{itemize}
\end{definicion}

\begin{ejemplo}[Pico de corriente inversa]
Diodo con $\Qrr=\SI{1}{\micro C}$, apagado con $\dd i/\dd t=\SI{100}{A/\micro s}$:
\[
  \IRM\approx\sqrt{2\cdot\num{1e-6}\cdot\num{1e8}}=\SI{14}{A},\qquad
  \trr\approx\frac{\IRM}{\dd i/\dd t}=\SI{0.14}{\micro s}.
\]
Esos \SI{14}{A} los entrega la llave que se está encendiendo (por ejemplo el MOSFET de la otra
rama de un puente), \emph{sumados} a la corriente de carga: la recuperación del diodo es la principal
pérdida de encendido de muchos convertidores.
\end{ejemplo}

Al \emph{encender} hay un fenómeno dual, la \textbf{recuperación directa}: la región de arrastre
todavía no está modulada, su resistencia es alta y la tensión tiene un sobrepico $V_{FR}$ de varios
volts hasta que la inyección doble la llena de portadores.

\subsection{Diodo Schottky de potencia}

Un contacto metal--semiconductor sobre la $n^-$ es un diodo \emph{unipolar}: conduce por emisión
termoiónica de mayoritarios (hoja de fórmulas: $I_S=AA^*T^2e^{-\phi_B/kT}$ \enhojatexto), no inyecta
huecos y por lo tanto \textbf{no tiene carga almacenada ni recuperación inversa} (sólo la carga de
la capacidad de juntura). A cambio no modula su región de arrastre: su caída sigue el límite
unipolar y en silicio sólo se usan hasta $\sim\SI{100}{V}$--$\SI{200}{V}$. En SiC, gracias al $\Ecr$
diez veces mayor, llegan a \SIrange{600}{1700}{V} (sección~\ref{sec:wbg}).

% =====================================================================
\section{Recombinación y vida media en dispositivos de potencia}
\label{sec:recomb}
% =====================================================================

En los dispositivos bipolares de potencia (diodo PiN, BJT, tiristor, GTO, IGBT) la vida media de los
portadores en exceso es \emph{el} parámetro de diseño: fija cuánto se modula la región de arrastre
(caída en conducción) y cuánta carga hay que sacar al apagar (velocidad y pérdidas de conmutación).

\subsection{Repaso: equilibrio y decaimiento de los portadores en exceso}

\diapo{4} En equilibrio valen la ley de acción de masas y la neutralidad \enhojatexto:
\[
  p_0n_0=n_i^2,\qquad p_0+N_D=n_0+N_A
  \quad\Longrightarrow\quad
  \text{tipo }p:\;\; p_0\approx N_A,\quad n_0\approx\frac{n_i^2}{N_A}.
\]
Fuera del equilibrio, con un exceso $\delta n$ y sin generación externa, la ecuación de
generación--recombinación térmica en baja inyección es
\[
  \frac{\dd(\delta n)}{\dd t}=-\frac{\delta n}{\tau}
  \quad\Longrightarrow\quad \delta n(t)=\delta n(0)\,e^{-t/\tau},
\]
con $\tau$ la vida media del portador minoritario en exceso.

\begin{ojo}[``El $\tau$ no es una constante'']
En baja inyección $\tau$ es una propiedad del material (trampas, temperatura). En un dispositivo de
potencia, en cambio, se trabaja con $\delta n$ que va de $\ll N_D$ (bloqueo) a $\gg N_D$
(conducción), y $\tau$ cambia con $\delta n$ \diapo{4}. Por eso la ecuación de arriba con $\tau$ fijo
sólo sirve para una corriente dada.
\end{ojo}

\subsection{Los tres mecanismos}

\paragraph{SRH (trampas).} Es el dominante en silicio con inyección moderada. Con trampas en el centro
de la banda, Pierret §3.3:
\[
  \tau_{SRH}=\tau_{p0}\frac{n_0+n_1+\delta n}{n_0+p_0+\delta n}+\tau_{n0}\frac{p_0+p_1+\delta n}{n_0+p_0+\delta n}
  \;\xrightarrow{\;\delta n\gg n_0\;}\;\tau_{HL}=\tau_{n0}+\tau_{p0}.
\]
En alta inyección la vida media \emph{aumenta} (de $\tau_{p0}$ a $\tau_{n0}+\tau_{p0}$): la
recombinación necesita capturar los dos portadores y ambos están en exceso.
\textbf{$\tau_{SRH}$ crece con la temperatura} \diapo{4} (aproximadamente $\propto T^{1.5}$--$T^2$):
al subir $T$ los portadores son emitidos térmicamente desde las trampas antes de completar la
recombinación. Consecuencia práctica: a mayor temperatura, más carga almacenada y apagados más
lentos y con más $\Qrr$.

\paragraph{Auger.} Tres partículas: un electrón y un hueco se recombinan y ceden la energía a un
tercer portador. La tasa es $R_A=(C_nn+C_pp)(np-n_i^2)$; en alta inyección ($n=p=\delta n$)
\[
  \tau_A=\frac{1}{C_a\,\delta n^2},\qquad C_a=C_n+C_p\approx\SI{1.66e-30}{cm^6/s}\text{ (Si)}.
\]
\begin{center}\small
\begin{tabular}{lcccc}
\toprule
$\delta n$ (\si{cm^{-3}}) & $10^{16}$ & $10^{17}$ & $10^{18}$ & $10^{19}$ \\
$\tau_A$ & \SI{6}{ms} & \SI{60}{\micro s} & \SI{0.6}{\micro s} & \SI{6}{ns} \\
\bottomrule
\end{tabular}
\end{center}
Como $\tau_{SRH}$ típico en potencia es de \SIrange{1}{10}{\micro s}, Auger domina recién por
encima de $\sim\SI{3e17}{cm^{-3}}$.

\paragraph{Radiativa.} Despreciable en Si y SiC (gap indirecto; unidad 6).

\begin{clave}[Vida media efectiva]
\[
  \frac{1}{\taueff}=\frac{1}{\tau_{SRH}}+\frac{1}{\tau_A}
  \qquad\text{modelo empírico de la diapo 4:}\qquad
  \tau=\frac{\tau_0}{1+(\delta n/n_b)^2}.
\]
Las dos expresiones tienen la misma forma: el término $(\delta n/n_b)^2$ es la Auger con
$n_b=1/\sqrt{C_a\tau_0}$. A densidades muy altas $\tau$ cae, la modulación empeora y las pérdidas en
conducción aumentan \diapo{4}: es lo que limita la densidad de corriente útil de un PiN o un tiristor.
\end{clave}

\begin{figure}[H]
\centering
\begin{tikzpicture}
\begin{axis}[
  width=0.75\textwidth, height=5.8cm,
  xmode=log, ymode=log,
  xmin=1e14, xmax=1e19, ymin=1e-8, ymax=1e-4,
  xlabel={$\delta n$ (cm$^{-3}$)}, ylabel={$\tau$ (s)},
  grid=major, grid style={gris!20},
  legend pos=south west, legend style={font=\small},
  domain=1e14:1e19, samples=60,
]
  \addplot[gris,dashed,line width=0.9pt] {1e-5};
  \addlegendentry{$\tau_{SRH}=\SI{10}{\micro s}$}
  \addplot[rojo,dashed,line width=0.9pt] {1/(1.66e-30*x^2)};
  \addlegendentry{$\tau_A=1/(C_a\delta n^2)$}
  \addplot[azul,line width=1.2pt] {1/(1e5+1.66e-30*x^2)};
  \addlegendentry{$\taueff$}
  \draw[dotted] (axis cs:2.45e17,1e-8) -- (axis cs:2.45e17,1e-4);
  \node[font=\small,anchor=west] at (axis cs:2.6e17,3e-5) {$n_b\approx\SI{2.5e17}{cm^{-3}}$};
\end{axis}
\end{tikzpicture}
\caption{Vida media efectiva en Si con $\tau_{SRH}=\SI{10}{\micro s}$. Para $\delta n<n_b$ manda SRH;
por encima, Auger reduce $\tau$ como $\delta n^{-2}$.}
\label{fig:tau-dn}
\end{figure}

\begin{ojo}[Dónde empieza la Auger \diapo{4}]
La diapositiva dice que $\tau$ disminuye por Auger ``para altas densidades cercanas a
$n_b\approx\SI{1e16}{cm^{-3}}$''. En silicio, con $C_a\approx\SI{1.7e-30}{cm^6/s}$, $\tau_A$ a
\SI{1e16}{cm^{-3}} es de milisegundos: no compite con SRH. El quiebre ocurre en
$n_b=1/\sqrt{C_a\tau_{SRH}}\approx\SIrange{1e17}{1e18}{cm^{-3}}$ según la vida SRH. En el modelo
empírico $n_b$ es un parámetro de ajuste, no una constante del material.
\end{ojo}

\subsection{El compromiso de la vida media}

\begin{intuicion}[Ni mucho ni poco]
\begin{itemize}
  \item $\tau$ \textbf{grande} $\Rightarrow$ $L_a$ grande, región de arrastre bien modulada,
  \textbf{baja caída en conducción}; pero mucha carga almacenada ($Q\approx I\tau$),
  \textbf{conmutación lenta} y gran $\Qrr$ \diapo{4}.
  \item $\tau$ \textbf{chico} $\Rightarrow$ apagado rápido, pero modulación pobre y \textbf{más
  pérdidas en conducción}.
\end{itemize}
El fabricante elige el punto de la curva $V_F$--$\Qrr$ (o $V_{CE(sat)}$--$E_{\mathrm{apag}}$ en IGBT)
según la frecuencia de trabajo. Un rectificador de red (\SI{50}{Hz}) usa $\tau$ largo; un diodo de
rueda libre de un inversor de \SI{20}{kHz}, $\tau$ corto.
\end{intuicion}

\paragraph{Control de vida media.} Se introducen deliberadamente centros de recombinación:
difusión de oro o platino (niveles profundos), irradiación con electrones (defectos en todo el
volumen) o con protones/helio (defectos \emph{localizados} a una profundidad elegida, por ejemplo cerca
de la juntura para mejorar la suavidad $S$ sin empeorar mucho $V_F$).

% =====================================================================
\section{Tiristor (SCR)}
\label{sec:tiristor}
% =====================================================================

El tiristor o \emph{rectificador controlado de silicio} (SCR) es una estructura de cuatro capas
$pnpn$ con tres terminales: ánodo (A), cátodo (K) y compuerta (G, conectada a la $p$ interna). Es un
diodo que \textbf{no conduce en directa hasta que se lo dispara} con un pulso de compuerta; una vez
encendido \emph{se queda encendido} aunque se quite el pulso, y sólo se apaga cuando la corriente de
ánodo cae por debajo de un mínimo (Pierret cap.~13; Mohan cap.~23).

\subsection{Estructura}

\begin{figure}[H]
\centering
\begin{tikzpicture}[x=1cm,y=1cm,>=Stealth]
  % capas (cátodo arriba, como en la diapo 8)
  \fill[azul!35] (0,4.6) rectangle (2.2,5.2);  \fill[azul!35] (3.8,4.6) rectangle (6.4,5.2);
  \fill[rojo!12] (0,3.6) rectangle (6.4,4.6);  \fill[rojo!12] (2.2,4.6) rectangle (3.8,5.2);
  \fill[azul!8]  (0,1.2) rectangle (6.4,3.6);
  \fill[rojo!12] (0,0.6) rectangle (6.4,1.2);
  \fill[rojo!35] (0,0)   rectangle (6.4,0.6);
  \draw (0,0) rectangle (6.4,5.2);
  \foreach \y in {0.6,1.2,3.6} \draw (0,\y) -- (6.4,\y);
  \draw (0,4.6) -- (2.2,4.6) -- (2.2,5.2); \draw (6.4,4.6) -- (3.8,4.6) -- (3.8,5.2);
  % contactos
  \fill[gris] (0.2,5.2) rectangle (2.0,5.32); \fill[gris] (4.0,5.2) rectangle (6.2,5.32);
  \fill[gris] (2.6,5.2) rectangle (3.4,5.32); \fill[gris] (0,-0.12) rectangle (6.4,0);
  \draw (1.1,5.32) -- (1.1,6.0) -- (5.1,6.0) -- (5.1,5.32);
  \draw (1.1,6.0) -- (0.4,6.0) node[circle,draw,inner sep=1.3pt,fill=white]{} node[left=2pt] {K};
  \draw (3.0,5.32) -- (3.0,5.6) node[circle,draw,inner sep=1.3pt,fill=white]{} node[right=2pt,font=\small] {G};
  \draw (3.2,-0.12) -- (3.2,-0.5) node[circle,draw,inner sep=1.3pt,fill=white]{} node[right=2pt] {A};
  % rótulos
  \node at (1.1,4.9) {$n_2^+$}; \node at (5.1,4.9) {$n_2^+$};
  \node at (1.1,4.1) {$p_2$}; \node at (1.1,2.4) {$n_1^-$}; \node at (1.1,0.9) {$p_1$}; \node at (1.1,0.3) {$p_1^+$};
  \node[font=\small] at (4.2,4.1) {$\SI{1e17}{cm^{-3}}$};
  \node[font=\small] at (4.2,2.4) {$\SIrange{1e13}{5e14}{cm^{-3}}$};
  \node[font=\small] at (4.2,0.9) {$\SI{1e17}{cm^{-3}}$};
  \node[font=\small] at (4.2,0.3) {$\SI{1e19}{cm^{-3}}$};
  % junturas
  \node[left,font=\small] at (0,4.6) {$J_3$};
  \node[left,font=\small] at (0,3.6) {$J_2$};
  \node[left,font=\small] at (0,1.2) {$J_1$};
  % espesores
  \draw[<->] (6.7,4.6) -- node[right,font=\small] {$\sim\SI{10}{\micro m}$, $\SI{1e19}{cm^{-3}}$} (6.7,5.2);
  \draw[<->] (6.7,3.6) -- node[right,font=\small] {\SIrange{30}{100}{\micro m}} (6.7,4.6);
  \draw[<->] (6.7,1.2) -- node[right,font=\small,align=left] {\SIrange{50}{1000}{\micro m}\\(bloquea la tensión)} (6.7,3.6);
  \draw[<->] (6.7,0) -- node[right,font=\small] {\SIrange{30}{50}{\micro m}} (6.7,1.2);
  % símbolo
  \begin{scope}[xshift=11.4cm,yshift=1.6cm]
    \draw (0,0) to[thyristor,l_=$ $] (0,2.2);
    \node[right] at (0,2.3) {A}; \node[right] at (0,-0.1) {K};
  \end{scope}
\end{tikzpicture}
\caption{Corte de un tiristor con los dopajes típicos de la diapo~8. La región $n_1^-$, poco dopada
y gruesa, es la región de arrastre: fija la tensión de bloqueo en los dos sentidos. A la derecha, el
símbolo.}
\label{fig:tiristor-estructura}
\end{figure}

\subsection{Modelo de dos transistores}

Cortando mentalmente las capas $n_1$ y $p_2$ en dos, el tiristor es un PNP ($Q_1=p_1n_1p_2$) y un
NPN ($Q_2=n_1p_2n_2$) conectados de modo que \textbf{el colector de cada uno alimenta la base del
otro} \diapo{5, 12}: realimentación positiva.

\begin{figure}[H]
\centering
\begin{minipage}{0.4\textwidth}\centering
\begin{tikzpicture}[x=1cm,y=1cm]
  \tikzset{capa/.style={draw,minimum width=1.4cm,minimum height=0.55cm,inner sep=0pt,font=\small}}
  % Q1 (PNP) a la izquierda
  \node[capa,fill=rojo!25] (a1) at (0,0) {$p_1$};
  \node[capa,fill=azul!10,below=0pt of a1] (a2) {$n_1$};
  \node[capa,fill=rojo!12,below=0pt of a2] (a3) {$p_2$};
  % Q2 (NPN) a la derecha, desplazado una capa hacia abajo
  \node[capa,fill=azul!10] (b1) at (2.0,-0.55) {$n_1$};
  \node[capa,fill=rojo!12,below=0pt of b1] (b2) {$p_2$};
  \node[capa,fill=azul!35,below=0pt of b2] (b3) {$n_2$};
  \draw[dashed,gris] (a2.east) -- (b1.west);
  \draw[dashed,gris] (a3.east) -- (b2.west);
  \draw (a1.north) -- ++(0,0.4) node[above] {A};
  \draw (b3.south) -- ++(0,-0.4) node[below] {K};
  \draw (b2.east) -- ++(0.5,0) node[right] {G};
  \node[font=\small] at (0,-1.75) {$Q_1$ (PNP)};
  \node[font=\small] at (2.0,0.05) {$Q_2$ (NPN)};
\end{tikzpicture}
\end{minipage}
\hfill
\begin{minipage}{0.55\textwidth}\centering
\begin{circuitikz}
  \draw (0,0) node[pnp,xscale=-1] (Q1) {};
  \draw (2.2,-1.6) node[npn] (Q2) {};
  \node[left] at (Q1.text) {$Q_1$};
  \node[right] at (Q2.text) {$Q_2$};
  \draw (Q1.E) -- ++(0,0.5) node[above] {A};
  \draw[->] ($(Q1.E)+(0.25,0.4)$) -- node[right,font=\small] {$\IA$} ($(Q1.E)+(0.25,0.05)$);
  \draw (Q1.C) |- (Q2.B);
  \draw (Q1.B) -| (Q2.C);
  \draw (Q2.E) -- ++(0,-0.5) node[below] {K};
  \draw (Q2.B) ++(-0.65,0) -- ++(0,-1.2) -- ++(-1.3,0) node[left] {G};
  \node[font=\small] at (1.35,0.25) {$I_{C2}=I_{B1}$};
\end{circuitikz}
\end{minipage}
\caption{Izquierda: partición de la estructura $p_1n_1p_2n_2$ en dos transistores que comparten
$n_1$ y $p_2$. Derecha: circuito equivalente \diapo{5, 12}, con $I_{B1}=I_{C2}$ e $I_{B2}=I_{C1}+\IG$. La corriente de colector de cada
transistor es la corriente de base del otro.}
\label{fig:dos-transistores}
\end{figure}

\paragraph{Derivación de la corriente de ánodo.}
Con los dos transistores en activa (estado de bloqueo directo), la relación de la hoja de fórmulas
$I_C=\alpha I_E+I_{CB0}$ \enhojatexto{} aplicada a cada uno da \diapo{12}
\[
  I_{C1}=\alpha_1\IA+I_{CO1},\qquad I_{C2}=\alpha_2I_K+I_{CO2},\qquad I_K=\IA+\IG ,
\]
donde las corrientes de fuga son las de Ebers--Moll con el emisor abierto,
$I_{CO}=I_{CB0}=I_{R0}(1-\alpha_F\alpha_R)$ (la diapo escribe $I_{CS}[1-\alpha_f\alpha_r]$). En el nodo de
ánodo, la corriente de emisor de $Q_1$ se reparte entre su colector y su base, y su base es el
colector de $Q_2$:
\[
  \IA=I_{C1}+I_{B1}=I_{C1}+I_{C2}=\alpha_1\IA+I_{CO1}+\alpha_2(\IA+\IG)+I_{CO2}.
\]
Despejando:

\begin{clave}[Corriente de ánodo del tiristor]
\[
  \IA=\frac{\alpha_2\IG+I_{CO1}+I_{CO2}}{1-(\alpha_1+\alpha_2)} .
\]
\begin{itemize}
  \item \textbf{Bloqueo}: $\alpha_1+\alpha_2\ll1$ y $\IA$ es apenas una corriente de fuga \diapo{12}.
  \item \textbf{Disparo}: cuando $\alpha_1+\alpha_2\to1$ el denominador se anula y la corriente crece
  sin límite: la realimentación se vuelve regenerativa y $\IA$ queda \textbf{limitada sólo por el
  circuito externo} \diapo{10}.
\end{itemize}
\end{clave}

\begin{intuicion}[¿Por qué los $\alpha$ cambian?]
Si los $\alpha$ fueran constantes, un tiristor con $\alpha_1+\alpha_2<1$ nunca se encendería. En
realidad los $\alpha$ \textbf{crecen con la corriente}: a corriente baja la corriente de emisor es
mayormente de recombinación en la zona desierta (Pierret §11.2.6), que no contribuye a la corriente
de colector, y la eficiencia de emisión es pobre. Además crecen con $\VAK$ \diapo{13}: al aumentar la
tensión, la zona desierta de $J_2$ se extiende en $n_1$ (achica la base de $Q_1$ $\Rightarrow$ sube
$\alpha_{pnp}$) y en $p_2$ (achica la base de $Q_2$ $\Rightarrow$ sube $\alpha_{npn}$). Por eso el
diseño usa $\alpha_1$ chico (base $n_1$ muy ancha) y $\alpha_2$ alto pero dependiente de la corriente.
\end{intuicion}

\begin{ejemplo}[Del bloqueo al disparo]
Con $\alpha_1=0.3$, $\alpha_2=0.5$ y $I_{CO1}+I_{CO2}=\SI{2}{nA}$, sin compuerta:
$\IA=\SI{2}{nA}/0.2=\SI{10}{nA}$ (bloqueo). Se inyecta $\IG=\SI{1}{mA}$ y, a medida que la corriente
sube, también los $\alpha$; tomando $\alpha_2\approx0.5$ en el numerador:
\begin{center}\small
\begin{tabular}{lcccc}
\toprule
$\alpha_1+\alpha_2$ & 0.8 & 0.95 & 0.99 & $\to1$ \\
$\IA$ & \SI{2.5}{mA} & \SI{10}{mA} & \SI{50}{mA} & limitada por el circuito \\
\bottomrule
\end{tabular}
\end{center}
Cada aumento de $\IA$ sube los $\alpha$, que suben $\IA$: el proceso se dispara solo.
\end{ejemplo}

\subsection{Característica $I$--$V$ y estados}

\begin{figure}[H]
\centering
\begin{tikzpicture}
\begin{axis}[
  width=0.9\textwidth, height=7.4cm,
  axis lines=middle, xmin=-7, xmax=7.2, ymin=-2.2, ymax=4.2,
  xtick=\empty, ytick=\empty,
  xlabel={$\VAK$}, ylabel={$\IA$},
  every axis plot/.append style={line width=1.1pt},
  clip=false,
]
  % estado de conducción
  \addplot[azul,smooth] coordinates {(0.65,1.1) (0.85,1.7) (1.1,2.8) (1.4,4.0)};
  % bloqueo directo IG=0, IG1, IG2
  \addplot[rojo,smooth] coordinates {(0,0) (2,0.15) (4,0.3) (5.5,0.55) (5.9,1.0)};
  \addplot[rojo,dashed,smooth] coordinates {(5.9,1.0) (4,1.15) (2,1.25) (0.75,1.3)};
  \addplot[naranja,smooth] coordinates {(0,0) (2,0.2) (3.5,0.45) (3.9,0.9)};
  \addplot[naranja,dashed,smooth] coordinates {(3.9,0.9) (2.5,1.15) (0.75,1.3)};
  \addplot[violeta,smooth] coordinates {(0,0) (1.4,0.3) (2.0,0.75)};
  \addplot[violeta,dashed,smooth] coordinates {(2.0,0.75) (1.4,1.1) (0.75,1.3)};
  % bloqueo inverso y avalancha
  \addplot[rojo,smooth] coordinates {(0,0) (-3,-0.15) (-5.5,-0.3) (-5.9,-0.8) (-6.05,-2.0)};
  % rótulos
  \node[font=\small,anchor=west] at (axis cs:5.95,1.0) {$V_{BO}$ ($\IG=0$)};
  \node[font=\small,anchor=west] at (axis cs:3.95,0.85) {$I_{G1}$};
  \node[font=\small,anchor=west] at (axis cs:2.05,0.7) {$I_{G2}$};
  \draw[gris,dotted] (axis cs:0,1.3) -- (axis cs:0.75,1.3);
  \node[font=\small,anchor=east] at (axis cs:0,1.3) {$I_H$};
  \node[font=\small,anchor=west,align=left] at (axis cs:1.5,3.2) {conducción\\($\VAK\approx\SIrange{1}{2}{V}$)};
  \node[font=\small,align=center] at (axis cs:3.6,-0.55) {bloqueo directo};
  \node[font=\small,align=center] at (axis cs:-3.3,0.45) {bloqueo inverso};
  \node[font=\small,anchor=east] at (axis cs:-6.1,-1.6) {$V_{RB}$ (avalancha)};
  \node[font=\small,rojo,align=center] at (axis cs:3.6,2.4) {resistencia negativa\\(inestable)};
  \draw[rojo,->] (axis cs:3.4,2.0) -- (axis cs:3.0,1.25);
\end{axis}
\end{tikzpicture}
\caption{Característica estática \diapo{5, 8}. Con $\IG=0$ el tiristor bloquea hasta la tensión de
ruptura directa $V_{BO}$; cuanto mayor $\IG$ ($I_{G2}>I_{G1}>0$), menor la tensión a la que se dispara. La rama
punteada (resistencia negativa) no es observable en régimen: el punto de trabajo salta de bloqueo a
conducción.}
\label{fig:tiristor-iv}
\end{figure}

\begin{definicion}[Corrientes de enganche y de mantenimiento]
\begin{itemize}
  \item \textbf{Corriente de enganche} $I_L$ (\emph{latching}): mínima $\IA$ que tiene que alcanzarse
  \emph{durante el pulso de compuerta} para que, al retirarlo, el tiristor quede encendido.
  \item \textbf{Corriente de mantenimiento} $I_H$ (\emph{holding}): mínima $\IA$ que mantiene el
  dispositivo en conducción; si $\IA<I_H$, $\alpha_1+\alpha_2$ cae por debajo de 1 y se apaga.
  Siempre $I_L>I_H$ (típicamente $I_L\approx2$--$3\,I_H$).
\end{itemize}
\end{definicion}

\paragraph{Bloqueo inverso} ($\VAK<0$) \diapo{12}. $J_1$ y $J_3$ en inversa, $J_2$ en directa.
$J_3$ tiene dopajes altos a ambos lados y rompe a pocos volts (\SIrange{20}{30}{V}), así que
\textbf{la que soporta la tensión es $J_1$}, y el límite lo pone el ancho de la zona $n_1$
(perforación hasta $J_2$) o su avalancha.

\paragraph{Bloqueo directo} ($\VAK>0$, sin disparo) \diapo{12}. $J_1$ y $J_3$ en directa, $J_2$ en
inversa. Por los dopajes, la zona desierta de $J_2$ se extiende en $n_1$: otra vez es la región
$n_1$ la que fija la tensión de bloqueo directa $V_{BO}$. Por eso un tiristor simétrico bloquea
aproximadamente la misma tensión en los dos sentidos.

\paragraph{Conducción} \diapo{13}. Las cuatro capas están inundadas de portadores minoritarios en
exceso; $J_2$ pasa a \textbf{directa} y los dos transistores del modelo están \textbf{saturados}.
La tensión es la suma de las tres junturas (la de $J_2$ con signo opuesto) más la región modulada:
\[
  \VAK(\mathrm{enc})=V_{J1}-V_{J2}+V_{J3}+V_{n^-}\approx\SIrange{0.8}{1}{V}\;(+\text{óhmica}) \qquad\text{\diapo{10, 13}}.
\]
Como en el diodo PiN, la $n_1$ está modulada y la caída casi no depende de la tensión de bloqueo:
es el dispositivo con menor caída por ampere para tensiones de varios kV.

\subsection{Diagramas de bandas}

\newcommand{\bandastir}[5]{% #1..#4: nivel de E_c en n2, p2, n1, p1; #5: hasta dónde llega J2 en n1
  \draw[bandac] (0,#1) -- (1.8,#1) to[out=0,in=180] (2.3,#2) -- (4.3,#2) to[out=0,in=180] (#5,#3)
     -- (9.3,#3) to[out=0,in=180] (9.8,#4) -- (12,#4);
  \draw[bandav] (0,{#1-1.1}) -- (1.8,{#1-1.1}) to[out=0,in=180] (2.3,{#2-1.1}) -- (4.3,{#2-1.1})
     to[out=0,in=180] (#5,{#3-1.1}) -- (9.3,{#3-1.1}) to[out=0,in=180] (9.8,{#4-1.1}) -- (12,{#4-1.1});
}
\begin{figure}[H]
\centering
\begin{tikzpicture}[x=0.95cm,y=0.95cm,>=Stealth]
  % cabecera de regiones
  \foreach \xa/\xb/\t in {0/2/{$n_2^+$ (K)}, 2/4.5/{$p_2$ (G)}, 4.5/9.5/{$n_1^-$}, 9.5/12/{$p_1^+$ (A)}}
     \node[font=\small] at ({(\xa+\xb)/2},2.0) {\t};
  \foreach \x/\j in {2/J_3, 4.5/J_2, 9.5/J_1} \node[font=\small,gris] at (\x,1.55) {$\j$};
  % (a) equilibrio
  \begin{scope}
    \bandastir{0.1}{1.0}{0.25}{1.05}{4.9}
    \draw[fermi] (0,0) -- (12,0) node[right,black,font=\small] {$E_F$};
    \node[font=\small,anchor=west] at (12.5,0.6) {(a) equilibrio};
  \end{scope}
  % (b) bloqueo directo
  \begin{scope}[yshift=-3.0cm]
    \bandastir{0.1}{0.85}{-1.25}{-0.6}{6.0}
    \node[font=\small,anchor=west] at (12.5,0.0) {(b) bloqueo directo};
    \node[font=\scriptsize,rojo,align=center] at (8.2,-0.4) {$J_2$ en inversa: cae casi toda $\VAK$};
    \draw[<->] (11.6,0.1) -- node[right,font=\small] {$q\VAK$} (11.6,-0.6);
    \draw[dashed,gris] (0,0.1) -- (11.8,0.1);
  \end{scope}
  % (c) conducción
  \begin{scope}[yshift=-6.7cm]
    \bandastir{0.1}{0.35}{0.0}{0.1}{4.9}
    \node[font=\small,anchor=west] at (12.5,0.0) {(c) conducción};
    \node[font=\small,verde!60!black,align=center] at (7,0.8) {todas las barreras bajas: inyección en las cuatro capas};
  \end{scope}
\end{tikzpicture}
\caption{Banda de conducción (azul) y de valencia (roja) del tiristor, cátodo a la izquierda como en
las diapos~9--11. (a) Equilibrio. (b) Bloqueo directo: $J_1$ y $J_3$ apenas en directa; la tensión
cae en la zona desierta de $J_2$, que se extiende en $n_1^-$. (c) Conducción: con las barreras
reducidas, los electrones del cátodo y los huecos del ánodo atraviesan toda la estructura
(proceso regenerativo, \diapo{10}). Esquemático: las escalas de energía no son exactas.}
\label{fig:tiristor-bandas}
\end{figure}

\subsection{Mecanismos de encendido}

Cualquier cosa que haga $\alpha_1+\alpha_2\to1$ enciende el tiristor:
\begin{enumerate}
  \item \textbf{Corriente de compuerta} (el normal) \diapo{13}. Una $\IG$ entrante polariza en
  directa $J_3$ y la $n_2^+$ inyecta electrones en $p_2$; éstos difunden por la base y son barridos
  por el campo de $J_2$ hacia $n_1$. En $n_1$ producen dos efectos: ensanchan la zona desierta de
  $J_2$ y, por neutralidad, obligan a $J_1$ a inyectar huecos desde el ánodo. Los huecos cruzan hacia
  $p_2$ y suman corriente de base para $Q_2$: realimentación.
  \item \textbf{Sobretensión}: si $\VAK>V_{BO}$, la multiplicación por avalancha en $J_2$ aporta la
  corriente que falta (Pierret generaliza $\IA=\dots/[1-M(\alpha_1+\alpha_2)]$). No destruye el
  dispositivo si la corriente está limitada, pero no es un modo de disparo controlado.
  \item \textbf{$\dd v/\dd t$}: $J_2$ en inversa es un capacitor $C_{J2}$. Un escalón rápido de tensión
  inyecta $i=C_{J2}\,\dd\VAK/\dd t$, que actúa como corriente de compuerta.
  \item \textbf{Temperatura}: las fugas $I_{CO}$ se duplican cada $\sim\SI{8}{\celsius}$ y los $\alpha$
  suben.
  \item \textbf{Luz} (LASCR): pares generados en $J_2$; se usa en tiristores de muy alta tensión
  disparados por fibra óptica (unidad 6).
\end{enumerate}

\begin{ejemplo}[Disparo espurio por $\dd v/\dd t$]
Con $C_{J2}=\SI{50}{pF}$ y un flanco de $\SI{200}{V/\micro s}$:
$i=\num{50e-12}\cdot\num{2e8}=\SI{10}{mA}$, del orden de la corriente de disparo de compuerta de un
tiristor de media potencia. Por eso la hoja de datos da un $(\dd v/\dd t)_{\max}$ y se usan
\emph{snubbers} RC en paralelo.
\end{ejemplo}

\subsection{Transitorio de encendido y $\dd i/\dd t$}

\begin{figure}[H]
\centering
\begin{tikzpicture}[x=1cm,y=1cm,>=Stealth]
  % iG
  \draw[->] (0,5.0) -- (10.5,5.0) node[right] {$t$};
  \node[left] at (0,5.5) {$i_G$};
  \draw[violeta,line width=1.1pt] (0,5.0) -- (1,5.0) -- (1,5.7) -- (5.5,5.7) -- (5.5,5.0) -- (10,5.0);
  % iA
  \draw[->] (0,2.6) -- (10.5,2.6) node[right] {$t$};
  \node[left] at (0,3.6) {$i_A$};
  \draw[azul,line width=1.1pt] (0,2.6) -- (2.2,2.6) -- (2.4,2.65) -- (4.0,3.9) -- (10,3.9);
  \draw[dashed,gris] (0,3.9) node[left,black,font=\small] {$I_o$} -- (4.0,3.9);
  \node[font=\small,anchor=east] at (3.1,3.45) {$\dd i_A/\dd t$};
  % vAK
  \draw[->] (0,0) -- (10.5,0) node[right] {$t$};
  \node[left] at (0,1.5) {$v_{AK}$};
  \draw[rojo,line width=1.1pt] (0,1.6) -- (2.3,1.6) -- (4.0,0.55) .. controls (5.5,0.25) and (7,0.12) .. (10,0.1);
  % tiempos
  \foreach \x in {1,2.3,4.0,8.5} \draw[dotted] (\x,-0.2) -- (\x,5.8);
  \draw[<->] (1,-0.35) -- node[below,font=\small] {$t_{d,\mathrm{enc}}$} (2.3,-0.35);
  \draw[<->] (2.3,-0.35) -- node[below,font=\small] {$t_r$} (4.0,-0.35);
  \draw[<->] (4.0,-0.35) -- node[below,font=\small] {$t_{ps}$} (8.5,-0.35);
\end{tikzpicture}
\caption{Encendido del tiristor \diapo{15}. $t_{d,\mathrm{enc}}$: retardo hasta que los $\alpha\to1$
cerca de la compuerta; $t_r$: subida de la corriente, con pendiente impuesta por el circuito;
$t_{ps}$: extensión del plasma, durante la cual el área encendida crece y $v_{AK}$ termina de bajar.}
\label{fig:tiristor-encendido}
\end{figure}

\begin{intuicion}[El plasma y el límite de $\dd i/\dd t$ \diapo{15, 16}]
La compuerta inyecta portadores en $p_2$ \emph{cerca de su contacto}: el tiristor se enciende
primero en una franja angosta junto a la compuerta, donde se forma una acumulación de portadores en
exceso (``plasma''). El plasma se expande lateralmente a $\sim\SI{0.1}{mm/\micro s}$, mucho más lento
que lo que puede crecer la corriente. Si $\dd i_A/\dd t$ es grande, toda la corriente pasa por la
pequeña área ya encendida: densidad de corriente enorme, punto caliente y destrucción. De ahí el
$(\dd i/\dd t)_{\max}$ del fabricante, que se respeta con una inductancia serie \diapo{13, 15}, y las
compuertas \emph{interdigitadas} o \emph{amplificadoras} de los tiristores grandes.
\end{intuicion}

\subsection{Apagado}

Un tiristor convencional \textbf{no se apaga desde la compuerta}: hay que llevar $\IA$ por debajo
de $I_H$ (conmutación natural, cuando la corriente alterna pasa por cero, o forzada por un circuito
auxiliar) y mantenerlo sin tensión directa el tiempo suficiente.

\begin{figure}[H]
\centering
\begin{tikzpicture}[x=1cm,y=1cm,>=Stealth]
  \draw[->] (0,0) -- (11,0) node[right] {$t$};
  \draw[->] (0,-2.4) -- (0,2.3);
  \draw[azul,line width=1.1pt] (0,1.6) -- (1.2,1.6) -- (3.0,-1.2) .. controls (3.3,-1.5) and (3.6,-0.4) .. (4.4,-0.05) -- (11,0);
  \node[azul,left] at (0,1.6) {$i_A$};
  \draw[rojo,line width=1.1pt] (0,0.1) -- (3.1,0.1) .. controls (3.4,-1.6) .. (3.9,-1.8) -- (6.6,-1.6) .. controls (7.0,0.2) .. (7.4,1.9) -- (11,1.9);
  \node[rojo,right] at (11,1.9) {$v_{AK}$};
  \node[rojo,font=\small] at (5.2,-2.1) {bloqueo inverso};
  \draw[dashed,gris] (2.25,0) -- (2.25,2.2);
  \draw[dashed,gris] (7.0,0) -- (7.0,2.2);
  \draw[<->] (2.25,2.1) -- node[above,font=\small] {$t_q$: tiempo de apagado} (7.0,2.1);
  \node[font=\small,align=left] at (9.2,1.0) {recién ahora puede\\volver la tensión directa};
\end{tikzpicture}
\caption{Apagado por conmutación \diapo{11}. Como en el diodo, la corriente se invierte mientras se
extraen los portadores almacenados (pico de corriente inversa al vaciarse $J_1$ y $J_3$); luego la
carga que queda en $J_2$ debe desaparecer por \emph{recombinación}, porque $J_2$ está en directa y
no la barre ningún campo. Si se reaplica tensión directa antes de $t_q$, el tiristor se reenciende.}
\label{fig:tiristor-apagado}
\end{figure}

\diapo{11, 14} Durante la conducción se inyectaron en $n_1$ y $p_2$ concentraciones de portadores
mucho mayores que las de equilibrio. Para apagar, la recombinación interna y el barrido de
portadores (acelerado si se invierte la polaridad, a costa de una considerable corriente inversa)
retiran la carga almacenada hasta sacar a los transistores de saturación y llevarlos a la zona
activa, donde la ganancia de lazo ya es menor que 1. El tiempo $t_q$ (decenas a cientos de
\si{\micro s}) es la razón por la que los tiristores se usan a frecuencia de red.

\begin{figure}[H]
\centering
\begin{tikzpicture}[x=1cm,y=1cm,>=Stealth,
    resis/.style={decorate,decoration={zigzag,amplitude=2pt,segment length=4pt}}]
  \fill[rojo!12] (0,0) rectangle (10,1.6);
  \fill[azul!35] (0.8,1.0) rectangle (4.4,1.6); \fill[azul!35] (5.6,1.0) rectangle (9.2,1.6);
  \fill[azul!8] (0,-1.2) rectangle (10,0);
  \draw (0,-1.2) rectangle (10,1.6); \draw (0,0) -- (10,0);
  \fill[gris] (0.8,1.6) rectangle (4.4,1.72); \fill[gris] (5.6,1.6) rectangle (9.2,1.72);
  \fill[gris] (0,1.6) rectangle (0.5,1.72); \fill[gris] (9.5,1.6) rectangle (10,1.72);
  \node[font=\small] at (2.6,1.3) {$n_2^+$ (cátodo)}; \node[font=\small] at (7.4,1.3) {$n_2^+$ (cátodo)};
  \node[font=\small] at (0.25,2.0) {G}; \node[font=\small] at (9.75,2.0) {G};
  \node[font=\small] at (5,0.3) {$p_2$}; \node[font=\small] at (5,-0.8) {$n_1$};
  % corriente de compuerta negativa: sale por los bordes
  \draw[violeta,resis] (2.4,0.55) -- (0.6,0.55);
  \draw[violeta,->] (0.6,0.55) -- (0.25,1.5);
  \draw[violeta,resis] (7.6,0.55) -- (9.4,0.55);
  \draw[violeta,->] (9.4,0.55) -- (9.75,1.5);
  % corriente de ánodo concentrada en el centro
  \foreach \x in {2.2,2.6,3.0,7.0,7.4,7.8} \draw[azul,->,line width=0.9pt] (\x,-1.1) -- (\x,0.95);
  \foreach \x in {1.2,3.8,6.2,8.8} \draw[azul!50,->] (\x,-1.1) -- (\x,-0.3);
  \node[font=\small,violeta,align=center] at (5,2.3) {$-\IG$ sólo saca portadores del borde del cátodo};
\end{tikzpicture}
\caption{Por qué una corriente de compuerta negativa no apaga un tiristor convencional \diapo{14}.
La corriente de compuerta circula lateralmente por la $p_2$, que tiene resistencia; sólo polariza en
inversa $J_3$ cerca del borde del cátodo. El centro de los cátodos (de área mucho mayor que la
compuerta) sigue conduciendo y la corriente se concentra allí. El GTO resuelve esto achicando y
entrelazando los cátodos (sección~\ref{sec:gto}).}
\label{fig:compuerta-negativa}
\end{figure}

% =====================================================================
\section{Triac y control de fase}
\label{sec:triac}
% =====================================================================

\subsection{El triac: un tiristor bidireccional}

El triac (\emph{triode for alternating current}) integra en una sola pastilla el equivalente de
\textbf{dos tiristores en antiparalelo con una compuerta común} \diapo{6}. Sus terminales principales
no se llaman ánodo y cátodo sino MT1 y MT2 (\emph{main terminals}); la compuerta G está del lado de MT1
y todas las tensiones se miden respecto de MT1. Es una estructura de cinco regiones ($npnpn$ vista en
un corte): según el signo de $V_{MT2-MT1}$ conduce una u otra mitad (Pierret §13.3).

\begin{figure}[H]
\centering
\begin{minipage}{0.3\textwidth}\centering
\begin{circuitikz}
  \draw (0,0) node[below] {MT1} to[triac,n=T] (0,2.6) node[above] {MT2};
  \draw (T.gate) -- ++(-0.6,0) node[left] {G};
\end{circuitikz}
\end{minipage}
\begin{minipage}{0.3\textwidth}\centering
\begin{circuitikz}
  \draw (0,0) node[below] {MT1} -- (0,0.3) -- (-0.8,0.3) to[thyristor,n=S1] (-0.8,2.3) -- (0,2.3) -- (0,2.6) node[above] {MT2};
  \draw (0,0.3) -- (0.8,0.3) to[thyristor,invert,n=S2] (0.8,2.3) -- (0,2.3);
\end{circuitikz}
\end{minipage}
\begin{minipage}{0.36\textwidth}\centering
\begin{tikzpicture}[x=1cm,y=1cm]
  \fill[azul!30] (0,2.2) rectangle (1.1,2.6);
  \fill[rojo!15] (1.1,2.2) rectangle (2.6,2.6);
  \fill[rojo!15] (0,1.8) rectangle (2.6,2.2);
  \fill[azul!8] (0,0.8) rectangle (2.6,1.8);
  \fill[rojo!15] (0,0.4) rectangle (2.6,0.8);
  \fill[rojo!15] (0,0) rectangle (1.5,0.4);
  \fill[azul!30] (1.5,0) rectangle (2.6,0.4);
  \fill[azul!30] (0,0.4) rectangle (0.5,0.62);
  \draw (0,0) rectangle (2.6,2.6);
  \node[font=\scriptsize] at (0.55,2.4) {$n$}; \node[font=\scriptsize] at (1.85,2.4) {$p$};
  \node[font=\scriptsize] at (1.3,2.0) {$p$}; \node[font=\scriptsize] at (1.3,1.3) {$n^-$};
  \node[font=\scriptsize] at (1.3,0.6) {$p$}; \node[font=\scriptsize] at (0.75,0.2) {$p$}; \node[font=\scriptsize] at (2.05,0.2) {$n$};
  \node[font=\scriptsize] at (0.25,0.51) {$n$};
  \fill[gris] (0,2.6) rectangle (2.6,2.7); \fill[gris] (0.6,-0.1) rectangle (2.6,0);
  \fill[gris] (0,-0.1) rectangle (0.4,0);
  \node[font=\small] at (1.3,2.95) {MT2}; \node[font=\small] at (1.6,-0.35) {MT1}; \node[font=\small] at (0.2,-0.35) {G};
\end{tikzpicture}
\end{minipage}
\caption{Triac: símbolo, equivalente con dos tiristores en antiparalelo y corte esquemático
\diapo{6}. Las regiones $n$ junto a cada terminal hacen de cátodo de uno u otro tiristor según la
polaridad; la $n$ junto a la compuerta permite dispararlo con $\IG$ de ambos signos.}
\label{fig:triac}
\end{figure}

\begin{figure}[H]
\centering
\begin{tikzpicture}
\begin{axis}[
  width=0.78\textwidth, height=6.4cm,
  axis lines=middle, xmin=-7, xmax=7.5, ymin=-4.2, ymax=4.2,
  xtick=\empty, ytick=\empty,
  xlabel={$V_{MT2-MT1}$}, ylabel={$I$},
  every axis plot/.append style={line width=1.1pt},
]
  \addplot[azul,smooth] coordinates {(0.65,1.1) (0.85,1.7) (1.1,2.8) (1.4,4.0)};
  \addplot[rojo,smooth] coordinates {(0,0) (2,0.15) (4,0.3) (5.5,0.55) (5.9,1.0)};
  \addplot[rojo,dashed,smooth] coordinates {(5.9,1.0) (3,1.2) (0.75,1.3)};
  \addplot[naranja,dashed,smooth] coordinates {(3.0,0.6) (2,1.1) (0.75,1.3)};
  \addplot[naranja,smooth] coordinates {(0,0) (2.6,0.3) (3.0,0.6)};
  \addplot[azul,smooth] coordinates {(-0.65,-1.1) (-0.85,-1.7) (-1.1,-2.8) (-1.4,-4.0)};
  \addplot[rojo,smooth] coordinates {(0,0) (-2,-0.15) (-4,-0.3) (-5.5,-0.55) (-5.9,-1.0)};
  \addplot[rojo,dashed,smooth] coordinates {(-5.9,-1.0) (-3,-1.2) (-0.75,-1.3)};
  \addplot[naranja,dashed,smooth] coordinates {(-3.0,-0.6) (-2,-1.1) (-0.75,-1.3)};
  \addplot[naranja,smooth] coordinates {(0,0) (-2.6,-0.3) (-3.0,-0.6)};
  \node[font=\small,anchor=south west] at (axis cs:5.9,1.05) {$V_{BO}$};
  \node[font=\small,anchor=east] at (axis cs:-5.95,-1.0) {$-V_{BO}$};
  \node[font=\small] at (axis cs:3.2,3.3) {cuadrante I};
  \node[font=\small] at (axis cs:-3.2,-3.3) {cuadrante III};
\end{axis}
\end{tikzpicture}
\caption{Característica del triac \diapo{6}: la del tiristor repetida en el tercer cuadrante. Se
dispara en los dos sentidos, con corriente de compuerta entrante o saliente.}
\label{fig:triac-iv}
\end{figure}

\begin{table}[H]
\centering\small
\begin{tabular}{cccl}
\toprule
Modo & $V_{MT2-MT1}$ & $\IG$ & Sensibilidad \\
\midrule
I$^+$   & $+$ & $+$ (entra) & la mejor \\
I$^-$   & $+$ & $-$ (sale)  & buena \\
III$^-$ & $-$ & $-$ (sale)  & buena \\
III$^+$ & $-$ & $+$ (entra) & mala (se evita) \\
\bottomrule
\end{tabular}
\caption{Modos de disparo del triac. En la práctica se usan I$^-$ y III$^-$ (compuerta siempre
negativa) o I$^+$ y III$^-$ (pulso del mismo signo que $V_{MT2}$).}
\label{tab:triac-cuadrantes}
\end{table}

\begin{intuicion}[Tiristor o triac \diapo{7}]
El triac es cómodo (un solo componente y un solo disparo para alterna), pero es más lento y soporta
menos $\dd v/\dd t$ de conmutación: al pasar la corriente por cero, la carga que quedó de un
semiciclo puede reencender el dispositivo en el semiciclo siguiente. Por eso se usa en \textbf{baja
potencia} (atenuadores de luz, control de motores universales, relés de estado sólido) y el
\textbf{tiristor en media, alta y muy alta potencia} (rectificadores controlados, transmisión HVDC,
arrancadores).
\end{intuicion}

\subsection{Control de fase}

La aplicación clásica de ambos es el \textbf{control de fase}: en cada semiciclo de la red se
dispara el dispositivo con un retardo $\alpha$ (ángulo de disparo) respecto del cruce por cero. El
dispositivo conduce desde $\omega t=\alpha$ hasta que la corriente pasa por cero (con carga
resistiva, en $\omega t=\pi$) y la carga recibe una ``porción'' de cada semiciclo \diapo{7}.

\begin{figure}[H]
\centering
\begin{circuitikz}
  \draw (0,0) to[sV,l={$v_s=V_m\sin\omega t$}] (0,2.6) -- (1.2,2.6) to[thyristor,n=S] (3.6,2.6)
        -- (5.0,2.6) to[R,l=$R$] (5.0,0) -- (0,0);
  \node[draw,font=\small,align=center] (ctl) at (2.9,3.9) {control de compuerta};
  \draw (S.gate) -- (S.gate |- ctl.south);
  \node at (5.9,2.3) {$+$}; \node at (5.9,1.3) {$v_o$}; \node at (5.9,0.3) {$-$};
\end{circuitikz}

\medskip
\begin{tikzpicture}
\begin{axis}[
  width=0.9\textwidth, height=4.0cm, name=a1,
  xmin=0, xmax=720, ymin=-1.15, ymax=1.15,
  xtick={0,60,180,360,420,540,720}, xticklabels={0,$\alpha$,$\pi$,$2\pi$,,$3\pi$,$4\pi$},
  ytick=\empty, ylabel={tiristor},
  tick label style={font=\scriptsize}, samples=361, domain=0:720,
]
  \addplot[gris,dashed] {sin(x)};
  \addplot[rojo,line width=1.1pt] {(mod(x,360)>=60 && mod(x,360)<=180) ? sin(x) : 0};
  \addplot[violeta,ycomb,line width=1.2pt] coordinates {(60,0.35) (420,0.35)};
\end{axis}
\begin{axis}[
  width=0.9\textwidth, height=4.0cm, at={(a1.south)}, anchor=north, yshift=-0.2cm,
  xmin=0, xmax=720, ymin=-1.15, ymax=1.15,
  xtick={0,60,180,240,360,420,540,600,720}, xticklabels={0,$\alpha$,$\pi$,$\pi{+}\alpha$,$2\pi$,,$3\pi$,,$4\pi$},
  ytick=\empty, ylabel={triac},
  tick label style={font=\scriptsize}, samples=361, domain=0:720,
  xlabel={$\omega t$},
]
  \addplot[gris,dashed] {sin(x)};
  \addplot[verde!60!black,line width=1.1pt] {(mod(x,180)>=60) ? sin(x) : 0};
  \addplot[violeta,ycomb,line width=1.2pt] coordinates {(60,0.35) (240,-0.35) (420,0.35) (600,-0.35)};
\end{axis}
\end{tikzpicture}
\caption{Control de fase con carga resistiva, $\alpha=\ang{60}$ como en la diapo~7. Arriba: tiristor
(media onda); los pulsos (violeta) en el semiciclo negativo no tienen efecto porque el tiristor está
en inversa. Abajo: triac (onda completa), disparado en los dos semiciclos.}
\label{fig:control-fase}
\end{figure}

Integrando la porción de senoide entre $\alpha$ y $\pi$:

\begin{clave}[Control de fase con carga resistiva]
\[
  \text{Tiristor (media onda):}\qquad
  V_{\mathrm{med}}=\frac{1}{2\pi}\int_\alpha^\pi V_m\sin\theta\dd\theta=\frac{V_m}{2\pi}(1+\cos\alpha)
\]
\[
  \text{Triac (onda completa):}\qquad
  V_{\mathrm{ef}}=\sqrt{\frac1\pi\int_\alpha^\pi V_m^2\sin^2\theta\dd\theta}
  =\frac{V_m}{\sqrt2}\sqrt{1-\frac{\alpha}{\pi}+\frac{\sin2\alpha}{2\pi}},
  \qquad P=\frac{V_{\mathrm{ef}}^2}{R}.
\]
\end{clave}

\begin{ejemplo}[Atenuador de \SI{220}{V} con triac]
Lámpara de $R=\SI{100}{\ohm}$ (supuesta constante) en la red de \SI{220}{V} ($V_m=\SI{311}{V}$):
\begin{center}\small
\begin{tabular}{lcccc}
\toprule
$\alpha$ & \ang{0} & \ang{60} & \ang{90} & \ang{120} \\
\midrule
$V_{\mathrm{ef}}$ (triac) & \SI{220}{V} & \SI{197}{V} & \SI{156}{V} & \SI{97}{V} \\
$P$ & \SI{484}{W} & \SI{389}{W} & \SI{242}{W} & \SI{95}{W} \\
$V_{\mathrm{med}}$ (tiristor, media onda) & \SI{99}{V} & \SI{74}{V} & \SI{50}{V} & \SI{25}{V} \\
\bottomrule
\end{tabular}
\end{center}
Con $\alpha=\ang{90}$ la potencia es exactamente la mitad: se entrega medio semiciclo de cada
uno, y por simetría de $\sin^2$ cada mitad tiene la misma energía.
\end{ejemplo}

\begin{intuicion}[El precio del control de fase]
La corriente de línea deja de ser senoidal: el encendido brusco en $\alpha$ genera armónicos e
interferencia (de ahí los filtros LC de los atenuadores) y el factor de potencia empeora aun con
carga resistiva. Para generar el pulso en $\alpha$, los atenuadores simples usan una red RC que
carga un capacitor hasta la tensión de disparo de un \emph{diac} (dos junturas en antiparalelo que
rompen a $\sim\SI{30}{V}$).
\end{intuicion}

% =====================================================================
\section{GTO (\emph{Gate Turn-Off Thyristor})}
\label{sec:gto}
% =====================================================================

El GTO es un tiristor modificado para que \textbf{pueda apagarse desde la compuerta} con una
corriente de compuerta negativa grande \diapo{17}. Se enciende como un tiristor; para apagarlo se
\emph{extrae} por la compuerta una fracción importante de la corriente de ánodo durante algunos
microsegundos.

\subsection{Diferencias estructurales con el tiristor}

\diapo{17} Para lograr la capacidad de desconexión, el GTO cambia tres cosas:
\begin{enumerate}
  \item \textbf{Estructura compuerta--cátodo muy interdigitada}: el cátodo se divide en cientos o
  miles de ``dedos'' angostos (de unos \SI{100}{\micro m}), cada uno rodeado de compuerta. Así ningún
  punto de la $p_2$ queda lejos de un contacto de compuerta y la caída lateral en $p_2$ durante el
  apagado es chica (justo el problema de la figura~\ref{fig:compuerta-negativa}).
  \item \textbf{Cátodos en forma de islas}: se graba el silicio alrededor de cada cátodo $n^+$, que
  queda como una meseta; la compuerta contacta la $p_2$ en el fondo del grabado.
  \item \textbf{Ánodo en cortocircuito}: zonas $n^+$ que atraviesan la $p^+$ del ánodo y conectan la
  $n_1$ directamente con el metal de ánodo.
\end{enumerate}

\begin{figure}[H]
\centering
\begin{tikzpicture}[x=1cm,y=1cm,>=Stealth]
  % capas
  \fill[rojo!12] (0,2.6) rectangle (10,3.4);
  \fill[azul!8] (0,0.6) rectangle (10,2.6);
  \fill[rojo!30] (0,0) rectangle (10,0.6);
  \foreach \x in {2.3,5.0,7.7} \fill[azul!40] (\x,0) rectangle (\x+0.6,0.6);
  % islas de cátodo
  \foreach \x in {0.8,3.5,6.2} {
    \fill[azul!35] (\x,3.4) rectangle (\x+1.6,4.0);
    \draw (\x,3.4) -- (\x,4.0) -- (\x+1.6,4.0) -- (\x+1.6,3.4);
    \fill[gris] (\x+0.1,4.0) rectangle (\x+1.5,4.12);
    \node[font=\small] at (\x+0.8,3.7) {$n_2^+$};
  }
  \foreach \x in {0.05,2.6,5.3,8.0} \fill[gris] (\x,3.4) rectangle (\x+0.5,3.52);
  \draw (0,0) rectangle (10,3.4); \draw (0,2.6) -- (10,2.6); \draw (0,0.6) -- (10,0.6);
  \fill[gris] (0,-0.12) rectangle (10,0);
  \node[font=\small] at (9.2,3.0) {$p_2$}; \node[font=\small] at (9.2,1.6) {$n_1^-$}; \node[font=\small] at (9.3,0.3) {$p_1^+$};
  \foreach \x in {2.3,5.0,7.7} \node[font=\scriptsize] at (\x+0.3,0.3) {$n^+$};
  \node[font=\small,above] at (1.6,4.15) {K}; \node[font=\small,above] at (4.3,4.15) {K}; \node[font=\small,above] at (7.0,4.15) {K};
  \node[font=\small,above] at (2.85,3.55) {G}; \node[font=\small,above] at (5.55,3.55) {G};
  \node[font=\small] at (5,-0.4) {A (metal de ánodo)};
  % flechas de huecos que salen por los cortos
  \draw[rojo,->,line width=0.8pt] (2.6,1.4) -- (2.6,0.65);
  \draw[rojo,->,line width=0.8pt] (5.3,1.4) -- (5.3,0.65);
  \node[rojo,font=\scriptsize,align=left,anchor=west] at (5.5,1.2) {los huecos en exceso\\salen por los $n^+$ del ánodo};
\end{tikzpicture}
\caption{Corte esquemático de un GTO \diapo{17}: cátodos en islas $n^+$ rodeados de contactos de
compuerta sobre la $p_2$, y ánodo con cortocircuitos $n^+$ que perforan la $p_1^+$.}
\label{fig:gto}
\end{figure}

\subsection{Ganancia de apagado}

Partimos del GTO encendido, con $Q_1$ y $Q_2$ saturados, y extraemos por la compuerta una corriente
$I_G'=-\IG>0$ \diapo{18}. La base de $Q_2$ recibe el colector de $Q_1$ menos lo que sale por la
compuerta:
\[
  I_{B2}=I_{C1}-I_G'=\alpha_1\IA-I_G'.
\]
Por otro lado, el colector de $Q_2$ es la base de $Q_1$, que es lo que queda de $\IA$ después de
restar su corriente de colector:
\[
  I_{C2}=I_{B1}=(1-\alpha_1)\IA .
\]
$Q_2$ sale de saturación (y se rompe la realimentación) cuando su corriente de base ya no alcanza
para sostener su corriente de colector:
\[
  I_{B2}<\frac{I_{C2}}{\beta_2},\qquad \beta_2=\frac{\alpha_2}{1-\alpha_2}
  \quad\Longrightarrow\quad
  \alpha_1\IA-I_G'<\frac{(1-\alpha_1)(1-\alpha_2)}{\alpha_2}\IA .
\]
Despejando $I_G'$:
\[
  I_G'>\IA\left[\alpha_1-\frac{(1-\alpha_1)(1-\alpha_2)}{\alpha_2}\right]
  =\IA\,\frac{\alpha_1+\alpha_2-1}{\alpha_2}.
\]

\begin{clave}[Ganancia de desconexión del GTO]
\[
  I_G'>\frac{\IA}{\boff},\qquad
  \boff=\frac{\alpha_2}{\alpha_1+\alpha_2-1}\qquad\text{\diapo{18}}
\]
Para que $\boff$ sea grande (poca corriente de compuerta para apagar):
\begin{itemize}
  \item $\alpha_2\approx1$: capa $p_2$ delgada y $n_2$ (cátodo) muy dopada, como un BJT de alto $\beta$.
  \item $\alpha_1$ chico: $n_1$ lo más gruesa posible y \textbf{vida media corta en $n_1$}, de modo que
  $\alpha_1+\alpha_2$ supere a 1 por poco.
\end{itemize}
\end{clave}

\begin{ejemplo}[Corriente de apagado]
Con $\alpha_2=0.95$ y $\alpha_1=0.3$: $\boff=0.95/(0.3+0.95-1)=3.8$. Para apagar $\IA=\SI{1000}{A}$
hay que extraer $I_G'>\SI{263}{A}$ durante el apagado. Aunque el pulso dura pocos microsegundos (la
energía es chica), el circuito de compuerta de un GTO es voluminoso: es su principal desventaja frente
al IGBT, y lo que llevó al IGCT (GTO con compuerta integrada de muy baja inductancia, $\boff\approx1$).
\end{ejemplo}

\subsection{El ánodo en cortocircuito}

\diapo{18, 19} Una vida media corta en $n_1$ baja $\alpha_1$ y acelera el apagado, pero aumenta las
pérdidas en conducción. El ánodo en cortocircuito resuelve el compromiso de otra manera:
\begin{itemize}
  \item Las zonas $n^+$ del ánodo \textbf{eliminan la barrera} para que los huecos en exceso de $n_1$
  salgan al metal: son un ``disipador'' de huecos. Durante el apagado, la carga almacenada se retira
  mucho más rápido que sólo por recombinación, sin tener que acortar tanto $\tau$.
  \item También reducen la eficiencia de inyección de $J_1$, es decir, bajan $\alpha_1$
  (exactamente lo que pide $\boff$).
  \item Precio: $J_1$ queda cortocircuitada y \textbf{el GTO no bloquea en inversa} \diapo{17}; la
  única juntura que lo intenta es $J_3$, que rompe a \SIrange{20}{30}{V} por sus dopajes altos. Y como
  no se le puede aplicar tensión inversa, tampoco hay corriente inversa que barra los portadores: la
  única forma de sacarlos es la recombinación interna y la difusión a través de los $n^+$
  \diapo{19}.
\end{itemize}
En circuitos que necesitan bloqueo inverso se agrega un diodo en serie.

\begin{ojo}[Redacción de la diapo 18]
La diapo dice: ``se reduce $I_{B2}$, se \emph{activaría} $Q_2$, el cual se empezaría a apagar por la
acción regenerativa''. Lo que ocurre es que al extraer corriente por la compuerta se reduce $I_{B2}$
y $Q_2$ \textbf{sale de saturación} y pasa a la región activa; recién entonces la realimentación
positiva actúa \emph{en sentido inverso}: menos $I_{C2}$ es menos base para $Q_1$, que entrega menos
$I_{C1}$, y así sucesivamente hasta el corte.
\end{ojo}

\subsection{Transitorio de apagado}

El apagado de un GTO tiene tres etapas: (1) \emph{almacenamiento}: la compuerta extrae carga de $p_2$
mientras $\IA$ todavía no cambia; (2) \emph{caída}: $J_3$ se recupera, $\IA$ cae bruscamente y
$\VAK$ sube; (3) \emph{cola}: la carga que quedó en $n_1$ (lejos de la compuerta) se va por
recombinación y por los cortos de ánodo, con una corriente de cola que, multiplicada por la tensión
ya alta, es una pérdida importante. Durante el apagado la corriente tiende a \emph{concentrarse} en
los dedos que tardan más en apagarse (filamentación), por lo que el GTO necesita un \emph{snubber}
que limite $\dd v/\dd t$.

% =====================================================================
\section{MOSFET de potencia}
\label{sec:mosfet}
% =====================================================================

El MOSFET de potencia es el mismo transistor de efecto de campo de la unidad de MOSFET, pero con
una estructura \textbf{vertical} que le permite bloquear cientos de volts y conducir decenas de
amperes: la corriente entra por la fuente arriba, cruza un canal superficial corto y \emph{baja}
por la región de arrastre hasta el drenaje, que es todo el dorso de la pastilla \diapo{20}. Es un
dispositivo \textbf{unipolar}: no hay carga almacenada y conmuta en decenas de nanosegundos.

\subsection{Estructura VDMOS}

\begin{figure}[H]
\centering
\begin{tikzpicture}[x=1.05cm,y=1.05cm,>=Stealth]
  % sustrato y deriva
  \fill[azul!35] (0,0) rectangle (10,0.8);
  \fill[azul!8] (0,0.8) rectangle (10,4.0);
  % cuerpos p
  \fill[rojo!18,rounded corners=10pt] (-0.3,2.8) rectangle (3.6,4.3);
  \fill[rojo!18,rounded corners=10pt] (6.4,2.8) rectangle (10.3,4.3);
  % fuentes n+
  \fill[azul!40,rounded corners=5pt] (1.5,3.5) rectangle (3.0,4.2);
  \fill[azul!40,rounded corners=5pt] (7.0,3.5) rectangle (8.5,4.2);
  % tapar lo que sobresale
  \fill[white] (-0.4,4.0) rectangle (10.4,4.4);
  \fill[white] (-0.4,2.7) rectangle (0,4.4); \fill[white] (10,2.7) rectangle (10.4,4.4);
  \draw (0,0) rectangle (10,4.0);
  % canales
  \draw[azul,line width=2pt] (3.0,3.96) -- (3.6,3.96);
  \draw[azul,line width=2pt] (6.4,3.96) -- (7.0,3.96);
  % óxido de compuerta, poli, óxido de aislación
  \fill[gris!30] (2.5,4.0) rectangle (7.5,4.12);
  \fill[naranja!45] (2.6,4.12) rectangle (7.4,4.5);
  \fill[gris!30] (2.35,4.0) rectangle (2.6,4.75); \fill[gris!30] (7.4,4.0) rectangle (7.65,4.75);
  \fill[gris!30] (2.35,4.5) rectangle (7.65,4.75);
  % metal de fuente
  \fill[gris] (0,4.0) -- (2.35,4.0) -- (2.35,4.75) -- (7.65,4.75) -- (7.65,4.0) -- (10,4.0) -- (10,4.95) -- (0,4.95) -- cycle;
  \fill[gris] (0,-0.15) rectangle (10,0);
  % rótulos de capas
  \node[font=\small] at (5,0.4) {sustrato $n^+$};
  \node[font=\small] at (8.6,1.7) {$n^-$};
  \node[font=\small,align=center] at (0.85,3.25) {cuerpo\\$p$};
  \node[font=\small,align=center] at (9.15,3.25) {cuerpo\\$p$};
  \node[font=\small] at (2.25,3.82) {$n^+$}; \node[font=\small] at (7.75,3.82) {$n^+$};
  \node[font=\small] at (5,4.31) {compuerta (poli)};
  % terminales
  \node[font=\small,white] at (1,4.75) {S};
  \draw (5,-0.15) -- (5,-0.5) node[below] {D};
  \draw (7.4,4.3) -- (10.6,4.3) node[right] {G};
  % anotaciones
  \draw[<-] (5.6,4.06) -- (6.2,5.4) node[above,font=\small] {óxido de compuerta};
  \draw[<-] (2.45,4.65) -- (1.6,5.4) node[above,font=\small] {óxido de aislación};
  \draw[<-] (3.3,3.9) -- (3.2,2.55) node[below,font=\small] {canal};
  \draw[<->] (10.75,0.8) -- node[right,font=\small,align=left] {región de\\arrastre de\\drenaje} (10.75,4.0);
  % camino de los electrones
  \draw[azul,->,line width=1.1pt] (2.4,3.85) .. controls (3.3,3.88) and (3.9,3.9) .. (4.4,3.6) -- (4.6,1.0);
  \draw[azul,->,line width=1.1pt] (7.6,3.85) .. controls (6.7,3.88) and (6.1,3.9) .. (5.6,3.6) -- (5.4,1.0);
  \node[azul,font=\small] at (5,2.3) {$e^-$};
  % zona JFET
  \draw[dashed,gris] (3.6,2.8) -- (3.6,4.0); \draw[dashed,gris] (6.4,2.8) -- (6.4,4.0);
  \node[font=\scriptsize,gris] at (5,3.3) {zona JFET};
  % parásitos
  \draw[rojo,line width=0.8pt] (1.2,2.85) -- (1.2,2.55);
  \node[rojo,font=\scriptsize,align=center] at (1.2,2.1) {diodo de cuerpo\\($p$--$n^-$)};
\end{tikzpicture}
\caption{Corte de una celda VDMOS (\emph{double-diffused}) \diapo{20}. Con $V_{GS}>V_{TH}$ se forma
una capa de inversión (canal) en la superficie del cuerpo $p$ bajo la compuerta; los electrones van
de la fuente $n^+$ al canal, de ahí a la región de arrastre y bajan hasta el drenaje. La metalización
de fuente cortocircuita la fuente $n^+$ con el cuerpo $p$. Un dispositivo real tiene miles de estas
celdas en paralelo.}
\label{fig:vdmos}
\end{figure}

\paragraph{Estructuras parásitas.} La estructura tiene, ``gratis'':
\begin{itemize}
  \item un \textbf{diodo de cuerpo} $p$--$n^-$ entre fuente (ánodo) y drenaje (cátodo), en antiparalelo
  con el canal: el MOSFET de potencia \emph{no bloquea en inversa}, y ese diodo sirve de diodo de
  rueda libre (aunque, siendo bipolar, tiene recuperación inversa lenta);
  \item un \textbf{BJT parásito} $n^+pn^-$ (fuente--cuerpo--deriva). Si se encendiera, el MOSFET
  perdería el control. Por eso el metal de fuente \textbf{cortocircuita fuente y cuerpo}: la base del
  parásito queda atada al emisor. Aun así, un $\dd v_{DS}/\dd t$ muy alto inyecta corriente capacitiva
  en el cuerpo que, a través de su resistencia lateral, puede polarizar la juntura base--emisor.
\end{itemize}
El símbolo de la diapo~20 incluye el diodo de cuerpo y, en el modelo detallado, las capacidades y el
BJT parásito.

\subsection{Física de operación y regiones}

\diapo{21} Con $V_{DS}>0$ y variando $V_{GS}$:
\begin{enumerate}
  \item $0<V_{GS}<V_{TH}$: la compuerta repele los huecos del cuerpo y forma una zona de
  \textbf{deplexión} bajo el óxido; no hay canal, $I_D\approx0$.
  \item $V_{GS}=V_{TH}$: comienza la capa de inversión.
  \item $V_{GS}>V_{TH}$: canal formado; $I_D$ crece con $V_{DS}$ (región óhmica) hasta que el canal se
  estrangula del lado del drenaje ($V_{DS}=V_{GS}-V_{TH}$) y la corriente satura.
\end{enumerate}

\begin{ojo}[``Zona de degradación'' \diapo{21}]
La primera figura de la diapo~21 ($V_{GS}$ chica, sin canal) está rotulada ``zona de degradación'':
se trata de la zona de \textbf{deplexión} (vaciamiento, \emph{depletion}) bajo la compuerta. No hay
ninguna degradación del dispositivo.
\end{ojo}

\begin{clave}[Ecuaciones del MOSFET canal $n$ \enhoja]
\[
  I_D=\mu_nC_{ox}\frac WL\left[(V_{GS}-V_{TH})V_{DS}-\frac{V_{DS}^2}{2}\right],\quad V_{DS}\le V_{GS}-V_{TH}
  \qquad\text{(óhmica)}
\]
\[
  I_D=\frac12\mu_nC_{ox}\frac WL(V_{GS}-V_{TH})^2,\quad V_{DS}>V_{GS}-V_{TH}
  \qquad\text{(estrangulamiento; ``activa'' en potencia)}
\]
$I_D=0$ si $V_{GS}<V_{TH}$ (corte). En la región óhmica con $V_{DS}$ chico,
$\RDSon\approx\left[\mu_nC_{ox}\frac WL(V_{GS}-V_{TH})\right]^{-1}+R_{\mathrm{resto}}$.
\end{clave}

\begin{figure}[H]
\centering
\begin{tikzpicture}
\begin{axis}[
  width=0.82\textwidth, height=6.4cm,
  xmin=0, xmax=12.5, ymin=0, ymax=15.5,
  axis lines=left, xtick={10}, xticklabels={$BV_{DSS}$}, ytick=\empty,
  xlabel={$V_{DS}$}, ylabel={$I_D$},
  every axis plot/.append style={line width=1.1pt}, domain=0:10, samples=80,
]
  \addplot[rojo] {(x<2) ? (2*x-x^2/2)*(1+0.01*x) : (2.0)*(1+0.01*x)};
  \addplot[naranja] {(x<3) ? (3*x-x^2/2)*(1+0.01*x) : (4.5)*(1+0.01*x)};
  \addplot[verde!60!black] {(x<4) ? (4*x-x^2/2)*(1+0.01*x) : (8.0)*(1+0.01*x)};
  \addplot[azul] {(x<5) ? (5*x-x^2/2)*(1+0.01*x) : (12.5)*(1+0.01*x)};
  \addplot[gris,dashed,domain=0:5.2] {x^2/2};
  \addplot[black,line width=1.3pt] coordinates {(0,0.05) (10,0.05) (10.4,13.5)};
  \node[font=\small,anchor=west] at (axis cs:10.5,11.5) {ruptura};
  \node[font=\small,anchor=south west] at (axis cs:8.2,13.3) {$V_{GS4}$};
  \node[font=\small,anchor=south west] at (axis cs:8.2,8.6) {$V_{GS3}$};
  \node[font=\small,anchor=south west] at (axis cs:8.2,4.8) {$V_{GS2}$};
  \node[font=\small,anchor=south west] at (axis cs:8.2,2.1) {$V_{GS1}$};
  \node[font=\small,gris,anchor=south] at (axis cs:5.2,13.6) {$V_{DS}=V_{GS}-V_{TH}$};
  \node[font=\small] at (axis cs:1.2,11) {óhmica};
  \node[font=\small,align=center] at (axis cs:5,10.6) {activa (estrangulamiento)};
  \node[font=\small] at (axis cs:4,0.8) {corte: $V_{GS}<V_{TH}$};
\end{axis}
\end{tikzpicture}
\caption{Característica de salida del MOSFET de potencia \diapo{20, 21}. Los nombres de las regiones
en la bibliografía de potencia (Mohan) son ``óhmica'' y ``activa''; en la de circuitos
(Sedra--Smith) son ``tríodo'' y ``saturación''. Como llave, el MOSFET se usa en corte y en la
óhmica.}
\label{fig:mosfet-salida}
\end{figure}

\begin{ojo}[Regiones en la diapo 20]
La diapo dice ``Región de estrangulamiento cuando $V_{GS}<V_{GS(th)}$''. Con $V_{GS}<V_{TH}$ el
dispositivo está en \textbf{corte} (no hay canal que estrangular). La región de estrangulamiento
(activa) es la de canal formado y $V_{DS}>V_{GS}-V_{GS(th)}$. La condición de la región óhmica de la
misma diapo ($V_{GS}-V_{GS(th)}>V_{DS}>0$) sí es correcta.
\end{ojo}

\subsection{Resistencia en conducción}

Siguiendo el camino de los electrones de la figura~\ref{fig:vdmos}:
\[
  \RDSon=R_{n^+}+R_{\mathrm{canal}}+R_{\mathrm{acum}}+R_{\mathrm{JFET}}+R_{\mathrm{arrastre}}+R_{\mathrm{sustrato}} .
\]
En un MOSFET de \SI{30}{V} domina el canal (por eso se usan estructuras \emph{trench}, con la
compuerta en una trinchera vertical que elimina la zona JFET y empaqueta más canal por
\si{cm^2}). En uno de \SI{600}{V} domina la región de arrastre y vale el límite unipolar
$\Ronsp\propto\VB^2$ de la sección~\ref{sec:deriva}: \textbf{la región de arrastre determina la
tensión de ruptura} \diapo{20} y, con ella, la resistencia mínima.

\begin{intuicion}[La superjuntura]
Los MOSFET ``CoolMOS'' (como el IPP90R340C3 de la diapo~35) rompen el límite unipolar: la región de
arrastre se hace con columnas alternadas $p$ y $n$ fuertemente dopadas. En bloqueo las columnas se
vacían lateralmente entre sí y el campo vertical queda casi uniforme (como en un PiN), aunque el dopaje
sea alto; en conducción la columna $n$ dopada conduce bien. El resultado es $\Ronsp\propto\VB$ en lugar
de $\VB^2$.
\end{intuicion}

\paragraph{Coeficiente de temperatura.} $\RDSon$ \textbf{aumenta} con la temperatura (la movilidad
baja como $T^{-2.3}$ aprox.): un MOSFET más caliente toma menos corriente. Esto hace que varios
MOSFET en paralelo repartan la corriente solos y que no tengan segunda ruptura, a diferencia de los
BJT de potencia.

\subsection{Capacidades y carga de compuerta}

La compuerta es un capacitor: el MOSFET no consume corriente de compuerta en régimen, pero en cada
conmutación hay que cargar y descargar sus capacidades \diapo{20, 35}:
\[
  C_{iss}=C_{GS}+C_{GD},\qquad C_{oss}=C_{DS}+C_{GD},\qquad C_{rss}=C_{GD}.
\]
$C_{GD}$ (capacidad Miller) es la más importante: durante la conmutación $v_{DS}$ cambia cientos de
volts mientras $v_{GS}$ queda fija, y toda la corriente de compuerta se usa para cargar $C_{GD}$.

\begin{figure}[H]
\centering
\begin{tikzpicture}[x=1cm,y=1cm,>=Stealth]
  \draw[->] (0,0) -- (8.2,0) node[right] {$Q_G$};
  \draw[->] (0,0) -- (0,3.6) node[above] {$v_{GS}$};
  \draw[azul,line width=1.2pt] (0,0) -- (2.2,1.8) -- (4.6,1.8) -- (7.6,3.1);
  \draw[dashed,gris] (0,1.8) node[left,black,font=\small] {$V_{\mathrm{plateau}}$} -- (2.2,1.8);
  \draw[dashed,gris] (0,0.9) node[left,black,font=\small] {$V_{TH}$} -- (1.1,0.9);
  \draw[dashed,gris] (2.2,0) -- (2.2,1.8); \draw[dashed,gris] (4.6,0) -- (4.6,1.8);
  \draw[dashed,gris] (7.6,0) -- (7.6,3.1); \draw[dashed,gris] (0,3.1) node[left,black,font=\small] {$V_{GG}$} -- (7.6,3.1);
  \draw[<->] (0,-0.35) -- node[below,font=\small] {$Q_{GS}$} (2.2,-0.35);
  \draw[<->] (2.2,-0.35) -- node[below,font=\small] {$Q_{GD}$ (Miller)} (4.6,-0.35);
  \draw[<->] (0,-1.0) -- node[below,font=\small] {$\Qg$ total} (7.6,-1.0);
  \node[font=\small,align=left,anchor=west] at (2.4,2.5) {meseta: $v_{DS}$ cae y\\$i_G$ carga $C_{GD}$};
\end{tikzpicture}
\caption{Curva de carga de compuerta. El tiempo de conmutación es aproximadamente
$t\approx Q/i_G$: con $\Qg=\SI{94}{nC}$ y un driver de \SI{1}{A}, unos \SI{94}{ns}.}
\label{fig:carga-compuerta}
\end{figure}

\begin{definicion}[Figura de mérito $\RDSon\cdot\Qg$]
Un MOSFET más grande (más celdas) tiene menos $\RDSon$ pero más $\Qg$, en proporción inversa. El
producto $\RDSon\Qg$ depende de la tecnología y no del tamaño: mide a la vez pérdidas de conducción
($\propto\RDSon$) y de conmutación y excitación ($\propto\Qg$). Cuanto menor, mejor \diapo{35}.
\end{definicion}

% =====================================================================
\section{IGBT (\emph{Insulated Gate Bipolar Transistor})}
\label{sec:igbt}
% =====================================================================

El IGBT combina lo mejor de los dos mundos: \textbf{se controla por tensión como un MOSFET}
(compuerta aislada, potencia de excitación casi nula) y \textbf{conduce como un dispositivo bipolar}
(región de arrastre modulada, caída casi independiente de la tensión de bloqueo) \diapo{22}. Es el
dispositivo dominante entre $\sim\SI{600}{V}$ y $\sim\SI{6.5}{kV}$: variadores de motores, trenes,
inversores fotovoltaicos y eólicos.

\subsection{Estructura}

Es un VDMOS al que se le cambió el sustrato $n^+$ por uno $p^+$ \diapo{22}:

\begin{figure}[H]
\centering
\begin{tikzpicture}[x=1.05cm,y=1.05cm,>=Stealth]
  \fill[rojo!35] (0,0) rectangle (10,0.6);
  \fill[azul!30] (0,0.6) rectangle (10,1.0);
  \fill[azul!8] (0,1.0) rectangle (10,4.0);
  \fill[rojo!18,rounded corners=10pt] (-0.3,2.8) rectangle (3.6,4.3);
  \fill[rojo!18,rounded corners=10pt] (6.4,2.8) rectangle (10.3,4.3);
  \fill[azul!40,rounded corners=5pt] (1.5,3.5) rectangle (3.0,4.2);
  \fill[azul!40,rounded corners=5pt] (7.0,3.5) rectangle (8.5,4.2);
  \fill[white] (-0.4,4.0) rectangle (10.4,4.4);
  \fill[white] (-0.4,2.7) rectangle (0,4.4); \fill[white] (10,2.7) rectangle (10.4,4.4);
  \draw (0,0) rectangle (10,4.0); \draw (0,0.6) -- (10,0.6); \draw (0,1.0) -- (10,1.0);
  \draw[azul,line width=2pt] (3.0,3.96) -- (3.6,3.96);
  \draw[azul,line width=2pt] (6.4,3.96) -- (7.0,3.96);
  \fill[gris!30] (2.5,4.0) rectangle (7.5,4.12);
  \fill[naranja!45] (2.6,4.12) rectangle (7.4,4.5);
  \fill[gris!30] (2.35,4.0) rectangle (2.6,4.75); \fill[gris!30] (7.4,4.0) rectangle (7.65,4.75);
  \fill[gris!30] (2.35,4.5) rectangle (7.65,4.75);
  \fill[gris] (0,4.0) -- (2.35,4.0) -- (2.35,4.75) -- (7.65,4.75) -- (7.65,4.0) -- (10,4.0) -- (10,4.95) -- (0,4.95) -- cycle;
  \fill[gris] (0,-0.15) rectangle (10,0);
  \node[font=\small] at (8.6,0.3) {$p^+$ (capa de inyección)};
  \node[font=\small] at (8.3,0.8) {$n^+$ (capa de compensación)};
  \node[font=\small] at (8.6,1.9) {$n^-$ (arrastre)};
  \node[font=\small,align=center] at (0.85,3.25) {cuerpo\\$p$};
  \node[font=\small] at (2.25,3.82) {$n^+$}; \node[font=\small] at (7.75,3.82) {$n^+$};
  \node[font=\small] at (5,4.31) {compuerta};
  \node[font=\small,white] at (1,4.75) {E (fuente)};
  \draw (5,-0.15) -- (5,-0.5) node[below] {C (drenaje)};
  \draw (7.4,4.3) -- (10.6,4.3) node[right] {G};
  % electrones
  \draw[azul,->,line width=1.1pt] (2.4,3.85) .. controls (3.3,3.88) and (3.9,3.9) .. (4.4,3.6) -- (4.6,1.1);
  \node[azul,font=\small] at (4.25,2.0) {$e^-$};
  % huecos
  \draw[rojo,->,line width=1.1pt] (5.4,0.5) -- (5.6,2.4) .. controls (5.8,2.9) and (6.6,3.0) .. (7.4,3.2) -- (8.7,3.3);
  \draw[rojo,->,line width=1.1pt] (2.9,0.5) -- (2.6,2.5) .. controls (2.5,3.0) and (1.9,3.2) .. (1.2,3.3) -- (0.4,3.6);
  \node[rojo,font=\small] at (6.0,1.6) {$h^+$};
  \draw[decorate,decoration={zigzag,amplitude=1.5pt,segment length=3pt},rojo!70!black] (7.2,3.38) -- (8.4,3.38);
  \node[rojo!70!black,font=\scriptsize] at (7.9,2.95) {$R_b$};
\end{tikzpicture}
\caption{Corte de una celda de IGBT (\emph{punch-through}) \diapo{22, 23}. Los electrones bajan
por el canal como en el MOSFET; al llegar a la $p^+$, ésta inyecta huecos que inundan la región de
arrastre. Los huecos que no se recombinan son recogidos por el cuerpo $p$ y fluyen lateralmente bajo
la fuente $n^+$ (resistencia $R_b$) hasta el contacto: es la corriente que puede encender el
tiristor parásito.}
\label{fig:igbt}
\end{figure}

\begin{ojo}[Nombres de los terminales]
Las diapositivas (siguiendo a Mohan) llaman a los terminales \emph{drenaje} y \emph{fuente}, como en
el MOSFET; las hojas de datos los llaman \emph{colector} (C) y \emph{emisor} (E), como en el BJT. Ojo
que el ``colector'' del IGBT es en realidad el \emph{emisor} del PNP interno (la $p^+$ que inyecta
huecos) y el ``emisor'' del IGBT es el colector del PNP. Los nombres hablan del uso externo, no de la
física interna.
\end{ojo}

\subsection{Funcionamiento: un MOSFET que excita un PNP}

Con $V_{GE}>V_{TH}$ el canal conduce y los electrones llegan a la región de arrastre $n^-$, que es la
\textbf{base del PNP} $p^+n^-p$ formado por la capa de inyección, la región de arrastre y el cuerpo.
La juntura $J_1$ ($p^+$--$n^-$, o $p^+$--$n^+$ buffer) queda en directa e inyecta huecos: la región
de arrastre se modula como en un diodo PiN y su resistencia cae varios órdenes de magnitud.
\textbf{La resistencia de la región de arrastre depende de la tensión de compuerta} \diapo{24},
porque la compuerta fija la corriente de electrones que alimenta la inyección.

\begin{figure}[H]
\centering
\begin{minipage}{0.4\textwidth}\centering
\begin{circuitikz}
  \draw (0,3.8) node[above] {C} -- (0,3.4) to[D,l_=$J_1$] (0,2.2) to[R,l_=$R_{\mathrm{arr}}$] (0,1.0) -- (0,0.85);
  \draw (0,0.4) node[nmos,anchor=D] (M) {};
  \draw (M.G) -- ++(-0.4,0) node[left] {G};
  \draw (M.S) -- ++(0,-0.3) node[below] {E};
\end{circuitikz}\\[-0.2em]
\small (a) aproximado
\end{minipage}
\hfill
\begin{minipage}{0.55\textwidth}\centering
\begin{circuitikz}
  \draw (2.2,2.8) node[pnp] (Q1) {};
  \draw (Q1.E) -- ++(0,0.4) node[above] {C};
  \draw (0.6,1.0) node[npn,xscale=-1] (Q2) {};
  \draw (Q2.C) |- (Q1.B);
  \draw (Q1.C) |- (Q2.B);
  \draw (Q1.C) -- (2.2,1.0) to[R,l=$R_b$,-*] (2.2,-0.6);
  \draw (Q2.E) |- (2.2,-0.6);
  \draw (-1.6,1.0) node[nmos] (M) {};
  \draw (M.D) |- (-1.0,2.8) to[R,l=$R_{\mathrm{arr}}$] (Q2.C |- Q1.B);
  \draw (M.S) |- (0.6,-0.6);
  \draw (M.G) -- ++(-0.4,0) node[left] {G};
  \draw (0.6,-0.6) -- (0.6,-0.9) node[below] {E};
  \node[right,font=\small] at (Q1.text) {PNP};
\end{circuitikz}\\[-0.2em]
\small (b) con el tiristor parásito
\end{minipage}
\caption{Circuitos equivalentes del IGBT \diapo{24}. (a) Para condiciones normales: una juntura en
directa, la región de arrastre modulada y el canal en serie. (b) Completo: el PNP principal (arriba) y el NPN
parásito (al centro) ($n^+$ fuente -- cuerpo $p$ -- $n^-$) forman un tiristor; $R_b$ es la resistencia lateral
del cuerpo.}
\label{fig:igbt-equiv}
\end{figure}

\begin{clave}[Caída en conducción del IGBT]
\[
  V_{CE(\mathrm{on})}=V_{J1}+V_{\mathrm{arrastre}}+I_C\,R_{\mathrm{canal}}\qquad\text{\diapo{24}}
\]
$V_{J1}\approx\SIrange{0.7}{0.8}{V}$ (por eso la característica de salida arranca recién después de
ese umbral, figura~\ref{fig:igbt-salida}); $V_{\mathrm{arrastre}}$ es chica gracias a la modulación.
Comparado con un MOSFET de igual tensión, a corriente alta el IGBT cae mucho menos; a corriente muy
baja el MOSFET gana, porque no tiene el escalón de \SI{0.7}{V}.
\end{clave}

\begin{figure}[H]
\centering
\begin{tikzpicture}
\begin{axis}[
  width=0.78\textwidth, height=5.8cm,
  xmin=0, xmax=8, ymin=0, ymax=15.5,
  axis lines=left, xtick={0.7}, xticklabels={$\approx0.7$\,V}, ytick=\empty,
  xlabel={$V_{CE}$}, ylabel={$I_C$},
  every axis plot/.append style={line width=1.1pt}, samples=80,
]
  \addplot[rojo,domain=0:8] {(x<0.7) ? 0 : ((x-0.7<2) ? (2*(x-0.7)-(x-0.7)^2/2) : 2)*(1+0.01*x)};
  \addplot[naranja,domain=0:8] {(x<0.7) ? 0 : ((x-0.7<3) ? (3*(x-0.7)-(x-0.7)^2/2) : 4.5)*(1+0.01*x)};
  \addplot[verde!60!black,domain=0:8] {(x<0.7) ? 0 : ((x-0.7<4) ? (4*(x-0.7)-(x-0.7)^2/2) : 8)*(1+0.01*x)};
  \addplot[azul,domain=0:8] {(x<0.7) ? 0 : ((x-0.7<5) ? (5*(x-0.7)-(x-0.7)^2/2) : 12.5)*(1+0.01*x)};
  \node[font=\small,anchor=south west] at (axis cs:6.5,13.1) {$V_{GE4}$};
  \node[font=\small,anchor=south west] at (axis cs:6.5,8.4) {$V_{GE3}$};
  \node[font=\small,anchor=south west] at (axis cs:6.5,4.7) {$V_{GE2}$};
  \node[font=\small,anchor=south west] at (axis cs:6.5,2.1) {$V_{GE1}$};
\end{axis}
\end{tikzpicture}
\caption{Característica de salida del IGBT \diapo{22}: como la del MOSFET pero desplazada por la
caída de $J_1$. La característica de transferencia ($I_C$ vs.\ $V_{GE}$) es la de un MOSFET de canal
$n$, con umbral $V_{GE(th)}$.}
\label{fig:igbt-salida}
\end{figure}

\subsection{PT, NPT y \emph{field-stop}}

\begin{itemize}
  \item \textbf{PT} (\emph{punch-through}, figura~\ref{fig:igbt}): con capa de compensación $n^+$
  (buffer). La región de arrastre es más delgada (campo trapezoidal, figura~\ref{fig:pt-npt}) y el
  buffer reduce la eficiencia de inyección del PNP y acelera el apagado. Precio: casi \textbf{no
  bloquea en inversa} ($J_1$ entre dos capas muy dopadas rompe a pocas decenas de volts).
  \item \textbf{NPT}: sin buffer, región de arrastre gruesa y $p^+$ delgada y poco eficiente.
  Bloquea en ambos sentidos, es robusto en cortocircuito y su $V_{CE(\mathrm{on})}$ tiene coeficiente
  de temperatura positivo (fácil de paralelar).
  \item \textbf{Field-stop} (moderno): combina ambos, con un buffer suave y una región de arrastre
  delgada; es el estándar actual.
\end{itemize}

\subsection{Latch-up: el tiristor parásito}

La estructura $n^+$(fuente)--$p$(cuerpo)--$n^-$--$p^+$ es un \textbf{tiristor} \diapo{23, 24}. Si el
NPN parásito se enciende, la suma de ganancias de los dos transistores llega a 1, el tiristor se
engancha y \textbf{la compuerta pierde el control}: el IGBT ya no se apaga y, si el circuito no limita
la corriente, se destruye.

\begin{intuicion}[¿Qué enciende al NPN parásito?]
La corriente de huecos que recoge el cuerpo $p$ circula lateralmente bajo la fuente $n^+$ por la
resistencia $R_b$ (figura~\ref{fig:igbt}). Esa caída polariza en directa la juntura
$n^+$--cuerpo, que es la base--emisor del NPN parásito. Cuando $I_hR_b\gtrsim\SI{0.7}{V}$, la fuente
inyecta electrones al cuerpo y se cierra el lazo regenerativo \diapo{24}.
\end{intuicion}

\begin{itemize}
  \item \textbf{Latch-up estático}: ocurre en conducción, por encima de una corriente máxima
  $I_{DM}$ que especifica el fabricante \diapo{24}.
  \item \textbf{Latch-up dinámico}: al apagar rápido, el canal se cierra, la corriente de electrones
  desaparece y \emph{toda} la corriente pasa a ser de huecos por el cuerpo, mientras la zona desierta
  de $J_2$ crece y acelera su barrido. Por eso el enganche se produce con corrientes \emph{menores}
  que $I_{DM}$ \diapo{24}. Corrientes capacitivas por $\dd v/\dd t$ alto tienen el mismo efecto.
  \item \textbf{Remedios}: el \textbf{cortocircuito cuerpo--fuente} en la metalización \diapo{22} (ata
  la base del NPN a su emisor), un cuerpo $p^+$ profundo bajo la fuente (baja $R_b$), fuentes $n^+$
  cortas, y un PNP de baja ganancia (buffer, control de vida media) para que $\alpha_{npn}+\alpha_{pnp}$
  no alcance 1.
\end{itemize}

\subsection{Apagado y corriente de cola}

Al quitar la tensión de compuerta el canal se cierra en nanosegundos y la corriente cae
bruscamente en la fracción de electrones. Pero la región de arrastre está llena de huecos y
electrones en exceso, y es la \emph{base} de un PNP que no tiene terminal propio: nadie puede extraer
esa carga. Desaparece por recombinación, dando una \textbf{corriente de cola} que decae con $\tau$,
mientras $v_{CE}$ ya está alta: es la principal pérdida de conmutación del IGBT y lo que limita su
frecuencia a $\sim\SIrange{1}{50}{kHz}$. Otra vez el compromiso de la sección~\ref{sec:recomb}: vida
media corta (o buffer) $\Rightarrow$ cola corta pero $V_{CE(\mathrm{on})}$ más alta.

\begin{table}[H]
\centering\small
\begin{tabularx}{\textwidth}{>{\raggedright\arraybackslash}p{3.3cm}>{\raggedright\arraybackslash}X>{\raggedright\arraybackslash}X}
\toprule
 & MOSFET de potencia (Si) & IGBT \\
\midrule
Portadores en conducción & sólo electrones (unipolar) & electrones $+$ huecos (bipolar) \\
Caída en conducción & $\RDSon I_D$, crece como $\VB^2$ & $\approx\SI{0.7}{V}+$ pequeña; casi independiente de $\VB$ \\
Tensiones típicas & hasta $\sim\SI{600}{V}$ (\SI{900}{V} superjuntura) & \SI{600}{V} a \SI{6.5}{kV} \\
Conmutación & muy rápida, sin cola & corriente de cola \\
Frecuencia típica & \SI{50}{kHz} a \si{MHz} & \SIrange{1}{50}{kHz} \\
Conducción inversa & diodo de cuerpo integrado & no (se agrega un diodo en antiparalelo) \\
Riesgo parásito & BJT parásito ($\dd v/\dd t$) & tiristor parásito (latch-up) \\
Control & por tensión, $\Qg$ & por tensión, $\Qg$ \\
\bottomrule
\end{tabularx}
\caption{MOSFET vs.\ IGBT.}
\label{tab:mosfet-igbt}
\end{table}

% =====================================================================
\section{Dispositivos de banda prohibida ancha (WBG): SiC}
\label{sec:wbg}
% =====================================================================

Todo lo anterior es silicio. Los materiales de \textbf{banda prohibida ancha} (\emph{wide band gap},
WBG) --- carburo de silicio (SiC) y nitruro de galio (GaN) --- tienen $\EG\approx\SI{3.3}{eV}$, tres
veces el del Si \diapo{25}. De ese único número salen sus dos grandes ventajas: \textbf{trabajan a
temperaturas mucho mayores} y \textbf{soportan campos diez veces mayores}.

\begin{table}[H]
\centering\small
\begin{tabular}{lcc}
\toprule
Material & $\EG$ (\si{\eV}) & $n_i$ a \SI{300}{K} (\si{cm^{-3}}) \\
\midrule
Ge & 0.66 & \num{2.4e13} \\
Si & 1.1 & \num{1.5e10} \\
GaAs & 1.4 & \num{1.8e6} \\
3C-SiC & 2.2 & 6.9 \\
6H-SiC & 3.0 & \num{2.3e-6} \\
4H-SiC & 3.26 & \num{8.2e-9} \\
GaN & 3.39 & \num{1.9e-10} \\
\bottomrule
\end{tabular}
\caption{Concentración intrínseca \diapo{25}. 3C, 6H y 4H son politipos del SiC (distintas
secuencias de apilamiento del cristal); en potencia se usa casi exclusivamente 4H-SiC por su mayor
$\EG$ y su mayor movilidad. Para el Si la diapo usa $n_i=\num{1.5e10}$; la tabla de constantes de
la materia usa \niSi, que es el valor más moderno.}
\label{tab:ni-wbg}
\end{table}

\subsection{Temperatura máxima de operación}

Una juntura funciona mientras el semiconductor sea extrínseco: $n_i\ll N_D$. Como
$n_i=\sqrt{N_cN_v}\,e^{-\EG/2kT}$ (Pierret §2.5), con $\EG$ triple el exponente es
triple.

\begin{ejemplo}[El mismo diodo a \SI{25}{\celsius} y a \SI{200}{\celsius} \diapo{26}]
Diodo de Si con $N_A=N_D=\SI{1e15}{cm^{-3}}$. A \SI{25}{\celsius}, $n_i\approx\SI{1e10}{cm^{-3}}$ y
los minoritarios son $n_i^2/N=\SI{1e5}{cm^{-3}}$. A \SI{200}{\celsius} (\SI{473}{K}), con
$N_c,N_v\propto T^{3/2}$ y $\EG(T)$ de Varshni ($\EG\approx\SI{1.08}{eV}$):
\[
  n_i(\SI{473}{K})=\sqrt{\num{3.22e19}\cdot\num{1.83e19}}\left(\tfrac{473}{300}\right)^{3/2}
  e^{-1.08/(2\cdot\num{8.61e-5}\cdot473)}\approx\SI{9e13}{cm^{-3}}\approx\SI{1e14}{cm^{-3}},
\]
y los minoritarios pasan a $\sim\SI{1e13}{cm^{-3}}$: \textbf{ocho órdenes de magnitud} más. La
corriente de saturación ($\propto n_i^2$) crece en la misma proporción y, poco más arriba, $n_i$
alcanza al dopaje y la juntura deja de existir.
\end{ejemplo}

\begin{figure}[H]
\centering
\begin{tikzpicture}
\begin{axis}[
  width=0.82\textwidth, height=6.6cm,
  ymode=log, xmin=0, xmax=600, ymin=1e-10, ymax=1e18,
  xlabel={$T$ (\si{\celsius})}, ylabel={$n_i$ (\si{cm^{-3}})},
  ytick={1e-10,1e-6,1e-2,1e2,1e6,1e10,1e14,1e18},
  grid=major, grid style={gris!20},
  legend pos=south east, legend style={font=\small},
  domain=0:600, samples=80,
  every axis plot/.append style={line width=1.1pt},
]
  \addplot[rojo] {sqrt(3.22e19*1.83e19)*((x+273)/300)^1.5*exp(-(1.17-4.73e-4*(x+273)^2/(x+273+636))/(2*8.61e-5*(x+273)))};
  \addlegendentry{Si}
  \addplot[azul] {sqrt(1.7e19*2.5e19)*((x+273)/300)^1.5*exp(-3.26/(2*8.61e-5*(x+273)))};
  \addlegendentry{4H-SiC}
  \addplot[verde!60!black] {sqrt(2.3e18*4.6e19)*((x+273)/300)^1.5*exp(-3.39/(2*8.61e-5*(x+273)))};
  \addlegendentry{GaN}
  \addplot[gris,dashed,line width=0.8pt] {1e15};
  \node[font=\small,anchor=north west] at (axis cs:300,1e15) {$N_D=\SI{1e15}{cm^{-3}}$ (arrastre)};
  \draw[rojo,dotted,line width=0.9pt] (axis cs:280,1e-10) -- (axis cs:280,1e15);
  \node[rojo,font=\small,anchor=east] at (axis cs:278,1e-4) {Si: $\approx\SI{280}{\celsius}$};
\end{axis}
\end{tikzpicture}
\caption{$n_i(T)$ calculada con $\sqrt{N_cN_v}(T/300)^{3/2}e^{-\EG/2kT}$ \diapo{27}. Para Si se usan
$N_c$, $N_v$ de la tabla de constantes y $\EG(T)$ de Varshni; para 4H-SiC y GaN, valores de
literatura con $\EG$ constante. El Si con $N_D=\SI{1e15}{cm^{-3}}$ se vuelve intrínseco a
$\sim\SI{280}{\celsius}$; el SiC y el GaN recién por encima de \SI{1000}{\celsius} (en la práctica los
limita el encapsulado, a $\sim\SIrange{200}{250}{\celsius}$).}
\label{fig:ni-T}
\end{figure}

\subsection{Comparación de materiales}

\begin{table}[H]
\centering\small
\begin{tabular}{lccccccc}
\toprule
Material & $\EG$ & $\mu_n$ & $\vsat$ & $\Ecr$ & $\lambda$ & $\epsilon_r$ & BFOM \\
 & (\si{\eV}) & (\si{cm^2/V.s}) & (\si{cm/s}) & (\si{V/cm}) & (\si{W/cm.K}) & & (rel.\ Si) \\
\midrule
Si      & 1.12 & 1350 & \num{1e7}   & \num{3e5}   & 1.5 & 11.8 & 1 \\
GaAs    & 1.42 & 8500 & \num{2e7}   & \num{4e5}   & 0.5 & 13.1 & 17 \\
3C-SiC  & 2.2  & 900  & \num{2e7}   & \num{1.2e6} & 4.5 & 9.7  & 35 \\
6H-SiC  & 3.0  & 370  & \num{2e7}   & \num{2.4e6} & 4.5 & 9.7  & 115 \\
4H-SiC  & 3.26 & 720  & \num{2e7}   & \num{2e6}   & 4.5 & 9.7  & 130 \\
GaN     & 3.39 & 900  & \num{2.5e7} & \num{3.3e6} & 1.3 & 9.0  & 680 \\
\bottomrule
\end{tabular}
\caption{Propiedades de los materiales \diapo{29} (se agregó $\epsilon_r$ y la figura de mérito de
Baliga $\epsilon\mu_n\Ecr^3$ relativa al Si, calculada con esos valores). En el GaN el valor de
$\mu_n$ es de volumen; en el canal 2D de un HEMT llega a \SIrange{1500}{2000}{cm^2/V.s}
(sección~\ref{sec:gan}).}
\label{tab:materiales}
\end{table}

\begin{figure}[H]
\centering
\begin{tikzpicture}[scale=1]
  % ejes: (nombre, max, ángulo)
  \def\R{2.6}
  \foreach \a/\lab in {90/{Campo crítico\\(MV/cm, máx.\ 4)},18/{Conductividad\\térmica (W/cm\,K, máx.\ 5)},
                     -54/{T. de fusión\\(\num{1000}\,\si{\celsius}, máx.\ 3)},-126/{Velocidad de sat.\\($10^7$\,cm/s, máx.\ 3)},
                     162/{Energía del gap\\(eV, máx.\ 4)}} {
    \draw[gris!50] (0,0) -- (\a:\R);
    \node[font=\scriptsize,align=center] at (\a:{\R+0.85}) {\lab};
  }
  \foreach \f in {0.25,0.5,0.75,1} \draw[gris!30] (90:\f*\R) -- (18:\f*\R) -- (-54:\f*\R) -- (-126:\f*\R) -- (162:\f*\R) -- cycle;
  % Si: Ec .3/4, k 1.5/5, Tf 1.42/3, vs 1/3, Eg 1.12/4
  \draw[naranja,line width=1.2pt,fill=naranja,fill opacity=0.12]
    (90:{0.3/4*\R}) -- (18:{1.5/5*\R}) -- (-54:{1.42/3*\R}) -- (-126:{1/3*\R}) -- (162:{1.12/4*\R}) -- cycle;
  % 4H-SiC: Ec 2, k 4.5, Tf 2.8, vs 2, Eg 3.26
  \draw[azul,line width=1.2pt,fill=azul,fill opacity=0.12]
    (90:{2/4*\R}) -- (18:{4.5/5*\R}) -- (-54:{2.8/3*\R}) -- (-126:{2/3*\R}) -- (162:{3.26/4*\R}) -- cycle;
  % GaN: Ec 3.3, k 1.3, Tf 2.5, vs 2.5, Eg 3.39
  \draw[rojo,line width=1.2pt,fill=rojo,fill opacity=0.12]
    (90:{3.3/4*\R}) -- (18:{1.3/5*\R}) -- (-54:{2.5/3*\R}) -- (-126:{2.5/3*\R}) -- (162:{3.39/4*\R}) -- cycle;
  % leyenda
  \begin{scope}[xshift=6.6cm,yshift=-0.6cm]
    \draw[naranja,line width=1.2pt] (0,0) -- (0.5,0) node[right,black,font=\small] {Si};
    \draw[azul,line width=1.2pt] (0,-0.45) -- (0.5,-0.45) node[right,black,font=\small] {4H-SiC};
    \draw[rojo,line width=1.2pt] (0,-0.9) -- (0.5,-0.9) node[right,black,font=\small] {GaN};
  \end{scope}
\end{tikzpicture}
\caption{Comparación de materiales \diapo{28}. Campo crítico y gap $\Rightarrow$ \textbf{alta
tensión}; conductividad térmica y temperatura de fusión $\Rightarrow$ \textbf{alta temperatura}
(el SiC disipa el calor tres veces mejor que el Si); velocidad de saturación $\Rightarrow$ \textbf{alta
frecuencia} (el GaN es el mejor).}
\label{fig:radar}
\end{figure}

\subsection{Dos uniones que soportan la misma tensión}

\diapo{30} Se construyen dos diodos $p^+n^-$, uno de Si y otro de SiC, para la misma $\VB$. Con las
fórmulas de la sección~\ref{sec:deriva}:
\[
  \frac{N_{D,\mathrm{SiC}}}{N_{D,\mathrm{Si}}}=\frac{\epsilon_{\mathrm{SiC}}\mathcal{E}_{c,\mathrm{SiC}}^2}{\epsilon_{\mathrm{Si}}\mathcal{E}_{c,\mathrm{Si}}^2},\qquad
  \frac{W_{\mathrm{Si}}}{W_{\mathrm{SiC}}}=\frac{\mathcal{E}_{c,\mathrm{SiC}}}{\mathcal{E}_{c,\mathrm{Si}}},\qquad
  \frac{R_{\mathrm{Si}}}{R_{\mathrm{SiC}}}=\frac{(\epsilon\mu_n\Ecr^3)_{\mathrm{SiC}}}{(\epsilon\mu_n\Ecr^3)_{\mathrm{Si}}}.
\]
Y si en cambio se construyen con la \emph{misma resistencia}, como $\Ronsp=4\VB^2/(\epsilon\mu_n\Ecr^3)$:
\[
  \frac{V_{B,\mathrm{SiC}}}{V_{B,\mathrm{Si}}}=\sqrt{\frac{(\epsilon\mu_n\Ecr^3)_{\mathrm{SiC}}}{(\epsilon\mu_n\Ecr^3)_{\mathrm{Si}}}}
  \qquad\text{\diapo{30}}.
\]

\begin{ejemplo}[Si vs.\ 4H-SiC para \SI{600}{V}]
Con los valores de la tabla~\ref{tab:materiales}:
\begin{center}\small
\begin{tabular}{lccc}
\toprule
 & $N_D$ & $W$ & $\Ronsp$ \\
\midrule
Si     & \SI{4.9e14}{cm^{-3}} & \SI{40}{\micro m} & \SI{38}{m\ohm.cm^2} \\
4H-SiC & \SI{1.8e16}{cm^{-3}} & \SI{6}{\micro m}  & \SI{0.29}{m\ohm.cm^2} \\
\midrule
cociente & 36.5 & 6.7 & 130 \\
\bottomrule
\end{tabular}
\end{center}
La región de arrastre del SiC se dopa 36 veces más y es casi 7 veces más delgada; su resistencia es
130 veces menor. A igual resistencia, el SiC bloquea $\sqrt{130}\approx11$ veces más tensión.
\end{ejemplo}

\begin{ojo}[Los números de la diapo 30]
La diapo da $N_{D,\mathrm{SiC}}=26.14\,N_{D,\mathrm{Si}}$, $R_{N,\mathrm{SiC}}=R_{N,\mathrm{Si}}/92.94$ y
un factor $9.6$ en tensión ($=\sqrt{92.94}$, consistente con la $R$). Con los valores de su propia
tabla (diapo~29: $\Ecr=\SI{2e6}{V/cm}$ y $\mu_n=720$ para 4H-SiC) se obtienen 36.5, 130 y 11.4. La
diferencia viene de los parámetros elegidos (el $\Ecr$ del SiC se cita entre 2 y \SI{3}{MV/cm}, y
depende del dopaje); lo robusto es el orden de magnitud: \textbf{$\sim10$ veces más tensión, o
$\sim100$ veces menos resistencia}.
\end{ojo}

\subsection{Diodos de SiC}

\diapo{31} Tres tipos:
\begin{itemize}
  \item \textbf{Schottky} (hasta $\approx\SIrange{600}{1200}{V}$, hoy \SI{1700}{V}): unipolar, sin
  carga almacenada ni recuperación inversa. Es el producto de SiC más maduro: reemplaza al diodo
  ultrarrápido de Si en fuentes con corrección de factor de potencia e inversores.
  \item \textbf{PN y PiN} (varios kV): bipolares, modulan la región de arrastre, pero con una caída
  de codo de casi \SI{3}{V} (ver ejemplo).
  \item \textbf{Mixtos} (JBS/MPS, tensiones intermedias): un contacto Schottky con islas $p^+$
  intercaladas. En inversa, las zonas desiertas de las islas se juntan y \emph{apantallan} la
  interfaz Schottky (menos fuga, más tensión); en directa conduce el Schottky y, ante sobrecorrientes,
  las islas inyectan como un PiN.
\end{itemize}
Regla general: \textbf{cuanto más Schottky, más rápido y menor caída en conducción, pero menor tensión
inversa y más fugas; cuanto más bipolar (PN), lo contrario} \diapo{31}.

\begin{ejemplo}[Por qué un diodo PN de SiC cae \SI{3}{V}]
La corriente de un diodo es $J=J_0(e^{V/V_T}-1)$ con $J_0\propto n_i^2$ (hoja de fórmulas,
$I_S$ \enhojatexto). Para la misma densidad de corriente,
\[
  V_{\mathrm{SiC}}-V_{\mathrm{Si}}=V_T\ln\frac{n_{i,\mathrm{Si}}^2}{n_{i,\mathrm{SiC}}^2}
  =2\cdot\SI{0.026}{V}\cdot\ln\frac{\num{1e10}}{\num{8.2e-9}}\approx\SI{2.2}{V}.
\]
Si un diodo de Si conduce a \SI{0.7}{V}, uno de 4H-SiC lo hace a $\approx\SI{2.9}{V}$: el gap grande
que permite bloquear tanto también encarece cada ampere en conducción bipolar. Por eso en SiC se
prefieren los dispositivos unipolares (Schottky, MOSFET, JFET) hasta varios kV.
\end{ejemplo}

\subsection{JFET de SiC}

\diapo{32} El JFET fue el primer transistor de SiC comercial, porque esquiva el problema del óxido
(abajo).
\begin{center}\small
\begin{tabularx}{0.95\textwidth}{XX}
\toprule
\textbf{Ventajas} & \textbf{Desventajas} \\
\midrule
Conducción unipolar (sin degradación bipolar). & Es \emph{normalmente cerrado} (\emph{normally-on}):
conduce con $V_{GS}=0$. \\
Rápido. & Se apaga con tensión de compuerta negativa. \\
Se fabrica sobre obleas $n$; estructura simple. & \\
\bottomrule
\end{tabularx}
\end{center}

\begin{ojo}[``Abierto'' y ``cerrado'']
En las diapositivas \emph{normalmente cerrado} significa lo que en inglés se dice
\emph{normally-on}: la llave está cerrada, es decir \textbf{conduce} sin tensión de compuerta (JFET,
HEMT de deplexión). \emph{Normalmente abierto} = \emph{normally-off} (MOSFET de enriquecimiento). Es la
convención de llaves y relés; no confundir ``abierto'' con ``canal abierto''.
\end{ojo}

\paragraph{¿Abiertos o cerrados? \diapo{33}} En el convertidor de la diapo~33 (puente de JFETs), al
energizar el circuito, antes de que el control genere la tensión negativa de compuerta, \textbf{los
JFET están todos conduciendo}: las dos llaves de una rama cortocircuitan la fuente (\emph{shoot-through}).
Un dispositivo normally-on exige una secuencia de arranque cuidadosa o una solución de circuito:

\begin{figure}[H]
\centering
\begin{circuitikz}
  \draw (0,3.0) node[njfet] (J) {};
  \draw (0,0.8) node[nmos] (M) {};
  \draw (J.D) -- ++(0,0.4) node[above] {D};
  \draw (J.S) -- (M.D);
  \draw (M.S) -- ++(0,-0.3) coordinate (s) -- ++(0,-0.3) node[below] {S};
  \draw (J.G) -- ++(-0.5,0) -- ++(0,-2.4) coordinate (gj) |- (s);
  \draw (M.G) -- ++(-0.35,0) node[left] {G};
  \node[right,font=\small] at (J.east) {JFET SiC (normally-on, alta tensión)};
  \node[right,font=\small] at (M.east) {MOSFET Si (normally-off, baja tensión)};
\end{circuitikz}
\caption{Cascodo JFET-SiC $+$ MOSFET-Si \diapo{32}. La compuerta del JFET está atada a la fuente del
conjunto. Con el MOSFET encendido, $V_{GS,\mathrm{JFET}}\approx0$ y el JFET conduce. Con el MOSFET
apagado, su $V_{DS}$ sube hasta que $V_{GS,\mathrm{JFET}}=-V_{DS,\mathrm{MOS}}$ alcanza la tensión de
estrangulamiento del JFET ($\sim\SI{-10}{V}$), que se apaga y toma el resto de la tensión. El
conjunto se comporta como una llave \emph{normally-off} controlada por la compuerta del MOSFET, que
sólo tiene que bloquear unas decenas de volts.}
\label{fig:cascodo}
\end{figure}

\subsection{MOSFET de SiC}

El SiC tiene una ventaja única entre los WBG: como el Si, \textbf{forma $\mathrm{SiO_2}$ por oxidación
térmica}, así que admite MOSFET de compuerta aislada. Pero \diapo{34}:
\begin{itemize}
  \item \textbf{Baja movilidad en el canal}: la interfaz $\mathrm{SiC/SiO_2}$ tiene muchas más trampas y
  cargas que la de Si (en parte por el carbono residual). Los electrones del canal tienen
  \SIrange{20}{50}{cm^2/V.s} frente a 720--\SI{900}{cm^2/V.s} en volumen: en un MOSFET de SiC la
  resistencia del canal pesa tanto como la de la región de arrastre.
  \item \textbf{El óxido es clave} y es difícil de lograr con buena calidad y confiabilidad (deriva de
  $V_{TH}$ con la temperatura y el uso).
  \item \textbf{No reemplazan directamente} a los MOSFET de Si: se excitan con $V_{GS}\approx+15$ a
  $+20$\,V para encender (vs.\ \SIrange{10}{12}{V}) y con $-3$ a $-5$\,V para asegurar el apagado.
\end{itemize}
La figura de la diapo~34 es un MOSFET \emph{lateral} (didáctico); los MOSFET de SiC comerciales son
verticales (DMOS o \emph{trench}) como el de la figura~\ref{fig:vdmos}.

\begin{ejemplo}[MOSFET SiC vs.\ MOSFET Si de \SI{900}{V} \diapo{35}]
\begin{center}\small
\begin{tabular}{lccc}
\toprule
 & $\RDSon$ & $\Qg$ & $\RDSon\Qg$ \\
\midrule
SiC (CREE C3M0065090J) & \SI{65}{m\ohm} & \SI{30}{nC} & \SI{1.95}{\ohm.nC} \\
Si superjuntura (Infineon IPP90R340C3) & \SI{340}{m\ohm} & \SI{94}{nC} & \SI{32.0}{\ohm.nC} \\
\bottomrule
\end{tabular}
\end{center}
Cociente: $32.0/1.95=16.4$, ``aproximadamente 16 veces mejor el SiC'' \diapo{35}. También las
capacidades: $C_{iss}=\SI{660}{pF}$ vs.\ \SI{2400}{pF}, $C_{oss}=\SI{60}{pF}$ vs.\ \SI{120}{pF}. El SiC
tiene cinco veces menos resistencia \emph{y} tres veces menos carga de compuerta a la vez.
\end{ejemplo}

% =====================================================================
\section{GaN: el transistor de alta movilidad (HEMT)}
\label{sec:gan}
% =====================================================================

\diapo{36} Con GaN se pueden hacer diodos parecidos a los de SiC, pero no es su interés principal.
Su dispositivo estrella es el \textbf{HEMT} (\emph{High Electron Mobility Transistor}), un
\textbf{transistor de heterounión}: funciona como un MESFET (un JFET con compuerta metálica, de
Schottky, en lugar de semiconductora), pero su canal no es una capa dopada sino un \textbf{gas
bidimensional de electrones} (2DEG) confinado en la interfaz entre dos semiconductores distintos. La
conductividad de ese canal es altísima y se logra con capacidades parásitas muy bajas.

\begin{ojo}[``Gas bidireccional'' \diapo{36}]
La diapo~36 dice ``gas \emph{bidireccional} de electrones''. Es gas \textbf{bidimensional} (2D gas,
2DEG): los electrones están confinados en un pozo de unos pocos nanómetros en la dirección
perpendicular a la interfaz y se mueven libremente en las otras dos. Las diapos~37 y 38 lo llaman
correctamente.
\end{ojo}

\subsection{La heterounión y el 2DEG}

Cuando se apoya un semiconductor de gap grande (AlGaN, o AlGaAs) sobre uno de gap menor (GaN, o
GaAs), la banda de conducción tiene un escalón $\Delta E_c$ en la interfaz. Del lado del material de
gap chico se forma un \textbf{pozo triangular} cuyo fondo queda \emph{por debajo} de $E_F$: allí se
acumulan electrones, que sólo pueden moverse paralelos a la interfaz.

\begin{figure}[H]
\centering
\begin{tikzpicture}[x=1cm,y=1cm,>=Stealth]
  % regiones
  \fill[gris!40] (-1.0,-4.2) rectangle (0,2.4);
  \fill[naranja!15] (0,-4.2) rectangle (2.2,2.4);
  \fill[azul!6] (2.2,-4.2) rectangle (9,2.4);
  \node[font=\small] at (-0.5,2.65) {metal};
  \node[font=\small] at (1.1,2.65) {AlGaN};
  \node[font=\small] at (5.6,2.65) {GaN (sin dopar)};
  % E_F
  \draw[fermi] (-1.0,0) -- (9,0) node[right,black,font=\small] {$E_F$};
  % Ec
  \draw[bandac] (0,1.9) -- (2.2,0.5) -- (2.2,-0.45) .. controls (2.6,-0.15) and (3.4,0.45) .. (5,0.65) -- (9,0.7);
  % Ev
  \draw[bandav] (0,-2.1) -- (2.2,-3.5) -- (2.2,-3.85) .. controls (2.6,-3.55) and (3.4,-2.95) .. (5,-2.75) -- (9,-2.7);
  % 2DEG
  \fill[azul!60,opacity=0.6] (2.2,-0.45) .. controls (2.35,-0.32) and (2.5,-0.2) .. (2.62,-0.1) -- (2.62,0) -- (2.2,0) -- cycle;
  \foreach \x/\y in {2.3/-0.15,2.42/-0.08,2.33/-0.32} \node[elec,minimum size=3.5pt] at (\x,\y) {};
  \draw[<-] (2.5,-0.12) -- (3.6,-0.9) node[right,font=\small,align=left] {2DEG: pozo triangular\\bajo $E_F$ ($\sim\SI{1e13}{cm^{-2}}$)};
  \draw[<->] (2.05,0.5) -- node[left,font=\small] {$\Delta E_c$} (2.05,-0.45);
  \node[azul,font=\small] at (7.5,1.0) {$E_c$};
  \node[rojo,font=\small] at (7.5,-2.4) {$E_v$};
  % cargas de polarización
  \foreach \y in {1.8,1.2,0.6,-1.8,-2.4,-3.0} \node[font=\scriptsize,verde!50!black] at (2.0,\y-0.15) {$+$};
  \node[font=\scriptsize,verde!50!black,align=center] at (1.0,-1.0) {carga de\\polarización\\$+\sigma_{\mathrm{pol}}$};
\end{tikzpicture}
\caption{Diagrama de bandas esquemático de una heterounión AlGaN/GaN bajo la compuerta. La
polarización del cristal deja una carga positiva fija en la interfaz; el campo resultante inclina la
banda del AlGaN y crea, del lado del GaN, un pozo triangular por debajo de $E_F$ donde se acumula el
2DEG, \textbf{sin necesidad de dopar} \diapo{38}.}
\label{fig:hemt-bandas}
\end{figure}

\paragraph{¿De dónde salen los electrones?}
\begin{itemize}
  \item \textbf{AlGaAs/GaAs} (HEMT tradicional, \diapo{37}): por \emph{dopaje modulado}. Se dopa el
  AlGaAs con donores; sus electrones caen al pozo del GaAs, que queda sin dopar. Los electrones quedan
  \emph{separados espacialmente} de los donores ionizados que los generaron: casi no hay dispersión
  por impurezas y la movilidad es muy alta (\SI{8500}{cm^2/V.s} a \SI{300}{K}, mucho más a baja
  temperatura). De ahí el nombre.
  \item \textbf{AlGaN/GaN}: el GaN y el AlGaN tienen estructura wurtzita, \emph{polar}: tienen
  polarización espontánea, y el AlGaN, al crecer tensionado sobre el GaN, además polarización
  piezoeléctrica. La discontinuidad de polarización deja en la interfaz una carga fija
  $\sigma_{\mathrm{pol}}/q\sim\SI{1e13}{cm^{-2}}$ que atrae una cantidad igual de electrones. Ambas
  capas pueden estar sin dopar \diapo{38}.
\end{itemize}

\begin{ejemplo}[Resistencia del canal 2DEG]
La resistencia de hoja de un canal de espesor despreciable con densidad superficial $n_s$ es
$R_\square=1/(qn_s\mu)$. Para el 2DEG de AlGaN/GaN, $n_s=\SI{1e13}{cm^{-2}}$ y
$\mu=\SI{2000}{cm^2/V.s}$:
\[
  R_\square=\frac{1}{\num{1.6e-19}\cdot\num{1e13}\cdot2000}\approx\SI{310}{\ohm}/\square .
\]
Un canal de inversión de un MOSFET de Si con la misma $n_s$ y $\mu\approx\SI{400}{cm^2/V.s}$ (la
movilidad superficial es mucho menor que la de volumen, \munSi) tendría \SI{1.6}{k\ohm}/$\square$: el
2DEG conduce cinco veces mejor con la misma carga, y por lo tanto con la misma capacidad de compuerta.
\end{ejemplo}

\subsection{El HEMT: estructura y control}

\begin{figure}[H]
\centering
\begin{tikzpicture}[x=1cm,y=1cm]
  \fill[rojo!25] (0,0) rectangle (10,0.8);
  \fill[gris!25] (0,0.8) rectangle (10,1.3);
  \fill[azul!8] (0,1.3) rectangle (10,3.0);
  \fill[naranja!40] (1.6,3.0) rectangle (8.4,3.35);
  \fill[azul!40,rounded corners=3pt] (0.3,2.4) rectangle (1.8,3.35);
  \fill[azul!40,rounded corners=3pt] (8.2,2.4) rectangle (9.7,3.35);
  \draw (0,0) rectangle (10,3.35);
  \draw[azul,line width=1.6pt,dashed] (1.8,2.95) -- (8.2,2.95);
  \fill[gris] (0.4,3.35) rectangle (1.7,3.75); \fill[gris] (8.3,3.35) rectangle (9.6,3.75);
  \fill[gris] (3.2,3.35) rectangle (4.2,3.75);
  \node[font=\small] at (1.05,4.0) {S}; \node[font=\small] at (3.7,4.0) {G}; \node[font=\small] at (8.95,4.0) {D};
  \node[font=\small] at (1.05,2.75) {$n^+$}; \node[font=\small] at (8.95,2.75) {$n^+$};
  \draw[<-] (7.4,3.2) -- (10.4,2.2) node[right,font=\small,align=left] {AlGaN\\($\sim\SI{20}{nm}$,\\sin dopar)};
  \node[font=\small] at (5,2.0) {GaN (sin dopar)};
  \node[font=\small] at (5,1.05) {capas de transición};
  \node[font=\small] at (5,0.4) {sustrato: Si, SiC o zafiro};
  \draw[<-] (5.6,2.92) -- (6.4,2.45) node[right,font=\small] {2DEG};
  \draw[<->] (4.2,4.45) -- node[above,font=\scriptsize] {$L_{GD}$: soporta la tensión} (8.3,4.45);
\end{tikzpicture}
\caption{HEMT de AlGaN/GaN \diapo{38}. Es un dispositivo \textbf{lateral}: la corriente circula por el
2DEG en el plano, y la tensión de bloqueo la soporta la distancia compuerta--drenaje $L_{GD}$ (y no el
espesor, como en los verticales). El GaN se crece sobre Si (barato, obleas grandes), SiC (mejor
disipación) o zafiro, con capas de transición que acomodan la diferencia de red cristalina.}
\label{fig:hemt-estructura}
\end{figure}

La compuerta (Schottky) controla el 2DEG como la de un MESFET \diapo{37}: con $V_{GS}=0$ el canal
existe y el transistor conduce; \textbf{polarizando la compuerta en inversa} ($V_{GS}<0$) se vacía el
pozo y se puede disminuir, e incluso eliminar, el 2DEG. El HEMT básico es por lo tanto un dispositivo
de \textbf{deplexión}, \emph{normalmente cerrado} (normally-on), con $V_{TH}\approx\SIrange{-3}{-5}{V}$.

\begin{intuicion}[RF vs.\ potencia \diapo{38}]
Con esta tecnología se construyeron primero amplificadores de radiofrecuencia (estaciones base,
radares), donde que el transistor sea normally-on no importa: trabaja polarizado en la región
lineal. En electrónica de potencia, en cambio, una llave que conduce cuando falta la tensión de
compuerta es inaceptable (el mismo problema del JFET de SiC, diapo~33). Hay que conseguir HEMT
\emph{normalmente abiertos}.
\end{intuicion}

\subsection{Cómo hacer un HEMT normalmente abierto}

\begin{figure}[H]
\centering
\begin{tikzpicture}[x=1cm,y=1cm]
  \foreach \dx/\tit in {0/{(a) Schottky (deplexión)},5.2/{(b) compuerta $p$-GaN},10.4/{(c) MIS con recess}} {
    \begin{scope}[xshift=\dx cm]
      \fill[azul!8] (0,0) rectangle (4.4,1.4);
      \node[font=\small] at (2.2,2.6) {\tit};
    \end{scope}
  }
  % (a)
  \fill[naranja!40] (0,1.4) rectangle (4.4,1.75);
  \fill[gris] (1.6,1.75) rectangle (2.8,2.15);
  \draw[azul,line width=1.5pt,dashed] (0,1.33) -- (4.4,1.33);
  \node[font=\scriptsize] at (2.2,0.6) {2DEG continuo: conduce con $V_{GS}=0$};
  % (b)
  \begin{scope}[xshift=5.2cm]
    \fill[naranja!40] (0,1.4) rectangle (4.4,1.75);
    \fill[rojo!40] (1.6,1.75) rectangle (2.8,2.05);
    \fill[gris] (1.7,2.05) rectangle (2.7,2.3);
    \draw[azul,line width=1.5pt,dashed] (0,1.33) -- (1.5,1.33);
    \draw[azul,line width=1.5pt,dashed] (2.9,1.33) -- (4.4,1.33);
    \node[font=\scriptsize,align=center] at (2.2,0.6) {la $p$-GaN levanta la banda:\\sin 2DEG bajo la compuerta};
  \end{scope}
  % (c)
  \begin{scope}[xshift=10.4cm]
    \fill[naranja!40] (0,1.4) rectangle (1.5,1.75); \fill[naranja!40] (2.9,1.4) rectangle (4.4,1.75);
    \fill[gris!40] (1.4,1.25) rectangle (3.0,1.4); \fill[gris!40] (1.4,1.4) rectangle (1.5,1.85); \fill[gris!40] (2.9,1.4) rectangle (3.0,1.85);
    \fill[gris] (1.5,1.4) rectangle (2.9,1.95);
    \draw[azul,line width=1.5pt,dashed] (0,1.33) -- (1.4,1.33);
    \draw[azul,line width=1.5pt,dashed] (3.0,1.33) -- (4.4,1.33);
    \node[font=\scriptsize,align=center] at (2.2,0.6) {se elimina el AlGaN bajo la compuerta;\\óxido $+$ canal tipo MOS};
  \end{scope}
\end{tikzpicture}
\caption{Compuertas de HEMT de GaN (sólo la zona de compuerta; naranja: AlGaN, rojo: $p$-GaN, gris
claro: dieléctrico). (a) Compuerta Schottky: normally-on. (b) Compuerta de GaN tipo $p$: su
potencial interno vacía el 2DEG bajo la compuerta a $V_{GS}=0$; $V_{TH}\approx\SIrange{1}{2}{V}$.
(c) Compuerta MIS con el AlGaN removido: se forma un canal por acumulación como en un MOSFET.}
\label{fig:hemt-compuertas}
\end{figure}

La diapo~39 clasifica los HEMT de AlGaN/GaN comerciales por tipo de compuerta y comportamiento:

\begin{table}[H]
\centering\small
\begin{tabularx}{\textwidth}{>{\raggedright\arraybackslash}p{3.0cm}>{\raggedright\arraybackslash}p{2.6cm}>{\raggedright\arraybackslash}X>{\raggedright\arraybackslash}p{3.4cm}}
\toprule
Comportamiento & Compuerta & Uso & Fabricantes (diapo 39) \\
\midrule
normalmente cerrado & metálica (MES) & RF (amplificación) & Cree, Nitronex, TriQuint \\
normalmente abierto & metálica sobre $p$-GaN & electrónica de potencia (directo) & EPC, Panasonic, GaN Systems \\
normalmente cerrado & MOS (MIS) & electrónica de potencia \textbf{en cascodo} con MOSFET de Si (fig.~\ref{fig:cascodo}) & Transphorm, MicroGaN, IR \\
normalmente abierto & MOS (MIS) & en desarrollo (``¿?'' en la diapo) & --- \\
\bottomrule
\end{tabularx}
\caption{Tipos de HEMT de AlGaN/GaN para potencia \diapo{39}.}
\label{tab:hemt-tipos}
\end{table}

\paragraph{Ejemplo de Infineon \diapo{40}.} El dispositivo de la diapo (desarrollado por International
Rectifier, hoy parte de Infineon) es un HEMT \emph{de deplexión} con compuerta MIS: un dieléctrico bajo
la compuerta (``Metal Insulator Semiconductor'') que reduce mucho la fuga de compuerta respecto de una
Schottky, fabricado con procesos CMOS estándar sobre sustrato de Si. Es \textbf{normalmente cerrado} y
está pensado para usarse \textbf{en cascodo} con un MOSFET de Si de baja tensión, exactamente como el
JFET de SiC (figura~\ref{fig:cascodo}): el conjunto es una llave normally-off. (La línea CoolGaN actual
de Infineon usa compuerta $p$-GaN, normally-off.)

\subsection{Lo que ganan y lo que cuesta el GaN}

\begin{itemize}
  \item \textbf{Velocidad}: $\Qg$ y capacidades muy chicas $\Rightarrow$ conmutación en pocos
  nanosegundos, convertidores de \SIrange{1}{10}{MHz} con inductores y capacitores diminutos (los
  cargadores USB-C de alta potencia son de GaN).
  \item \textbf{Sin recuperación inversa}: no hay juntura $pn$ ni diodo de cuerpo; el HEMT conduce en
  inversa por el propio canal (con una caída de unos volts) y no almacena carga.
  \item \textbf{Limitaciones}: es lateral (hasta \SI{650}{V} típicos; los de \SI{1200}{V} recién
  aparecen), la compuerta $p$-GaN tolera poca tensión (\SIrange{6}{7}{V}) y exige drivers cuidados,
  y las trampas del material producen \emph{colapso de corriente} ($\RDSon$ dinámica: la resistencia
  sube justo después de bloquear alta tensión). Su conductividad térmica (\SI{1.3}{W/cm.K}) es menor
  que la del SiC, por lo que para alta potencia se lo crece sobre SiC.
\end{itemize}

% =====================================================================
\clearpage
\section{Resumen y repaso}
% =====================================================================

\subsection{Los dispositivos de un vistazo}

\begin{table}[H]
\centering\small
\begin{tabularx}{\textwidth}{>{\raggedright\arraybackslash}p{2.0cm}>{\raggedright\arraybackslash}p{1.6cm}>{\raggedright\arraybackslash}p{2.6cm}>{\raggedright\arraybackslash}p{2.2cm}>{\raggedright\arraybackslash}X}
\toprule
Dispositivo & Tipo & Control & Bloqueo & Lo esencial \\
\midrule
Diodo PiN & bipolar & no controlado & directo no, inverso sí & modulación de conductividad; recuperación inversa ($\Qrr$, $\trr$) \\
Schottky & unipolar & no controlado & inverso (Si $\lesssim\SI{200}{V}$, SiC kV) & sin carga almacenada \\
Tiristor & bipolar & enciende por $\IG$; apaga sólo con $\IA<I_H$ & directo e inverso & $\IA=\frac{\alpha_2\IG+I_{CO}}{1-(\alpha_1+\alpha_2)}$; kV y kA, frecuencia de red \\
Triac & bipolar & $\IG$ de cualquier signo & ambos sentidos & dos tiristores antiparalelo; baja potencia, control de fase \\
GTO & bipolar & enciende y apaga por $\IG$ & directo (inverso casi nulo) & $\boff=\frac{\alpha_2}{\alpha_1+\alpha_2-1}$; ánodo en cortocircuito \\
MOSFET & unipolar & tensión $V_{GS}$ & directo; diodo de cuerpo en inverso & $\Ronsp\propto\VB^2$; muy rápido \\
IGBT & bipolar & tensión $V_{GE}$ & directo (NPT también inverso) & MOSFET $+$ PNP; cola de corriente; latch-up \\
MOSFET SiC & unipolar & tensión ($+15/{-5}$\,V) & directo & $\Ecr$ $\times$7--10; canal de baja movilidad \\
JFET SiC & unipolar & tensión negativa & directo & normally-on $\Rightarrow$ cascodo \\
HEMT GaN & unipolar, lateral & tensión & directo & 2DEG por polarización; normally-on salvo $p$-GaN; MHz \\
\bottomrule
\end{tabularx}
\caption{Resumen de los dispositivos de la unidad.}
\label{tab:resumen-dispositivos}
\end{table}

\needspace{0.62\textheight}
\subsection{Formulario}

\begin{tcolorbox}[base,colback=white,colframe=verde,title={Formulario}]
\renewcommand{\arraystretch}{1.8}
\small
\begin{tabularx}{\linewidth}{>{\raggedright\arraybackslash}p{0.2\linewidth}X}
Región de arrastre (NPT) & $\VB=\dfrac{\epsilon\Ecr^2}{2qN_D}$,\quad $W=\dfrac{2\VB}{\Ecr}$,\quad
$\Ronsp=\dfrac{W}{q\mu_nN_D}=\dfrac{4\VB^2}{\epsilon\mu_n\Ecr^3}$ \\
Región de arrastre (PT) & $\VB=\Ecr W_d-\dfrac{qN_DW_d^2}{2\epsilon}$ \\
Modulación & $\sigma=q(\mu_n+\mu_p)n$ con $n=p\gg N_D$;\quad $D_a=\dfrac{2D_nD_p}{D_n+D_p}$,\quad $L_a=\sqrt{D_a\tau_{HL}}$ \\
Recuperación inversa & $\Qrr\approx\tfrac12\IRM\trr$,\quad $\IRM\approx\sqrt{2\Qrr\,\dd i/\dd t}$,\quad $S=t_b/t_a$ \\
Vida media & $\tau_{HL}=\tau_{n0}+\tau_{p0}$,\quad $\tau_A=\dfrac{1}{C_a\delta n^2}$,\quad
$\tau=\dfrac{\tau_0}{1+(\delta n/n_b)^2}$ \\
Tiristor & $I_C=\alpha I_E+I_{CO}$ \enhojatexto,\quad $\IA=\dfrac{\alpha_2\IG+I_{CO1}+I_{CO2}}{1-(\alpha_1+\alpha_2)}$,\quad
$\VAK(\mathrm{enc})=V_{J1}-V_{J2}+V_{J3}+V_{n^-}$ \\
$\dd v/\dd t$ & $i=C_{J2}\,\dd\VAK/\dd t$ \\
GTO & $I_G'>\IA/\boff$,\quad $\boff=\dfrac{\alpha_2}{\alpha_1+\alpha_2-1}$ \\
Control de fase & $V_{\mathrm{med}}=\dfrac{V_m}{2\pi}(1+\cos\alpha)$ (media onda),\quad
$V_{\mathrm{ef}}=\dfrac{V_m}{\sqrt2}\sqrt{1-\dfrac\alpha\pi+\dfrac{\sin2\alpha}{2\pi}}$ (onda completa) \\
MOSFET \enhojatexto & $I_D=\mu_nC_{ox}\frac WL\left[(V_{GS}-V_{TH})V_{DS}-\frac{V_{DS}^2}{2}\right]$;\quad
$I_D=\frac12\mu_nC_{ox}\frac WL(V_{GS}-V_{TH})^2$ \\
Capacidades & $C_{iss}=C_{GS}+C_{GD}$,\quad $C_{oss}=C_{DS}+C_{GD}$,\quad $C_{rss}=C_{GD}$;\quad FOM $=\RDSon\Qg$ \\
IGBT & $V_{CE(\mathrm{on})}=V_{J1}+V_{\mathrm{arrastre}}+I_CR_{\mathrm{canal}}$;\quad latch-up si $I_hR_b\gtrsim\SI{0.7}{V}$ \\
WBG & $n_i=\sqrt{N_cN_v}\,e^{-\EG/2kT}$;\quad BFOM $=\epsilon\mu_n\Ecr^3$;\quad
igual $R$: $\dfrac{V_{B2}}{V_{B1}}=\sqrt{\dfrac{\mathrm{BFOM}_2}{\mathrm{BFOM}_1}}$ \\
Codo de un diodo & $\Delta V=V_T\ln\dfrac{n_{i1}^2}{n_{i2}^2}$ (misma $J$, distinto material) \\
2DEG & $R_\square=\dfrac{1}{qn_s\mu}$ \\
\end{tabularx}
\end{tcolorbox}

\subsection{Preguntas de repaso}

\begin{enumerate}
  \item \textbf{¿Por qué todos los dispositivos de potencia tienen una región $n^-$ gruesa y poco
  dopada?}\\
  Porque la tensión de bloqueo es el área bajo el campo, y el campo no puede superar $\Ecr$: para
  bloquear $\VB$ hace falta un ancho $W=2\VB/\Ecr$ y un dopaje bajo $N_D=\epsilon\Ecr^2/(2q\VB)$ para que
  la zona desierta se extienda en ese ancho.
  \item \textbf{¿Por qué el MOSFET de Si deja de servir a alta tensión y el IGBT no?}\\
  En el MOSFET la región de arrastre conduce con sus electrones de equilibrio: $\Ronsp\propto\VB^2$.
  En el IGBT la $p^+$ inyecta huecos que modulan la región de arrastre y su resistencia deja de
  depender del dopaje.
  \item \textbf{¿Qué es la inyección doble en un diodo PiN?}\\
  En alta inyección ($\delta p>N_D$) la carga de los huecos inyectados por la $p^+$ atrae electrones
  desde la $n^+$, que inyecta a su vez: la región de arrastre se llena de pares con $\delta n=\delta p$
  desde ambos extremos.
  \item \textbf{¿Por qué un diodo de potencia sigue conduciendo un tiempo después de invertirle la
  tensión?}\\
  Porque tiene carga almacenada ($\approx I_F\tau$) en la región de arrastre; mientras haya portadores
  en el borde de la juntura, ésta sigue en directa. La corriente inversa dura hasta extraer o
  recombinar esa carga ($\Qrr$).
  \item \textbf{¿Qué compromiso fija la vida media en un dispositivo bipolar de potencia?}\\
  $\tau$ largo: buena modulación y baja caída, pero mucha carga y conmutación lenta. $\tau$ corto: lo
  contrario. Además $\tau$ cambia con la temperatura (SRH crece) y con la inyección (Auger la reduce).
  \item \textbf{Explicar con el modelo de dos transistores por qué un tiristor se ``engancha''.}\\
  El colector de cada transistor es la base del otro; si $\alpha_1+\alpha_2\to1$ la ganancia del lazo
  llega a 1 y la corriente crece sin límite propio. Una vez encendido, la corriente de ánodo sostiene
  los $\alpha$ altos aunque se retire $\IG$, mientras $\IA>I_H$.
  \item \textbf{¿Qué otras cosas, además de la compuerta, pueden encender un tiristor?}\\
  Superar $V_{BO}$ (avalancha en $J_2$), un $\dd v/\dd t$ alto ($C_{J2}\dd v/\dd t$ actúa como
  $\IG$), la temperatura (más fugas) y la luz.
  \item \textbf{¿Por qué hay que limitar el $\dd i/\dd t$ al encender un tiristor?}\\
  Porque el encendido empieza en una franja junto a la compuerta y el plasma se expande a
  $\sim\SI{0.1}{mm/\micro s}$; si la corriente sube más rápido, se concentra en esa franja y la
  destruye.
  \item \textbf{¿Por qué una corriente de compuerta negativa no apaga un tiristor común y sí un GTO?}\\
  En el tiristor, la caída lateral en $p_2$ hace que la compuerta sólo afecte el borde del cátodo y el
  centro sigue conduciendo. El GTO tiene cátodos angostos e interdigitados (poca caída lateral), un
  $\alpha_1$ chico (ánodo en cortocircuito, vida corta) para que $\alpha_1+\alpha_2$ supere a 1 por poco,
  y así $\boff$ es razonable.
  \item \textbf{¿Por qué el GTO no bloquea en inversa?}\\
  Porque los cortocircuitos $n^+$ del ánodo anulan la juntura $J_1$; sólo queda $J_3$, que rompe a
  \SIrange{20}{30}{V}.
  \item \textbf{¿Qué es el latch-up de un IGBT y cómo se evita?}\\
  El encendido del tiristor parásito $n^+pn^-p^+$ cuando la corriente de huecos por $R_b$ polariza la
  juntura fuente--cuerpo. Se evita con el corto cuerpo--fuente, un cuerpo $p^+$ profundo, fuentes
  cortas y un PNP de baja ganancia. El dinámico ocurre al apagar, a corriente menor que la estática.
  \item \textbf{¿De dónde viene la corriente de cola del IGBT?}\\
  De la carga almacenada en la región de arrastre, que es la base sin terminal del PNP: al cerrarse el
  canal nada la extrae, y desaparece por recombinación con constante $\tau$.
  \item \textbf{¿Por qué el SiC permite hacer diodos Schottky de \SI{1200}{V} y el Si no?}\\
  Con $\Ecr$ $\sim7$--10 veces mayor, la región de arrastre para la misma tensión se dopa unas 36 veces
  más y es 7 veces más delgada: $\Ronsp$ baja $\sim100$ veces y un dispositivo unipolar vuelve a ser
  competitivo.
  \item \textbf{¿Por qué un diodo PN de SiC tiene una caída de casi \SI{3}{V}?}\\
  Porque $I_S\propto n_i^2$ es $\sim10^{36}$ veces menor que en Si: hace falta
  $\Delta V=V_T\ln(n_{i,\mathrm{Si}}^2/n_{i,\mathrm{SiC}}^2)\approx\SI{2.2}{V}$ más para la misma
  corriente.
  \item \textbf{¿Qué es el 2DEG y por qué en AlGaN/GaN no hace falta dopar?}\\
  Un gas de electrones confinado en un pozo triangular en la heterointerfaz, libre en dos dimensiones.
  En AlGaN/GaN la polarización espontánea y piezoeléctrica deja una carga fija positiva en la interfaz
  que atrae $\sim\SI{1e13}{cm^{-2}}$ electrones.
  \item \textbf{¿Por qué un JFET de SiC o un HEMT de GaN de deplexión se usan en cascodo?}\\
  Porque son normally-on: sin tensión de compuerta conducen, lo que en un puente cortocircuita la
  fuente al arrancar. El cascodo con un MOSFET de Si de baja tensión da una llave normally-off.
\end{enumerate}

\subsection{Correcciones a las diapositivas}

\begin{table}[H]
\centering\small
\begin{tabularx}{\textwidth}{>{\raggedright\arraybackslash}p{1.3cm}>{\raggedright\arraybackslash}X>{\raggedright\arraybackslash}X}
\toprule
Diapo & Dice & Corrección \\
\midrule
4 & $\tau$ disminuye por Auger para $\delta n$ cercano a $n_b\approx\SI{1e16}{cm^{-3}}$ & En Si la Auger
compite con SRH recién en $\sim\SIrange{1e17}{1e18}{cm^{-3}}$; $n_b$ es un parámetro de ajuste \\
18 & Al reducir $I_{B2}$ ``se activaría $Q_2$, el cual se empezaría a apagar'' & $Q_2$ \emph{sale de
saturación} y pasa a activa; recién ahí la realimentación actúa en sentido inverso hasta el corte \\
20 & ``Región de estrangulamiento cuando $V_{GS}<V_{GS(th)}$'' & Eso es \textbf{corte}; el
estrangulamiento (activa) es $V_{DS}>V_{GS}-V_{GS(th)}$ con canal formado \\
21 & ``Zona de degradación'' & Zona de \textbf{deplexión} (vaciamiento) bajo la compuerta \\
25 & $n_i(\mathrm{Si})=\num{1.5e10}$ & Valor antiguo; la tabla de constantes de la materia usa \niSi \\
30 & $N_D$ $\times26.14$, $R$ $/92.94$, tensión $\times9.6$ & Con los datos de la diapo 29: $\times36.5$,
$/130$, $\times11.4$; lo robusto es el orden de magnitud \\
34 & MOSFET de SiC dibujado lateral & Didáctico; los comerciales son verticales (DMOS/trench) \\
36 & ``Gas bidireccional de electrones'' & Gas \textbf{bidimensional} (2DEG) \\
\bottomrule
\end{tabularx}
\caption{Errores e imprecisiones de la presentación corregidos en el apunte.}
\label{tab:correcciones}
\end{table}


\end{document}
